\documentclass[11pt]{article}

\usepackage[margin=1in]{geometry}
\usepackage{amsmath,amssymb}
\usepackage{array,longtable}
\usepackage[hidelinks]{hyperref}

\setlength{\parindent}{0pt}
\setlength{\parskip}{0.6em}
\renewcommand{\arraystretch}{1.18}

\title{Master Note: Assignment, Sorting, and Nested Wage Models}
\author{}
\date{}

\begin{document}
\maketitle

\section{Task 1: Is \(C^w_{\mathrm{assign}}\) identified when the product benchmark uses missing cells?}

\subsection{1.1 Sources and answer}

This section uses two source documents:

\begin{itemize}
\item \textbf{LOO-v2:} \emph{Firm Wage Effects and Assignment Beyond Additivity},
  especially pp.\ 9--11, 15--20, 27, and 39--41.
\item \textbf{BS20:} Borovi\v{c}kov\'a and Shimer (2020),
  \emph{High Wage Workers Work for High Wage Firms}, especially pp.\ 5--7 and
  18--24.
\end{itemize}

The short answer has four parts.

\begin{enumerate}
\item The marginal distributions \(P_I\) and \(P_J\), and therefore their
  product \(P_I\otimes P_J\), can be learned from observed assignment
  frequencies.  The product \emph{measure} is not the unidentified object.
  \textit{Source: LOO-v2, p.\ 10, equation (18), and p.\ 15, equation (42).}
\item The variance
  \(\operatorname{Var}_{P_I\otimes P_J}(m_{IJ})\) generally is not identified
  on an incomplete worker--firm graph because it evaluates the maintained wage
  schedule \(m_{ij}\) at product-supported pairs that never occur.
  \textit{Source: LOO-v2, pp.\ 19--20, Proposition 4, Corollary 1, and
  equation (58).}
\item Consequently, \(C^w_{\mathrm{assign}}\) is not point identified on the
  unrestricted incomplete schedule.  This conclusion is conditional on the
  estimand, target population, support, and model class.  It is not a claim
  that identification is impossible after imposing valid completion
  restrictions.  \textit{Source: LOO-v2, pp.\ 19--20 and p.\ 27.}
\item The BS20 observed-match estimator does not identify the original
  \(C^w_{\mathrm{assign}}(m)\).  It identifies a different estimand: the
  covariance, and hence the correlation, between a worker's observed-match
  expected wage and a firm's observed-match expected wage.  It can therefore
  supply (i) an alternative observed-support sorting statistic or (ii) an
  estimator of the original target only after an additional bridge restriction
  connects the full pairwise schedule \(m_{ij}\) to the BS types.
  \textit{Source for the BS estimand: BS20, pp.\ 5--7, equations (1)--(6).
  Source for the estimator: BS20, pp.\ 21--24, Auxiliary Assumptions 1--4 and
  equations (18)--(25).  The non-equivalence and bridge statements are derived
  below.}
\end{enumerate}

\subsection{1.2 Complete notation chart}

\begin{longtable}{p{0.19\textwidth}p{0.53\textwidth}p{0.21\textwidth}}
\textbf{Symbol} & \textbf{Meaning} & \textbf{Source or status} \\
\hline
\(\mathcal I,\mathcal J\) & Sets of workers and firms. & LOO-v2, p.\ 9. \\
\(I,J\) & Numbers of workers and firms; also random worker and firm identities
when written inside expectations.  Context distinguishes the uses. &
LOO-v2, pp.\ 9--10. \\
\(i,j\) & Particular worker and firm indices.  The pair
\((i_0,j_0)\) denotes a selected missing cell. & LOO-v2, p.\ 9; the subscript
zero is derived notation. \\
\(\mathcal E\) & Observed bipartite edge set:
\((i,j)\in\mathcal E\) when the pair occurs in the data. & LOO-v2, p.\ 9. \\
\(\mathcal T\) & Declared target set on which the maintained schedule is
defined; often \(\operatorname{supp}(P_I\otimes P_J)\). & LOO-v2, pp.\ 9--10. \\
\(Y\) & Observed log wage. & LOO-v2, p.\ 9; BS20, pp.\ 5, 18--20. \\
\(m_{ij}\) & Maintained systematic wage for pair \((i,j)\), including target
pairs that may be unobserved. & LOO-v2, pp.\ 9--10, equations (14)--(15). \\
\(m_{\mathcal E}\) & Vector of maintained wage cells on observed edges. &
LOO-v2, p.\ 19, equation (53). \\
\(P^{\mathrm{obs}}_{IJ}\) & Observed joint worker--firm assignment
distribution. & LOO-v2, pp.\ 9, 15. \\
\(W=(w_{ij})\) & Matrix representation of \(P^{\mathrm{obs}}_{IJ}\), with
\(w_{ij}=P^{\mathrm{obs}}_{IJ}(I=i,J=j)\). & LOO-v2, p.\ 40. \\
\(P_I,P_J\) & Worker and firm marginals of \(P^{\mathrm{obs}}_{IJ}\). &
LOO-v2, pp.\ 10, 15. \\
\(p_i,q_j\) & Marginal probabilities
\(p_i=\sum_jw_{ij}\) and \(q_j=\sum_iw_{ij}\). & Derived from the definition
of marginals; LOO-v2, pp.\ 10, 40. \\
\(p,q\) & Vectors collecting \(p_i\) and \(q_j\). & Derived notation. \\
\(P^\times\) & Independent-assignment benchmark
\(P^\times=P_I\otimes P_J\), whose cell probability is \(p_iq_j\). &
LOO-v2, p.\ 10, equation (18). \\
\(P,R,r_{ij}\) & A generic probability distribution \(P\), or its probability
matrix \(R=(r_{ij})\), used when stating formulas that apply to either the
observed or product assignment. & Derived notation. \\
\(\operatorname{supp}(P)\) & Cells receiving positive probability under a
distribution \(P\). & Standard notation used in LOO-v2, p.\ 20. \\
\(\mathbb E_P,\operatorname{Var}_P,\operatorname{Cov}_P\) & Expectation,
variance, and covariance under distribution \(P\). & Standard operators. \\
\(\bar m_\times\) & Product-population mean
\(\mathbb E_{P^\times}[m_{IJ}]\). & Derived notation. \\
\(C^w_{\mathrm{assign}}(m)\) & One half of the difference between the observed
and product-assignment variances of the same fixed wage schedule. &
LOO-v2, p.\ 15, equation (42). \\
\(R_{\mathcal E}\) & Selection matrix that extracts observed cells from the
vectorized full schedule. & LOO-v2, p.\ 19, equation (53). \\
\(\mathcal K_{\mathcal E}\) & Missing-cell perturbation space
\(\ker(R_{\mathcal E})\). & LOO-v2, p.\ 19, equation (54). \\
\(e_{i_0j_0}\) & Coordinate vector equal to one at cell \((i_0,j_0)\) and zero
elsewhere. & Derived notation; LOO-v2 uses the same construction on p.\ 39. \\
\(\delta\) & Arbitrary scalar change to one missing cell. & Derived notation. \\
\(m^\delta\) & Observationally equivalent schedule
\(m+\delta e_{i_0j_0}\). & Derived notation. \\
\(\omega_0\) & Product probability \(p_{i_0}q_{j_0}\) on the selected missing
cell. & Derived notation. \\
\(\mu_0\) & Baseline value \(m_{i_0j_0}\) in the selected missing cell. &
Derived notation. \\
\(a_i,b_j,h_{ij}\) & Product-weighted worker component, firm component, and
double-centered interaction in
\(m_{ij}=\mu+a_i+b_j+h_{ij}\). & LOO-v2, p.\ 11, equations (19)--(23). \\
\(\mu\) & Grand product-weighted mean in the canonical decomposition. &
LOO-v2, p.\ 11, equation (19). \\
\(\mathcal X,\mathcal Y\) & BS20 worker- and firm-characteristic spaces. &
BS20, pp.\ 4--6. \\
\(x,y,z\) & BS20 worker characteristics, firm characteristics, and
match-specific wage determinant; \(x_i\) and \(y_j\) are the characteristics
of worker \(i\) and firm \(j\). & BS20, pp.\ 4--6. \\
\(w(x,y,z)\) & BS20 log-wage function. & BS20, pp.\ 4--6. \\
\(F,G\) & Cross-sectional distributions of worker and firm characteristics. &
BS20, pp.\ 4--7. \\
\(\Phi_x,\Psi_y\) & Conditional distributions of the firm matched to worker
characteristic \(x\), and of the worker matched to firm characteristic \(y\).
& BS20, pp.\ 4--6. \\
\(\lambda(x),\mu(y)\) & BS20 worker and firm types: expected log wage received
conditional on \(x\), and expected log wage paid conditional on \(y\). &
BS20, p.\ 5, equations (1)--(2). \\
\(\lambda_i,\mu_j\) & Shorthand for \(\lambda(x_i)\) and \(\mu(y_j)\). &
Derived notation. \\
\(\bar w\) & Common cross-sectional mean log wage. & BS20, p.\ 6,
equation (4). \\
\(\sigma_\lambda,\sigma_\mu\) & Cross-sectional standard deviations of BS20
worker and firm types; their squares are the corresponding variances. &
BS20, p.\ 6, equation (5); p.\ 20, equation (16). \\
\(c_{\mathrm{BS}}\) & Covariance between \(\lambda(X)\) and \(\mu(Y)\) among
observed matches. & BS20, p.\ 6, equation (6); pp.\ 20--21, equation (17). \\
\(\rho_{\mathrm{BS}}\) & BS20 correlation
\(c_{\mathrm{BS}}/(\sigma_\lambda\sigma_\mu)\). & BS20, p.\ 5,
equation (3). \\
\(M_i,N_j\) & Numbers of observed matches for worker \(i\) and firm \(j\). &
BS20, pp.\ 19, 21. \\
\(m,m',n,n'\) & BS20 focal and other-match indices from the worker and firm
perspectives.  These are indices, not the wage schedule \(m_{ij}\). &
BS20, pp.\ 19--20. \\
\(n(i,m)\) & Firm-side index of the same match that worker \(i\) labels as
match \(m\). & BS20, pp.\ 19--20. \\
\(w^w_{i,m},w^f_{j,n}\) & Average log wage in a match, indexed from the worker
and firm perspectives. & BS20, p.\ 19. \\
\(t^w_{i,m},t^f_{j,n}\) & Durations of those matches. & BS20, p.\ 19. \\
\(T_i^w,T_j^f\) & Total observed employment time of worker \(i\), and total
observed employment supplied by firm \(j\). & BS20, p.\ 19. \\
\(\varepsilon^w_{i,m},\varepsilon^f_{j,n}\) & Worker-side and firm-side
match-wage errors in the BS20 auxiliary equations. & BS20, p.\ 21,
Auxiliary Assumptions 1--4. \\
\(\widehat\lambda_i,\widehat{\lambda_i^2}\) & BS20 estimators of a worker's
type and squared type. & BS20, p.\ 22, equation (18). \\
\(\widehat\mu_j,\widehat{\mu_j^2}\) & BS20 estimators of a firm's type and
squared type. & BS20, p.\ 23, equation (21). \\
\(\widehat c_{i,m}\) & Leave-focal-match-out product estimating the type
product for the worker and firm in focal match \((i,m)\). & BS20, p.\ 23,
equation (23). \\
\(\widehat{\bar w},\widehat c_{\mathrm{BS}},\widehat\rho_{\mathrm{BS}}\) &
BS20 estimators of the mean wage, type covariance, and type correlation. &
BS20, pp.\ 23--24, equations (19), (24), and (25). \\
\(\tau\) & BS20 replication index: the numbers of each worker and firm
characteristic are multiplied by \(\tau\), and consistency is studied as
\(\tau\to\infty\). & BS20, p.\ 18. \\
\(s_{ij}\) & Constructed additive BS type-score schedule
\(\lambda_i+\mu_j-\bar w\). & Derived below. \\
\end{longtable}

\subsection{1.3 What is observed and what the product benchmark requires}

Write the observed assignment weights and their marginals as
\[
w_{ij}=P^{\mathrm{obs}}_{IJ}(I=i,J=j),\qquad
p_i=\sum_jw_{ij},\qquad q_j=\sum_iw_{ij}.
\]
Therefore the independent-assignment cell weights
\[
P^\times(I=i,J=j)=p_iq_j
\]
are determined by the observed joint assignment distribution.
\textit{Source: LOO-v2, p.\ 10, equation (18), and p.\ 40.}

For any probability matrix \(R=(r_{ij})\), the variance of the fixed schedule is
\[
\operatorname{Var}_{R}(m_{IJ})
=\sum_{i,j}r_{ij}m_{ij}^{2}
-\left(\sum_{i,j}r_{ij}m_{ij}\right)^2.
\]
\textit{Derived directly from the definition of variance.}

The target is
\[
C^w_{\mathrm{assign}}(m)
=\frac12\left\{
\operatorname{Var}_{P^{\mathrm{obs}}_{IJ}}(m_{IJ})
-\operatorname{Var}_{P^\times}(m_{IJ})
\right\}.
\tag{T1.1}
\]
\textit{Source: LOO-v2, p.\ 15, equation (42).}

If \(w_{ij}>0\), then \((i,j)\in\mathcal E\), so the observed variance in
(T1.1) uses observed-edge schedule values.  By contrast, \(p_iq_j\) can be
positive even when \(w_{ij}=0\).  The product variance then uses \(m_{ij}\) in
a missing cell.  Thus the problem is not that the product weights are unknown.
The problem is that the product weights multiply unknown wage levels and
squares.  \textit{Derived from (T1.1); the same distinction is stated in
LOO-v2, p.\ 20, equation (58).}

In actual wage data, even \(m_{ij}\) on an observed edge is a conditional mean
rather than a single raw wage.  The identification argument below gives the
observed-edge conditional means to the researcher and still finds
nonidentification.  Sampling noise in estimating those means is therefore a
second problem, not the source of the missing-cell failure.
\textit{Source for the noiseless identification experiment: LOO-v2, p.\ 19.}

\subsection{1.4 Formal missing-cell result}

\textbf{Proposition 1 (nonidentification of the original assignment target).}
Suppose \((i_0,j_0)\notin\mathcal E\) but
\(\omega_0=p_{i_0}q_{j_0}>0\).  On the unrestricted class of complete wage
schedules, \(C^w_{\mathrm{assign}}(m)\) is not point identified from
\((W,m_{\mathcal E})\).

\textit{Proof.}
For any scalar \(\delta\), define
\[
m^\delta=m+\delta e_{i_0j_0}.
\]
Because the changed cell is unobserved,
\[
R_{\mathcal E}m^\delta=R_{\mathcal E}m
\quad\text{and}\quad
\operatorname{Var}_{P^{\mathrm{obs}}_{IJ}}(m^\delta_{IJ})
=\operatorname{Var}_{P^{\mathrm{obs}}_{IJ}}(m_{IJ}).
\tag{T1.2}
\]
Let \(\bar m_\times=\mathbb E_{P^\times}[m_{IJ}]\) and
\(\mu_0=m_{i_0j_0}\).  Direct expansion gives
\[
\begin{aligned}
&\operatorname{Var}_{P^\times}(m^\delta_{IJ})
-\operatorname{Var}_{P^\times}(m_{IJ})\\
&\qquad
=2\omega_0\delta(\mu_0-\bar m_\times)
+\omega_0(1-\omega_0)\delta^2.
\end{aligned}
\tag{T1.3}
\]
Combining (T1.1)--(T1.3),
\[
\begin{aligned}
C^w_{\mathrm{assign}}(m^\delta)
-C^w_{\mathrm{assign}}(m)
=-\omega_0\delta(\mu_0-\bar m_\times)
-\frac12\omega_0(1-\omega_0)\delta^2.
\end{aligned}
\tag{T1.4}
\]
When \(0<\omega_0<1\), the right side is not constant in \(\delta\) and tends
to \(-\infty\) as \(|\delta|\to\infty\).  Hence observationally equivalent
schedules imply different values of the target.  This proves failure of point
identification and shows why, without a bound on missing wages, the identified
set is unbounded below. \(\square\)

\textit{Derivation: equations (T1.2)--(T1.4) expand the single-cell
perturbation used in LOO-v2, p.\ 39, Appendix A.2.  The general
kernel-invariance criterion is LOO-v2, p.\ 19, Proposition 4.  The
nonidentification and unboundedness conclusion is LOO-v2, p.\ 20,
Corollary 1.}

The qualification ``on the unrestricted class'' is essential.  The target
becomes identified if every product-supported cell is observed, or if a
defended model uniquely determines all target-relevant missing-cell contrasts.
Low rank, additivity, grouped types, parametric response surfaces, monotonicity
plus sufficiently sharp bounds, or a restricted opportunity-set benchmark are
possible restrictions, but each changes the model or the estimand.
\textit{Source: LOO-v2, pp.\ 20--27, especially Sections 5.4, 6, and 7.6.
The list is a synthesis of those sections.}

\subsection{1.5 The cell problem survives the BS minimum-match requirement}

Consider three workers and three firms with observed edges
\[
\mathcal E
=\{(1,1),(1,2),(2,2),(2,3),(3,3),(3,1)\}.
\tag{T1.5}
\]
Give each edge probability \(1/6\).  Every worker has two observed firms and
every firm has two observed workers, while the worker and firm marginals are
both uniform.  Hence every one of the nine cells has product weight \(1/9\),
including the three missing cells \((1,3),(2,1),(3,2)\).
\textit{Derived from (T1.5).}

Now take any nonconstant observed-edge wages and compare two complete schedules
that agree on all six observed edges but differ at \((1,3)\).  Every observed
wage and every observed-match input to the BS estimator is identical under the
two schedules.  Nevertheless, equation (T1.4), with
\(\omega_0=1/9\), shows that the two schedules generally imply different
values of \(C^w_{\mathrm{assign}}\).  Thus having at least two matches for
every worker and firm is enough for the BS repeated-measurement construction
under its auxiliary independence assumptions, but it is not enough to recover
an unrestricted missing pairwise wage cell.
\textit{The graph calculation is derived.  Source for the BS requirement
\(M_i,N_j\geq2\): BS20, p.\ 21.}

\subsection{1.6 What BS20 identifies from observed matches}

BS20 define
\[
\lambda(x)
=\int_{\mathcal Y}\int_0^1w(x,y,z)\,dz\,d\Phi_x(y),
\qquad
\mu(y)
=\int_{\mathcal X}\int_0^1w(x,y,z)\,dz\,d\Psi_y(x).
\tag{T1.6}
\]
The first integral averages a worker's wage over the firms and shocks that the
worker encounters under the observed matching process.  The second averages a
firm's wage over its observed worker distribution and shocks.  Their sorting
estimand is
\[
c_{\mathrm{BS}}
=\operatorname{Cov}_{P^{\mathrm{obs}}_{IJ}}
\bigl(\lambda(I),\mu(J)\bigr),
\qquad
\rho_{\mathrm{BS}}
=\frac{c_{\mathrm{BS}}}{\sigma_\lambda\sigma_\mu}.
\tag{T1.7}
\]
\textit{Source: BS20, pp.\ 5--7, equations (1)--(6).}

The types in (T1.6) are expectations, not directly observed characteristics.
BS20 estimate their second moments without requiring many observations for
each unit.  For a worker with \(M_i\) matches,
\[
\widehat\lambda_i
=\frac1{M_i}\sum_{m=1}^{M_i}w^w_{i,m},
\qquad
\widehat{\lambda_i^2}
=\frac{1}{M_i(M_i-1)}
\sum_{m=1}^{M_i}\sum_{m'\ne m}
w^w_{i,m}w^w_{i,m'}.
\tag{T1.8}
\]
Distinct-match cross-products remove the conditional match-error variance that
would contaminate a squared single-match wage.  The firm-side estimators are
the same construction with \(N_j\) and \(w^f_{j,n}\).
\textit{Source: BS20, pp.\ 22--23, equations (18) and (21), and
Propositions 3--4.  The contamination statement follows by expanding the
auxiliary mean-zero wage equations on p.\ 21.}

For a focal observed match \(m\) of worker \(i\) with firm \(j\), BS20 use
\[
\widehat c_{i,m}
=
\left(\frac{\sum_{m'\ne m}w^w_{i,m'}}{M_i-1}\right)
\left(\frac{\sum_{n'\ne n(i,m)}w^f_{j,n'}}{N_j-1}\right).
\tag{T1.9}
\]
Under their Auxiliary Assumption 3, the two other-match averages are
independent conditional on the relevant characteristics, so
\[
\mathbb E[\widehat c_{i,m}\mid x_i,y_j]
=\lambda(x_i)\mu(y_j).
\tag{T1.10}
\]
Duration-weighting (T1.9) over focal observed matches and subtracting
\(\widehat{\bar w}^{\,2}\) gives \(\widehat c_{\mathrm{BS}}\); combining this
with the two estimated type variances gives
\(\widehat\rho_{\mathrm{BS}}\).
\textit{Source: BS20, pp.\ 23--24, equations (23)--(25) and Proposition 5.}

This is ``observed-data only'' in a precise but limited sense.  The estimator
uses wages, durations, identities, and repeated observed partners.  It does
not query \(w(x,y,z)\) for a worker--firm pair that never forms.  Its
identification nevertheless uses substantive restrictions: at least two
matches per worker and firm; conditional i.i.d.\ mean-zero match errors;
restrictions relating duration and wage errors; cross-side independence of
other-match errors; and a law of large numbers across units.  BS20 explicitly
describe some of these assumptions as strong and design alternative samples
and an alternative estimator around them.
\textit{Source: BS20, pp.\ 21--24, Auxiliary Assumptions 1--4 and the
discussion following them.}

\subsection{1.7 Why the BS estimator does not recover the original product variance}

The project's product-weighted worker component is
\[
a_i
=\sum_jq_jm_{ij}-\mu.
\tag{T1.11}
\]
It averages worker \(i\)'s wage across the \emph{same} reference distribution
\(q\) for every worker.  A discrete version of the BS worker type is
\[
\lambda_i
=\sum_jP^{\mathrm{obs}}(J=j\mid I=i)m_{ij},
\tag{T1.12}
\]
after integrating any within-match shock.  It averages across the worker's
\emph{own observed conditional partner distribution}.  The weights in
(T1.11) and (T1.12) are generally different.  The firm-side distinction is
symmetric.
\textit{Source for (T1.11): LOO-v2, p.\ 11, equations (19)--(22).
Source for (T1.12): the discrete form of BS20, p.\ 5, equation (1).}

Moreover, the product variance requires
\[
\sum_{i,j}p_iq_jm_{ij}^2
\quad\text{and}\quad
\left(\sum_{i,j}p_iq_jm_{ij}\right)^2.
\tag{T1.13}
\]
Neither the worker conditional means \(\lambda_i\), the firm conditional means
\(\mu_j\), nor their observed-match covariance determines the pairwise
squares in (T1.13).  The six-edge construction in Section 1.5 proves this:
all BS observed-match inputs can be held fixed while a product-supported
missing cell, and therefore (T1.13), changes.
\textit{Equation (T1.13) is derived from the definition of variance; the
non-equivalence follows from Sections 1.4--1.5.}

There is also an estimand mismatch even when wages are additive.  BS20 show
that under a correctly specified AKM wage equation and linear conditional
expectations, the \emph{correlation} between BS types can equal the AKM
correlation because the BS types are linear transformations of the AKM
effects.  This does not state that the covariances are equal, and it does not
recover a nonadditive product variance.
\textit{Source: BS20, pp.\ 8--10, Proposition 1 and its discussion.}

\subsection{1.8 The exact sense in which BS can help}

Define the additive BS type-score schedule
\[
s_{ij}=\lambda_i+\mu_j-\bar w.
\tag{T1.14}
\]
Because the observed and product assignments have the same marginal
distributions of \(\lambda_I\) and \(\mu_J\),
\[
\begin{aligned}
C^w_{\mathrm{assign}}(s)
&=\frac12\left\{
\operatorname{Var}_{P^{\mathrm{obs}}_{IJ}}(\lambda_I+\mu_J)
-\operatorname{Var}_{P^\times}(\lambda_I+\mu_J)
\right\}\\
&=\operatorname{Cov}_{P^{\mathrm{obs}}_{IJ}}
(\lambda_I,\mu_J)
=c_{\mathrm{BS}}.
\end{aligned}
\tag{T1.15}
\]
\textit{Derived by expanding both variances.  The marginal variances cancel,
and independence under \(P^\times\) makes the product covariance zero.}

Equation (T1.15) explains exactly how BS avoid the missing-cell problem.  Once
\(\lambda_i\) and \(\mu_j\) are identified from repeated observed matches, the
constructed value \(s_{ij}\) is available for every recombination of worker
and firm types.  No unrestricted pair-specific missing value is needed.
However, this replaces the original schedule \(m_{ij}\) by the additive score
\(s_{ij}\).  It resolves the support problem by changing the target's wage
object, not by nonparametrically recovering the missing cells of \(m\).
\textit{Derived from (T1.14)--(T1.15) and the BS estimator in
Section 1.6.}

Therefore BS can contribute in three defensible ways.

\begin{enumerate}
\item \textbf{Alternative estimand.} Report
  \((c_{\mathrm{BS}},\rho_{\mathrm{BS}})\) as observed-match vertical sorting
  in wage types.  This statistic answers a coherent question without claiming
  to equal \(C^w_{\mathrm{assign}}(m)\).
  \textit{Source for the interpretation: BS20, pp.\ 5--7.}
\item \textbf{Exact benchmark for a constructed additive schedule.} Interpret
  \(c_{\mathrm{BS}}\) through (T1.15).  It is exactly the assignment
  contribution for \(s_{ij}\), so it is a useful observed-data benchmark
  beside the completion-dependent estimate for \(m_{ij}\).
  \textit{Derived in (T1.15).}
\item \textbf{Bridge model.} If an independently justified calibration shows
  that the maintained schedule coincides, up to a constant, with the
  constructed score \(s_{ij}\), then
  \(C^w_{\mathrm{assign}}(m)=c_{\mathrm{BS}}\).  This is a strong fixed-point
  restriction because the BS types are themselves conditional expectations
  generated by the observed matching process.  More flexible bridges can
  make \(m_{ij}\) a known or estimable low-dimensional function of BS types,
  but that bridge is an additional completion restriction and must be tested
  on observed cells.  BS20 alone does not impose it.
  \textit{The equality is derived from variance invariance to constants and
  (T1.15).  The warning follows from the nonidentification proof.}
\end{enumerate}

The original nonadditive target has four channels:
\[
\begin{aligned}
C^w_{\mathrm{assign}}(m)
={}&\operatorname{Cov}_{\mathrm{obs}}(a_I,b_J)
+\operatorname{Cov}_{\mathrm{obs}}(a_I,h_{IJ})
+\operatorname{Cov}_{\mathrm{obs}}(b_J,h_{IJ})\\
&+\frac12\left\{
\operatorname{Var}_{\mathrm{obs}}(h_{IJ})
-\operatorname{Var}_{P^\times}(h_{IJ})
\right\}.
\end{aligned}
\tag{T1.16}
\]
A single BS type covariance does not recover the two interaction-alignment
covariances or the change in interaction dispersion.  Hence substituting
\(c_{\mathrm{BS}}\) for the original \(C^w_{\mathrm{assign}}(m)\) would remove
the very nonadditive assignment channels that motivate the project.
\textit{Source: LOO-v2, p.\ 16, Proposition 3 and equation (43).}

\subsection{1.9 Final determination}

\begin{longtable}{p{0.39\textwidth}p{0.12\textwidth}p{0.40\textwidth}}
\textbf{Question} & \textbf{Answer} & \textbf{Reason} \\
\hline
Are \(P_I\), \(P_J\), and \(P_I\otimes P_J\) observable or estimable? &
Yes. & They are determined by observed assignment marginals. \\
Is
\(\operatorname{Var}_{P_I\otimes P_J}(m_{IJ})\) identified on an incomplete
graph with unrestricted \(m\)? &
No. & Product weights fall on unobserved cells whose wage levels and squares
are unrestricted. \\
Is \(C^w_{\mathrm{assign}}(m)\) then identified? &
No, generically. & A missing-cell perturbation leaves all observed data fixed
but changes the target; its identified set is unbounded below without bounds
on missing wages. \\
Do two or more matches per worker and firm solve this problem? &
No. & They permit repeated-measurement estimation of unit-level BS types but
do not reveal every pair-specific counterfactual wage. \\
Does BS20 identify the original \(C^w_{\mathrm{assign}}(m)\)? &
No. & BS20 estimates the covariance/correlation of observed-match expected
wage types, a different estimand. \\
Can BS20 still help? &
Yes. & It supplies an observed-data alternative and exactly estimates the
assignment contribution for the constructed additive BS type-score schedule.
It can estimate the original target only with an explicit bridge from that
score to the full wage schedule. \\
\end{longtable}

The recommended empirical treatment is therefore to report the BS type
covariance or correlation as a separate observed-match benchmark, report the
support and model-propagated shares of the original
\(C^w_{\mathrm{assign}}(m)\), and never label the BS statistic as an estimator
of the original target unless the bridge restriction is stated and defended.
\textit{The separation recommendation follows from the two estimands and from
LOO-v2, p.\ 27, Sections 7.6--7.7.}

\section{Task 2: What is the BS object, what is estimated, and how did the later paper change the project?}

\subsection{2.1 Version convention and main answer}

There are two distinct Borovi\v{c}kov\'a--Shimer papers.

\begin{itemize}
\item \textbf{BS20:} \emph{High Wage Workers Work for High Wage Firms},
  February 13, 2020.  This is a reduced-form measurement paper with structural
  models used as validation laboratories.
\item \textbf{BS24:} \emph{Assortative Matching and Wages: The Role of
  Selection}, circulated as a Richmond Fed and NBER working paper in November
  2024.  The local project file is a later revision dated October 23, 2025.
  The authors' research page lists the paper as revise-and-resubmit at the
  \emph{Quarterly Journal of Economics}; it should therefore not be described
  as a published QJE article in this note.
\end{itemize}

\textit{Sources: BS20 title page; Richmond Fed Working Paper 24-14 and NBER
Working Paper 33184, November 2024; BS24 local revision title page; Katar\'ina
Borovi\v{c}kov\'a's research page.  The official 2024 working paper is
\url{https://www.richmondfed.org/-/media/RichmondFedOrg/publications/research/working_papers/2024/wp24-14.pdf}.}

The central answer is:

\begin{enumerate}
\item In BS20, the \textbf{object of interest and estimand} is
  \[
  \rho_{\mathrm{BS}}
  =
  \operatorname{Corr}_{P^{\mathrm{obs}}_{IJ}}
  \bigl(\lambda(I),\mu(J)\bigr),
  \]
  the correlation between the expected log-wage type of a worker and the
  expected log-wage type of the worker's observed employer.
  \textit{Source: BS20, pp.\ 1, 5--7, equations (1)--(6).}
\item The \textbf{estimator} is
  \[
  \widehat\rho_{\mathrm{BS}}
  =
  \frac{\widehat c_{\mathrm{BS}}}
  {\sqrt{\widehat\sigma_\lambda^2\widehat\sigma_\mu^2}},
  \]
  where the covariance and variances are estimated with distinct-match
  cross-products.  The estimator is a random function of the sample; it is not
  the economic object.
  \textit{Source: BS20, pp.\ 22--24, equations (18)--(25).}
\item The main object in BS24 is not a replacement estimator for
  \(\rho_{\mathrm{BS}}\).  The paper promotes the two-sided search model used
  as one laboratory in BS20 into the main economic model.  Its central object
  is the \textbf{selection mechanism} that maps potential meetings into
  observed matches and thereby jointly shapes wages and sorting.
  \textit{Source: BS20, pp.\ 10--15, Section 3.2; BS24, pp.\ 4--17,
  Sections 2--4.  The statement that one model is promoted from laboratory to
  main model is a derived cross-paper comparison.}
\item For policy, BS24 distinguishes the expected surplus of a random
  \emph{meeting} from the average surplus of an accepted \emph{match}.  The
  October 2025 revision expresses the worker-value versions as the intent to
  treat, \(\operatorname{ITT}_x(y)\), and average treatment on the treated,
  \(\operatorname{ATT}_x(y)\), and explicitly identifies the ITT as the
  parameter of interest.  Observed wage components are linked more closely to
  the ATT, not the ITT.
  \textit{Source: BS24, pp.\ 31--34, Section 8; BS24/25 revision,
  pp.\ 9--14, equations (17)--(28).}
\end{enumerate}

Thus the later paper is an intellectual continuation of the 2020 search-model
laboratory, but it is not a continuation of the 2020 estimator.  BS20 asks,
``How much observed-match sorting is there?''  BS24 asks, ``What structural
selection process can generate observed wage and sorting patterns, and what
policy object do observed wage effects recover?''  \textit{Derived synthesis
from the cited objectives and results.}

\subsection{2.2 Notation chart for Task 2}

The chart retains all Task 1 notation.  The meeting rate is written
\(\rho_{\mathrm{meet}}\) here to prevent confusion with the BS20 sorting
correlation \(\rho_{\mathrm{BS}}\).  BS24 writes the meeting-rate parameter as
\(\rho\).

\begin{longtable}{p{0.20\textwidth}p{0.52\textwidth}p{0.21\textwidth}}
\textbf{Symbol} & \textbf{Meaning} & \textbf{Source or status} \\
\hline
\(x,y\) & Ordered worker and firm characteristics/types in BS24.  In BS20
these may be vector-valued characteristics before they are mapped into wage
types. & BS20, pp.\ 4--6; BS24, pp.\ 4--5. \\
\(\lambda(x),\mu(y)\) & BS20 expected log-wage worker and firm types.  These
remain the canonical symbols for the BS20 estimand. & BS20, p.\ 5,
equations (1)--(2). \\
\(\bar w\) & BS20 common mean log wage.  This scalar is distinct from the
BS24 type-specific reservation wage \(\bar w_x\). & BS20, p.\ 6,
equation (4). \\
\(\sigma_\lambda^2,\sigma_\mu^2\) & Variances of BS20 worker and firm wage
types. & BS20, pp.\ 6, 20, equations (5) and (16). \\
\(c_{\mathrm{BS}}\) & Covariance of BS20 wage types in observed matches. &
BS20, pp.\ 6, 20--21, equations (6) and (17). \\
\(\rho_{\mathrm{BS}}\) & BS20 sorting estimand
\(c_{\mathrm{BS}}/(\sigma_\lambda\sigma_\mu)\). & BS20, p.\ 5,
equation (3). \\
\(\rho_{\mathrm{AKM}}\) & Correlation between AKM worker and firm effects in
observed matches, used by BS20 as a comparator. & BS20, pp.\ 7--9. \\
\(\widehat{\bar w}\) & BS20 estimator of the common mean wage. & BS20,
p.\ 23, equation (19). \\
\(\widehat\sigma_\lambda^2,\widehat\sigma_\mu^2\) & BS20 estimators of the
two type variances. & BS20, pp.\ 23--24, equations (20) and (22). \\
\(\widehat c_{\mathrm{BS}},\widehat\rho_{\mathrm{BS}}\) & BS20 covariance and
correlation estimators. & BS20, p.\ 24, equations (24)--(25). \\
\(M_i,N_j\) & Numbers of distinct observed matches for worker \(i\) and firm
\(j\). & BS20, pp.\ 19--22. \\
\(m,m',n,n'\) & Focal and other-match indices in BS20.  They are not the
project wage schedule \(m_{ij}\). & BS20, pp.\ 19--23. \\
\(\tau\) & BS20 replication index governing the number of workers and firms;
the number of observations per unit may remain bounded as
\(\tau\to\infty\). & BS20, p.\ 18. \\
\(f_{x,y}\) & BS24 deterministic production scale for worker type \(x\) and
firm type \(y\). & BS24, pp.\ 5--7. \\
\(z,G,g\) & Match-specific productivity, its cumulative distribution, and
its density. & BS24, pp.\ 5--6. \\
\(z_0,\theta\) & Pareto scale and shape parameters.  The exact finite-mean
results use \(\theta>1\). & BS24, pp.\ 10--12. \\
\(r,\delta,\gamma\) & Discount rate, exogenous separation rate, and worker
Nash-bargaining share. & BS24, pp.\ 4--6. \\
\(\rho_{\mathrm{meet}}\) & Canonical symbol for the BS24 meeting-rate
parameter, written \(\rho\) in that paper. & BS24, p.\ 5. \\
\(u_x,v_y\) & Steady-state measures of unemployed type-\(x\) workers and
vacant type-\(y\) jobs. & BS24, pp.\ 5, 8--9. \\
\(\bar w_x,\bar\pi_y\) & Reservation wage of a type-\(x\) worker and
reservation profit of a type-\(y\) firm. & BS24, pp.\ 6--8. \\
\(\bar z_{x,y}\) & Match-acceptance productivity threshold
\((\bar w_x+\bar\pi_y)/f_{x,y}\). & BS24, p.\ 7, equation (6). \\
\(p_{x,y}\) & Canonical shorthand for the match-acceptance probability
\(1-G(\bar z_{x,y})\). & Derived from BS24 equation (6). \\
\(V^s_{x,y}(z)\) & Total surplus created by an \((x,y,z)\) match. & BS24,
pp.\ 6--7, equation (5). \\
\(W_{x,y}(z)\) & Equilibrium wage paid in an accepted
\((x,y,z)\) match. & BS24, pp.\ 7--8, equation (9). \\
\(\bar\phi_{x,y}\) & Steady-state measure of accepted matches between types
\(x\) and \(y\). & BS24, p.\ 8, equation (13). \\
\(\omega_{i,j}\) & Observed wage of worker \(i\) at firm \(j\) in the BS24/25
analytic wage equation. & BS24/25 revision, pp.\ 10--12. \\
\(\alpha_x^*,\psi_y^*\) & BS24/25 type-level worker and firm components of
accepted-match wages. & BS24/25 revision, pp.\ 11--13,
equations (21)--(22). \\
\(\varepsilon_{i,j}\) & Conditional mean-zero accepted-match wage residual in
the Pareto model. & BS24/25 revision, pp.\ 11--13,
equations (21)--(23). \\
\(\bar V^s_{x,y}\) & BS24 average surplus conditional on an \((x,y)\) meeting
becoming a match. & BS24, pp.\ 31--32, equation (33). \\
\(\bar V^{s,m}_{x,y}\) & BS24 expected surplus from a random \((x,y)\)
meeting before match quality is known.  The superscript \(m\) is the paper's
notation for ``meeting,'' not the project wage schedule. & BS24, pp.\ 31--32,
equation (32). \\
\(\operatorname{ATT}_x(y)\) & Worker-value gain from an \((x,y)\) meeting
conditional on acceptance in the 2025 revision. & BS24/25 revision,
pp.\ 9--11, equation (19). \\
\(\operatorname{ITT}_x(y)\) & Worker-value gain from creating a random
\((x,y)\) meeting before knowing whether it will be accepted. & BS24/25
revision, pp.\ 9--11, equations (17)--(18). \\
\(\alpha_i,\psi_j\) & Individual worker and firm fixed effects in the
conventional two-way fixed-effects regression.  These are estimators'
parameters, distinct from the type-level structural components
\(\alpha_x^*,\psi_y^*\). & BS24/25 revision, pp.\ 15--17,
equation (30). \\
\(\widehat\alpha_i,\widehat\psi_j\) & OLS estimates of the individual
two-way fixed effects. & BS24/25 revision, pp.\ 16--18. \\
\(B\) & Multiplicative attenuation factor for estimated firm-component
differences under endogenous separations in the 2025 revision. & BS24/25
revision, pp.\ 26--28, Proposition 7. \\
\end{longtable}

\subsection{2.3 BS20: the economic object, estimand, estimator, and estimate}

BS20 begin with a cross-section of observed matches.  They define
\[
\lambda(x)
=
\mathbb E[Y\mid \text{worker characteristic}=x,\text{ match observed}],
\tag{T2.1}
\]
\[
\mu(y)
=
\mathbb E[Y\mid \text{firm characteristic}=y,\text{ match observed}],
\tag{T2.2}
\]
where the paper writes the expectations as integrals over the conditional
partner and match-shock distributions.  \textit{Source: BS20, p.\ 5,
equations (1)--(2); equations (T2.1)--(T2.2) are conditional-expectation
restatements.}

The population target is
\[
\rho_{\mathrm{BS}}
=
\frac{
\mathbb E_{P^{\mathrm{obs}}_{IJ}}
[(\lambda(I)-\bar w)(\mu(J)-\bar w)]
}{
\sqrt{
\mathbb E[(\lambda(I)-\bar w)^2]\,
\mathbb E[(\mu(J)-\bar w)^2]
}
}.
\tag{T2.3}
\]
\textit{Source: BS20, pp.\ 5--7, equations (3)--(6).}

The hierarchy is therefore:

\begin{longtable}{p{0.24\textwidth}p{0.31\textwidth}p{0.35\textwidth}}
\textbf{Layer} & \textbf{Object} & \textbf{Role in BS20} \\
\hline
Economic concept & Sorting of high expected-wage workers into high
expected-wage firms & Question the paper wants to measure. \\
Estimand/object of interest & \(\rho_{\mathrm{BS}}\) & Fixed population
functional defined before seeing a particular sample. \\
Component estimands & \(\bar w,\sigma_\lambda^2,\sigma_\mu^2,c_{\mathrm{BS}}\)
& Moments needed to construct \(\rho_{\mathrm{BS}}\). \\
Estimator & \(\widehat\rho_{\mathrm{BS}}\), constructed from
\(\widehat{\bar w}\), \(\widehat\sigma_\lambda^2\),
\(\widehat\sigma_\mu^2\), and \(\widehat c_{\mathrm{BS}}\) & Random statistic
computed from the observed sample. \\
Estimate & A realized number such as the paper's preferred Austrian estimates
near \(0.47\) for men and \(0.43\) for women & Empirical realization of the
estimator, not the estimand itself. \\
Structural models & AKM, two-sided search, and discrete choice laboratories &
Used to test whether \(\rho_{\mathrm{BS}}\) moves with an economically intuitive
notion of sorting; not imposed in the estimator. \\
\end{longtable}

\textit{Sources: BS20, pp.\ 1--4 for the contributions and empirical
estimates; pp.\ 5--17 for the estimand and model laboratories; pp.\ 18--24 for
the estimator.}

Two nearby objects are not the main empirical target.  First,
\(\rho_{\mathrm{AKM}}\) is a comparator based on additive worker and firm
effects.  Second, the correlation between primitive productivity
characteristics \(x\) and \(y\) is available inside a structural laboratory
but is generally infeasible and sensitive to monotone changes in the units of
those characteristics.  BS20 choose wage types precisely to obtain a cardinal,
one-dimensional, wage-data-based sorting measure.
\textit{Source: BS20, pp.\ 5, 8--15, especially the discussion on p.\ 12.}

\subsection{2.4 The key insight in the BS20 consistency result}

The central statistical problem is an incidental-measurement problem.  The
number of workers and firms is large, but each worker and firm may have only a
small number of observed matches.  Therefore the individual sample averages
\(\widehat\lambda_i\) and \(\widehat\mu_j\) need not converge to their latent
types.  Naively squaring those noisy averages would add sampling-error
variance.  \textit{Source: BS20, pp.\ 18--24 and pp.\ 48--51.}

For two distinct worker matches \(m\ne m'\), the auxiliary equation and
conditional independence give
\[
w^w_{i,m}=\lambda_i+\varepsilon^w_{i,m},
\qquad
\mathbb E[
w^w_{i,m}w^w_{i,m'}\mid\lambda_i]
=\lambda_i^2.
\tag{T2.4}
\]
Thus
\[
\widehat{\lambda_i^2}
=
\frac{1}{M_i(M_i-1)}
\sum_m\sum_{m'\ne m}
w^w_{i,m}w^w_{i,m'}
\tag{T2.5}
\]
is unbiased for \(\lambda_i^2\), even when \(M_i\) remains small.  The same
construction gives an unbiased score for \(\mu_j^2\).
\textit{Source: BS20, p.\ 22, equation (18), and pp.\ 49--50, proof of
Proposition 3.  Equation (T2.4) is the direct expansion of their auxiliary
wage equation.}

For a focal observed worker--firm match, BS20 multiply the worker's average
wage in other matches by the firm's average wage in other matches.  Auxiliary
Assumption 3 makes those two leave-focal-match-out errors independent, so the
product is unbiased for \(\lambda_i\mu_j\).  This yields an unbiased match-level
score for the covariance numerator.
\textit{Source: BS20, p.\ 23, equation (23), and pp.\ 50--51, proof of
Proposition 5.}

The consistency argument then has four linear steps:

\begin{enumerate}
\item Construct noisy but unbiased scores for the first moments, second
  moments, and cross-moment of latent types using distinct matches.
\item Weight those scores by employment durations so the sample moments target
  the same cross-sectional population moments used in the estimand.
\item Replicate the economy: as \(\tau\to\infty\), the number of workers and
  firms grows while match counts and durations remain bounded.  Auxiliary
  Assumption 4 supplies the cross-unit independence needed for a law of large
  numbers.  The variance of each aggregate score is proportional to
  \(1/\tau\) and therefore vanishes.
\item Use consistency of ratios, differences, squares, and the continuous
  mapping in the correlation formula to obtain
  \(\widehat\rho_{\mathrm{BS}}\overset{p}{\to}\rho_{\mathrm{BS}}\), provided
  both type variances are positive.
\end{enumerate}

\textit{Source: BS20, pp.\ 48--52, proofs of Propositions 3 and 5; the four
steps are a linear reconstruction of those proofs.}

The key insight is therefore:
\[
\boxed{
\text{consistent aggregate latent moments do not require consistent
individual latent types.}
}
\tag{T2.6}
\]
The distinct-match products remove the finite-observation bias at the score
level; the large number of units removes the remaining score noise at the
aggregate level.  \textit{Derived synthesis of BS20 Propositions 3--5 and
Appendix B.1.}

This leave-out operation should not be confused with KSS leave-one-out
correction.  BS20 leave out the focal match to create conditional independence
between repeated measurements of latent types.  KSS leave-out methods remove
own-observation bias in quadratic functions of estimated fixed effects on a
mobility graph.  The algebraic idea of cross-fitting is related, but the
nuisance problem and estimand are different.
\textit{Derived comparison from BS20, pp.\ 22--24, and the KSS description in
BS20, pp.\ 38--39.}

\subsection{2.5 BS24: the structural mechanism and its hierarchy of objects}

BS24 take a worker type \(x\), a firm type \(y\), and a match shock \(z\).
The pair accepts a match exactly when
\[
z\geq\bar z_{x,y},
\qquad
\bar z_{x,y}
=
\frac{\bar w_x+\bar\pi_y}{f_{x,y}}.
\tag{T2.7}
\]
Accepted matches obey the structural wage equation
\[
W_{x,y}(z)
=
\bar w_x
+\gamma\bigl(zf_{x,y}-\bar w_x-\bar\pi_y\bigr).
\tag{T2.8}
\]
Their steady-state mass satisfies
\[
\delta\bar\phi_{x,y}
=
\rho_{\mathrm{meet}}u_xv_y
\bigl[1-G(\bar z_{x,y})\bigr].
\tag{T2.9}
\]
\textit{Source: BS24, pp.\ 6--8, equations (6), (9), and (13).  The meeting
rate has been renamed \(\rho_{\mathrm{meet}}\) only to avoid collision with
\(\rho_{\mathrm{BS}}\).}

Equations (T2.7)--(T2.9) are the core of the later paper.  The threshold
selects which potential meetings appear in administrative wage data; the wage
equation maps accepted productivity into pay; and the steady-state equation
maps pair-specific acceptance probabilities into the observed match
distribution.  One selection rule therefore generates both the accepted-wage
surface and assortative matching.
\textit{Source: BS24, pp.\ 5--17; the three-equation summary is also stated in
the October 2025 revision, p.\ 5.}

Under a Pareto match shock, the October 2025 revision obtains the exact
accepted-wage representation
\[
\omega_{i,j}
=
\alpha_{x_i}^*
+\psi_{y_j}^*
+\varepsilon_{i,j},
\qquad
\mathbb E[\varepsilon_{i,j}\mid x_i,y_j]=0.
\tag{T2.10}
\]
The worker and firm components are functions of reservation wages and profits,
while the residual comes from the selected match shock.  Importantly, the
production surface \(f_{x,y}\) cancels from the accepted-wage distribution
after conditioning on the threshold in the Pareto case.  Hence complementary
production can coexist with additive accepted wages.
\textit{Source: BS24/25 revision, pp.\ 11--13, Proposition 2 and
equations (21)--(23).}

For sorting, BS24 do not make a Pearson correlation the primitive target.
They derive conditions under which the full match matrix
\((\bar\phi_{x,y})\) has the monotone likelihood ratio order.  With strictly
increasing and weakly log-supermodular production, high-productivity workers
are relatively more likely to match with high-productivity firms.  Because
wage types are increasing in productivity under the same conditions, this
also implies positive matching between high-wage workers and high-wage firms.
\textit{Source: BS24, pp.\ 15--17, Propositions 4--5; BS24/25 revision,
pp.\ 19--20, Proposition 4.}

The later paper's object hierarchy is:

\begin{longtable}{p{0.27\textwidth}p{0.29\textwidth}p{0.34\textwidth}}
\textbf{Layer} & \textbf{Object} & \textbf{Status} \\
\hline
Structural primitives & \(f_{x,y},G,\gamma,r,\delta,\rho_{\mathrm{meet}}\) &
Assumed or calibrated ingredients of the model. \\
Selection mechanism & \(\bar z_{x,y}\) and \(p_{x,y}\) & Determines which
meetings become observed matches. \\
Observed wage interface & \(\alpha_x^*,\psi_y^*,\varepsilon_{i,j}\) or
average accepted-match wages & Equilibrium outcome visible in administrative
wage data. \\
Sorting implication & MLR ordering of \(\bar\phi_{x,y}\) & Structural theorem;
stronger than merely requiring a positive scalar correlation. \\
Selected policy object & Match surplus \(\bar V^s_{x,y}\), or worker-value
\(\operatorname{ATT}_x(y)\) in the revision & Conditions on acceptance and is
therefore closer to observed accepted wages. \\
Ultimate policy object & Meeting surplus \(\bar V^{s,m}_{x,y}\), or
worker-value \(\operatorname{ITT}_x(y)\) in the revision & Values a meeting
before knowing whether it becomes a match; this is the policy-relevant object. \\
Empirical estimators studied & Two-way fixed-effects OLS, BLM grouping, and
KSS variance-component corrections & Existing estimators whose structural
meaning and selection bias are analyzed; they are not a new estimator of the
meeting surplus or ITT. \\
\end{longtable}

\textit{Sources: BS24, Sections 2--8; BS24/25 revision, Sections 2--6.  The
layer labels are derived organization.}

\subsection{2.6 ATT, ITT, and what the wage estimator estimates}

The October 2025 revision makes the policy distinction especially explicit:
\[
\operatorname{ITT}_x(y)
=
\gamma\int_0^\infty V^s_{x,y}(z)g(z)\,dz,
\tag{T2.11}
\]
\[
\operatorname{ATT}_x(y)
=
\frac{
\gamma\int_{\bar z_{x,y}}^\infty
V^s_{x,y}(z)g(z)\,dz
}{
1-G(\bar z_{x,y})
}.
\tag{T2.12}
\]
Because surplus is zero when the meeting is rejected,
\[
\operatorname{ITT}_x(y)
=
p_{x,y}\operatorname{ATT}_x(y).
\tag{T2.13}
\]
\textit{Source: BS24/25 revision, pp.\ 9--11, equations (17)--(20).
Equation (T2.13) uses \(p_{x,y}=1-G(\bar z_{x,y})\).}

The ATT is selected: it asks how valuable the meeting was among pairs whose
shock was high enough to form a match.  The ITT values every induced meeting,
including rejected meetings whose realized gain is zero.  A policy that changes
meeting opportunities generally targets the ITT.  Administrative wage data
only contain accepted matches, so the acceptance probability \(p_{x,y}\) is
the missing link.
\textit{Source: BS24, pp.\ 31--34; BS24/25 revision, pp.\ 9--14.}

With exogenous separations and the exact additive wage equation, conventional
two-way fixed-effects estimates can recover differences in the firm wage
component under the paper's regularity and support conditions.  Those
differences map to differential ATT after discounting for \(r+\delta\).  They
do not recover differential ITT because the wage regression does not reveal
the pair-specific acceptance probability.
\textit{Source: BS24/25 revision, pp.\ 13--18, equations (24)--(31).}

The 2025 revision adds an important further warning.  When separations are
endogenous, movers are selected from a nonrepresentative part of the
within-match shock distribution.  Since firm effects are learned from movers,
even the firm-component difference and differential ATT are attenuated:
\[
\mathbb E[\widehat\psi_{j_1}-\widehat\psi_{j_2}]
=
B\bigl(\psi^*_{y_{j_1}}-\psi^*_{y_{j_2}}\bigr),
\qquad B<1.
\tag{T2.14}
\]
BLM grouping and KSS leave-out corrections address classification or
incidental-parameter bias, not this mobility-selection bias.
\textit{Source: BS24/25 revision, pp.\ 25--29, Proposition 7 and the
discussion following it.}

Therefore the later paper does not present one estimator whose estimand is the
paper's ultimate policy object.  It instead diagnoses a chain:
\[
\boxed{
\text{fixed-effects estimator}
\longrightarrow
\text{accepted-wage component}
\longrightarrow
\text{ATT under conditions}
\not\!\longrightarrow
\text{ITT without acceptance/meeting structure}.
}
\tag{T2.15}
\]
\textit{Derived synthesis of BS24/25 Sections 2--5.}

\subsection{2.7 The exact continuation from BS20 to BS24}

The connection has three steps.

\begin{enumerate}
\item BS20 Section 3.2 already contains random search, two-sided worker and
  firm heterogeneity, Nash bargaining, match-specific productivity shocks,
  endogenous acceptance thresholds, and a Pareto numerical exercise.
  Its purpose there is to ask whether \(\rho_{\mathrm{BS}}\) behaves sensibly
  inside a known structural economy.
  \textit{Source: BS20, pp.\ 10--15.}
\item BS24 makes that environment the paper's central model.  It derives
  analytical Pareto results for accepted wages and the match distribution,
  then adds on-the-job search, quantitative calibration, selection tests, and
  policy comparisons between meetings and matches.
  \textit{Source: BS24, pp.\ 4--35.}
\item Under the BS24 sufficient conditions, the full match matrix has the MLR
  order and the worker and firm wage types are increasing in productivity.
  Consequently the BS20 covariance and correlation are positive.  Thus BS24
  supplies one structural explanation for a positive BS20 estimand, but it
  does not change how BS20 estimates that estimand.
  \textit{Source: BS24, pp.\ 15--17; the final implication for
  \(\rho_{\mathrm{BS}}\) is derived.}
\end{enumerate}

The continuation is therefore:
\[
\boxed{
\text{BS20: define and estimate a model-robust sorting statistic}
\quad\longrightarrow\quad
\text{BS24: explain one structural mechanism behind that statistic and its
policy limits}.
}
\tag{T2.16}
\]

\subsection{2.8 What BS24 adds and what it gives up relative to BS20}

\begin{longtable}{p{0.21\textwidth}p{0.35\textwidth}p{0.35\textwidth}}
\textbf{Dimension} & \textbf{What BS24 adds} & \textbf{What is lost relative
to BS20} \\
\hline
Economic mechanism & A complete equilibrium account of selective hiring:
meetings, acceptance thresholds, bargaining, match formation, and sorting. &
The headline result is no longer robust across a broad class of unrelated
structural models. \\
Production and sorting & Sufficient conditions connecting productive
complementarity, MLR sorting, and high-wage/high-productivity rankings. &
Ordered one-dimensional types and restrictions such as monotonicity and
log-supermodularity replace BS20's wage-type reduction of possibly
multidimensional characteristics. \\
Wage equation & Exact or quantitative results showing why selected wages can
look additive or submodular despite complementary production. & Exact analytic
results rely heavily on the shock family, especially Pareto; accepted wages
can reveal little about the underlying production surface. \\
Policy interpretation & Separates the value of a random meeting from the value
of an accepted match; the revision labels the worker-value objects ITT and ATT.
& The policy-relevant meeting object is not identified from administrative
wages without random-search and acceptance-probability structure. \\
Mobility & Adds on-the-job search, selection diagnostics, and in the 2025
revision endogenous separations and an analytical fixed-effect attenuation
factor. & Existing BLM/KSS corrections do not solve the structural selection
bias, and the paper does not supply a comparably assumption-light replacement
estimator. \\
Empirical discipline & Calibrates the model to BLM and KSS wage decompositions
and shows why standard event-study diagnostics may miss selection. & The
empirical exercise is model-based rather than a direct nonparametric estimator
like \(\widehat\rho_{\mathrm{BS}}\). \\
Statistical contribution & Clarifies what conventional fixed-effect
estimators can and cannot estimate structurally. & It does not extend the
BS20 repeated-measurement consistency theorem to structural productivity,
meeting surplus, or ITT. \\
Scope & Can discuss production, surplus, policy, and equilibrium selection. &
Conclusions require substantially more economic structure and are less
portable as a semistructural interface. \\
\end{longtable}

\textit{Sources: BS20, pp.\ 1--5 and Sections 2--4; BS24, Introduction and
Sections 2--8; BS24/25 revision, Sections 2--6.  The gains/losses language is a
derived comparison, not a quotation from either paper.}

\subsection{2.9 Which paper is better for the semistructural project?}

The user's initial instinct is substantially correct, with one qualification.
BS20 is the better \emph{core semistructural measurement reference}; BS24 is
the better \emph{structural interpretation and warning reference}.

BS20 is closer to the semistructural project for four reasons:

\begin{enumerate}
\item It declares a reduced-form economic functional,
  \(\rho_{\mathrm{BS}}\), rather than treating one full structural model as
  true.
\item It shows that the same functional is meaningful in several model
  laboratories.
\item It supplies an implementable estimator using observed matches and a
  transparent set of repeated-measurement assumptions.
\item Its consistency theorem targets an aggregate functional even when
  individual latent effects remain noisy, which is close in spirit to the
  project's target-functional approach.
\end{enumerate}

\textit{Source: BS20, pp.\ 1--5 and Sections 3--4; the semistructural
assessment is derived.}

The qualification is that BS20 does not solve the project's missing-cell
problem.  As established in Task 1, it obtains robustness by changing the wage
object to worker and firm observed-match types.  It does not identify the
complete nonadditive pairwise schedule \(m_{ij}\) or
\(C^w_{\mathrm{assign}}(m)\).
\textit{Derived from Task 1 and BS20's definitions.}

BS24 belongs one layer to the right.  It is valuable because it proves that:

\begin{itemize}
\item an additive accepted-wage interface need not imply additive production;
\item observed assignment can be generated by productive complementarity
  operating through selection;
\item accepted-match wage effects can correspond to an ATT while failing to
  recover the policy-relevant ITT;
\item event-study symmetry and leave-out bias correction need not eliminate
  structural mobility-selection bias.
\end{itemize}

\textit{Source: BS24, Sections 3--8; BS24/25 revision, Sections 3--6.}

The recommended architecture is therefore:
\[
\boxed{
\begin{array}{c}
\text{Observed matched wages and mobility}\\
\downarrow\\
\text{Semistructural interfaces: project functionals and BS20 }
(c_{\mathrm{BS}},\rho_{\mathrm{BS}})\\
\downarrow\ \text{only under an explicit structural map}\\
\text{BS24 selection model: productivity, match surplus, ATT, and ITT}.
\end{array}
}
\tag{T2.17}
\]

BS20 should be used as a benchmark estimator and alternative sorting
functional.  BS24 should be used as one model in the admissible structural
class, as a proof that accepted-wage additivity is not economically
diagnostic, and as a sensitivity framework for policy interpretation.  The
project should not use BS24's random-search/Pareto structure to identify its
headline wage-schedule objects unless it is willing to become a substantially
more structural paper.
\textit{Derived recommendation from the comparison above and the
semistructural definition in the project handoff.}

\subsection{2.10 Final classification}

\begin{longtable}{p{0.30\textwidth}p{0.29\textwidth}p{0.31\textwidth}}
\textbf{Question} & \textbf{BS20} & \textbf{BS24/25} \\
\hline
Main economic question & How much do expected-wage worker and firm types sort
in observed matches? & How can selection jointly generate assortative matching
and apparently additive/submodular accepted wages, and what does this imply
for policy? \\
Main object of interest & \(\rho_{\mathrm{BS}}\). & Selection mechanism;
for policy, meeting surplus or ITT rather than accepted-match surplus or ATT. \\
Main estimator & New repeated-match estimator
\(\widehat\rho_{\mathrm{BS}}\). & No single new estimator of the ultimate
object; analyzes conventional fixed-effects, BLM, and KSS procedures. \\
What observed wages recover & Moments of expected observed-match wage types
under auxiliary independence assumptions. & Accepted-wage components and,
under additional conditions, differential ATT; not ITT by wages alone. \\
Structural commitment & Limited; models validate the measure. & High; the
search, selection, bargaining, shock, and steady-state structure drive the
interpretation. \\
Best role in this project & Semistructural benchmark and estimation logic. &
Structural nesting example, interpretation boundary, and selection-bias
warning. \\
\end{longtable}

\clearpage
\section{Task 3: The exact connection between BS and the project}

\subsection{3.1 Scope, order, and conclusions}

This section compares three objects in one direction:
\[
\boxed{
\begin{array}{c}
\text{BS24/25 structural primitives and equilibrium}\\
\Downarrow\\
\text{accepted matches and accepted wages}\\
\Downarrow\\
\text{project wage-schedule and assignment interface}\\
\Downarrow\\
\text{project functionals and the BS20 sorting functional}.
\end{array}}
\tag{T3.1}
\]
The arrows are maps, not claims that the objects are identical.  Following
them in this order prevents four common mistakes: equating production with
wages, equating meetings with accepted matches, equating the accepted-match
mean with a realized wage, and equating a population functional with its
estimator.

\textit{Source: BS20, Sections 2--4; BS24/25 revision, Sections 2--4;
LOO-v2, Sections 2--6.  The ordered map and the four warnings are derived.}

The conclusions are:

\begin{enumerate}
\item \textbf{BS20 is functionally contained in the project interface.}
  Whenever the project supplies an accepted-match conditional mean schedule
  \(m\) and an observed assignment law \(P^{\mathrm{obs}}_{IJ}\), it also
  supplies the BS20 population types and
  \(\rho_{\mathrm{BS}}\).  This is a functional-inclusion claim, not a claim
  that the project's static wage equation contains BS20's full repeated-match
  data-generating process.
\item \textbf{The later BS model is a structural generator of the project
  interface.}  Its search equilibrium generates both an accepted-match
  assignment distribution and an accepted-wage distribution.  Taking a
  conditional mean of the latter gives the project's maintained wage
  schedule.
\item \textbf{The exact Pareto additivity result is in wage levels.}
  Under the Pareto assumptions, the mean accepted wage in levels is additive,
  so its project interaction is exactly zero.  The original project studies
  log wages.  The mean accepted log-wage surface is generally nonadditive
  and, in BS, strictly submodular under the stated ordering conditions.
\item \textbf{The project can represent the finite BS mean log-wage surface
  exactly if its interaction rank is allowed to be large enough.}  A fixed
  low-rank specification is an additional empirical restriction; BS
  submodularity by itself does not imply low rank.
\item \textbf{Neither full model nests the other.}  The project does not
  contain BS's search, bargaining, acceptance, steady-state, and dynamic
  observation laws.  Conversely, the Pareto BS model does not generate every
  low-rank wage schedule admitted by the project.
\item \textbf{The closest BS analogue to the project's semistructural
  approach is BS20, not the entire later paper.}  BS20 defines a portable
  population functional and uses structural models as laboratories.  The
  later paper provides a valuable structural interpretation and stress test,
  but its conclusions about production, ATT, and ITT require the full search
  model.
\end{enumerate}

\textit{Sources: BS20, pp.\ 5--24; BS24/25 revision, pp.\ 5--18,
especially equations (6), (9), (13), and Propositions 2--3; LOO-v2,
pp.\ 9--21 and 37--44.  All nesting classifications are derived below.}

\subsection{3.2 Notation chart for Task 3}

This chart inherits all notation from Tasks 1 and 2 without changing it.  It
lists every new symbol used in Task 3 and repeats the old symbols needed to
read the map without looking backward.  In particular,
\(\rho_{\mathrm{meet}}\) remains the meeting rate and
\(\rho_{\mathrm{BS}}\) remains the sorting correlation.

\begin{longtable}{p{0.20\textwidth}p{0.52\textwidth}p{0.21\textwidth}}
\textbf{Symbol} & \textbf{Meaning} & \textbf{Source or status} \\
\hline
\(\mathcal I,\mathcal J;i,j\) & Project worker and firm sets; individual
worker and firm indices. & LOO-v2, p.\ 9. \\
\(I,J\) & Random worker and firm identities when used inside probabilities,
expectations, and variances. & LOO-v2, pp.\ 9--10. \\
\(k,\ell;x',y'\) & Dummy worker and firm indices, and dummy BS worker- and
firm-type indices, used only inside sums. & Derived indexing notation. \\
\(x_i,y_j\) & BS structural types of worker \(i\) and firm \(j\). &
BS24/25 revision, pp.\ 4--5. \\
\(x_1,x_2,y_1,y_2\) & Two ordered worker types and two ordered firm types used
to state the BS log-wage submodularity inequality. & BS24/25 revision,
pp.\ 17--18. \\
\(\mathcal X,\mathcal Y;X,Y\) & Finite worker- and firm-type sets; their
numbers of elements. & BS24/25 revision, pp.\ 4--5. \\
\(P^{\mathrm{obs}}_{IJ},W,w_{ij}\) & Observed individual assignment law, its
matrix, and cell probabilities. & LOO-v2, pp.\ 9, 40. \\
\(P_I,P_J\) & Worker and firm marginals of the observed individual assignment
law. & LOO-v2, pp.\ 10, 15. \\
\(p_i,q_j\) & Individual worker and firm marginals of
\(P^{\mathrm{obs}}_{IJ}\). & LOO-v2, pp.\ 10, 40. \\
\(P^\times\) & Project product benchmark
\(P_I\otimes P_J\), with cell mass \(p_iq_j\). & LOO-v2, p.\ 10. \\
\(Y,m_{ij}\) & Project log-wage outcome and maintained conditional mean log
wage for pair
\((i,j)\), including target pairs when a completion is imposed. & LOO-v2,
pp.\ 9--10. \\
\(\mu,a_i,b_j,h_{ij}\) & Product-weighted grand mean, worker component, firm
component, and double-centered interaction in
\(m_{ij}=\mu+a_i+b_j+h_{ij}\). & LOO-v2, p.\ 11. \\
\(\mathbf u_i,\mathbf v_j,\Lambda\) & Project worker and firm interaction
factor vectors and coefficient matrix in
\(h_{ij}=\mathbf u_i'\Lambda\mathbf v_j\).  Boldface prevents confusion with
BS unemployment and vacancy measures. & LOO-v2, pp.\ 21--23; boldface is a
notational clarification. \\
\(r_0\) & Maintained project interaction rank in Task 3.  It is distinct from
the BS discount rate \(r\). & Derived clarification of LOO-v2's \(r\). \\
\(Q_F,H_F,A_h,C^w_{\mathrm{assign}}\) & Project average within-worker
cross-firm dispersion, heterogeneity in firm contrasts, mean favorable-cell
assignment, and fixed-schedule variance assignment contribution. & LOO-v2,
pp.\ 12--17. \\
\(\lambda_i^{\mathrm{BS}},\mu_j^{\mathrm{BS}}\) & BS20 expected log-wage type
of individual worker \(i\) and firm \(j\).  Superscripts distinguish the firm
type from the project's scalar \(\mu\). & BS20, p.\ 5; superscripts are
derived clarification. \\
\(c_{\mathrm{BS}},\rho_{\mathrm{BS}}\) & BS20 covariance and correlation
between worker and firm wage types in observed matches. & BS20, pp.\ 5--7. \\
\(\bar w,\sigma_\lambda,\sigma_\mu\) & BS20 common observed mean log wage and
standard deviations of its worker and firm wage types.  The scalar
\(\bar w\) is distinct from the later paper's reservation wage
\(\bar w_x\). & BS20, pp.\ 5--7. \\
\(T_{\mathrm{BS}}(m,P^{\mathrm{obs}}_{IJ})\) & Derived name for the mapping
from a wage schedule and assignment law to \(\rho_{\mathrm{BS}}\). & Derived
notation. \\
\(\mathcal S_{\mathrm{interface}}\) & Ordered pair
\((m,P^{\mathrm{obs}}_{IJ})\), used as a name for the common empirical
interface. & Derived notation. \\
\(f_{x,y}\) & BS deterministic type-pair production scale. & BS24/25
revision, pp.\ 5--7. \\
\(z,G,g,z_0,\theta\) & Match-specific productivity draw, its CDF and density,
and the Pareto scale and shape parameters. & BS24/25 revision, pp.\ 5,
11--12. \\
\(t\) & A possible realized match-quality value used as the upper argument
of the truncated CDF. & Derived dummy value. \\
\(r,\delta,\gamma\) & Discount rate, exogenous separation rate, and worker
Nash-bargaining share.  In Task 3, \(\delta\) always has this BS meaning; the
Task 1 missing-cell perturbation called \(\delta\) was local to that proof. &
BS24/25 revision, pp.\ 4--6; collision clarification is derived. \\
\(\rho_{\mathrm{meet}}\) & BS meeting-rate parameter, denoted \(\rho\) in
the paper. & BS24/25 revision, p.\ 5; canonical renaming from Task 2. \\
\(u_x,v_y\) & Steady-state measures of unemployed type-\(x\) workers and
vacant type-\(y\) jobs. & BS24/25 revision, pp.\ 5, 8. \\
\(\bar w_x,\bar\pi_y\) & Type-\(x\) reservation wage and type-\(y\)
reservation profit. & BS24/25 revision, pp.\ 6--8. \\
\(\bar z_{x,y},p_{x,y}\) & Acceptance threshold and acceptance probability,
\(\bar z_{x,y}=(\bar w_x+\bar\pi_y)/f_{x,y}\) and
\(p_{x,y}=1-G(\bar z_{x,y})\). & BS24/25 revision, pp.\ 6--8; the \(p\)
shorthand is canonical from Task 2. \\
\(V^s_{x,y}(z),W_{x,y}(z)\) & Total match surplus and equilibrium accepted
wage for a realized \(z\). & BS24/25 revision, pp.\ 6--8. \\
\(\bar\phi_{x,y},\bar\Phi\) & Steady-state mass of accepted \((x,y)\)
matches and total accepted-match mass
\(\bar\Phi=\sum_{x,y}\bar\phi_{x,y}\). & First symbol: BS24/25 revision,
p.\ 8; total is derived. \\
\(\Pi^{\mathrm{obs}}_{x,y}\) & Observed joint type distribution among
accepted matches, \(\bar\phi_{x,y}/\bar\Phi\). & Derived from BS steady-state
match masses. \\
\(\pi_x,\kappa_y\) & Worker- and firm-type marginals of
\(\Pi^{\mathrm{obs}}_{x,y}\).  These are not the reservation profit
\(\bar\pi_y\). & Derived notation. \\
\(\omega_{i,j},\varepsilon_{i,j}\) & Realized accepted wage in levels and
its Pareto-model conditional mean-zero residual. & BS24/25 revision,
pp.\ 11--13. \\
\(\alpha_x^*,\psi_y^*\) & BS Pareto worker- and firm-type components of the
accepted wage in levels. & BS24/25 revision, pp.\ 11--13. \\
\(m^{L}_{ij}\) & Conditional mean accepted wage in levels for individual
pair \((i,j)\). & Derived from BS Proposition 2. \\
\(A_i,B_j,\bar A,\bar B\) & Individual lifts of the BS worker and firm
level-wage components and their product-weighted marginal means. & Derived
notation for the level-wage bridge. \\
\(\mu^L,a_i^L,b_j^L,h_{ij}^L\) & Canonical project decomposition applied to
the conditional mean accepted wage in levels. & Derived notation. \\
\(Q_F^L,H_F^L,A_h^L,C_{\mathrm{assign}}^{w,L}\) & Project functionals applied
to the level-wage schedule rather than the original log-wage schedule. &
Derived notation. \\
\(\ell_{x,y},m^{\log}_{ij}\) & Conditional mean accepted log wage in a type
cell and its individual-pair lift,
\(m^{\log}_{ij}=\ell_{x_i,y_j}\). & First object: BS24/25 revision,
pp.\ 17--18; individual lift is derived. \\
\(\ell_{x\cdot},\ell_{\cdot y},\ell_{\cdot\cdot}\) & Product-weighted
worker-type row mean, firm-type column mean, and grand mean of
\(\ell_{x,y}\), using \(\kappa_y\) and \(\pi_x\). & Derived canonical
centering. \\
\(M^\ell\) & Double-centered type-cell mean log-wage matrix with element
\(M^\ell_{x,y}=\ell_{x,y}-\ell_{x\cdot}-\ell_{\cdot y}
+\ell_{\cdot\cdot}\). & Derived from LOO-v2 canonical decomposition. \\
\(K_\ell,L_\ell,\Sigma_\ell,R_\ell\) & Rank and thin singular-value
decomposition of \(M^\ell\):
\(K_\ell=\operatorname{rank}(M^\ell)\) and
\(M^\ell=L_\ell\Sigma_\ell R_\ell'\). & Standard linear algebra; derived
application. \\
\(\mathbf u_i^\ell,\mathbf v_j^\ell\) & Individual project factors obtained
by lifting rows of \(L_\ell\Sigma_\ell^{1/2}\) and
\(R_\ell\Sigma_\ell^{1/2}\) from types to individuals. & Derived. \\
\(I_{K_\ell}\) & \(K_\ell\)-dimensional identity matrix used as the
interaction coefficient after the SVD factorization. & Standard notation. \\
\(K\) & An independent draw from the firm marginal \(P_J\) when \(J\) and
\(K\) index a pair of firms in \(H_F\). & LOO-v2, p.\ 13. \\
\end{longtable}

\subsection{3.3 The common empirical interface}

The project begins after a match has been observed.  Its basic outcome is the
conditional mean
\[
m_{ij}
=
\mathbb E[
  Y\mid I=i,J=j,\text{accepted match observed}
],
\tag{T3.2}
\]
and its second basic object is the probability
\[
w_{ij}
=
\Pr(I=i,J=j\mid\text{accepted match observed}).
\tag{T3.3}
\]
Equation (T3.2) is a wage-schedule object; equation (T3.3) is an assignment
object.  Neither equation says why a pair forms.

\textit{Source: LOO-v2, pp.\ 9--10 and 37--38.  The explicit
accepted-match conditioning aligns the LOO object with BS and is derived.}

This pair of objects,
\[
\mathcal S_{\mathrm{interface}}
\equiv
\bigl(m,P^{\mathrm{obs}}_{IJ}\bigr),
\tag{T3.4}
\]
is the exact interface shared by the project and the two BS papers:

\begin{itemize}
\item BS20 takes an observed matched economy and projects it to worker and firm
  expected-wage types.
\item BS24/25 starts before matching and structurally generates the
  distribution in (T3.4).
\item The project takes (T3.4), decomposes \(m\), completes it under a
  maintained model when support is missing, and evaluates declared
  functionals.
\end{itemize}

\textit{Source: BS20, Sections 2 and 4; BS24/25 revision, Sections 2--4;
LOO-v2, Sections 2--8.  Naming (T3.4) the common interface is derived.}

The interface is deliberately narrower than a full data-generating process.
It does not, by itself, specify rejected meetings, unemployment, vacancies,
profits, durations, repeat-match dependence, or counterfactual equilibrium
responses.  Those omissions are why functional inclusion at this interface
must not be called full structural nesting.
\textit{Derived from the objects present in BS24/25 equations (1)--(16) and
absent from LOO-v2 equation (14).}

\subsection{3.4 Step 1: the project decomposes the interface}

Fix the product reference distribution
\[
P^\times=P_I\otimes P_J,
\qquad
P^\times(I=i,J=j)=p_iq_j.
\tag{T3.5}
\]
The project defines
\[
\mu=\sum_i\sum_jp_iq_jm_{ij},
\tag{T3.6}
\]
\[
a_i=\sum_jq_jm_{ij}-\mu,
\tag{T3.7}
\]
\[
b_j=\sum_ip_im_{ij}-\mu,
\tag{T3.8}
\]
and
\[
h_{ij}=m_{ij}-\mu-a_i-b_j.
\tag{T3.9}
\]
Substitution of (T3.6)--(T3.8) into (T3.9) gives the equivalent double-centering
formula
\[
h_{ij}
=
m_{ij}
-\sum_kq_km_{ik}
-\sum_\ell p_\ell m_{\ell j}
+\sum_\ell\sum_kp_\ell q_km_{\ell k}.
\tag{T3.10}
\]
Thus
\[
m_{ij}=\mu+a_i+b_j+h_{ij}.
\tag{T3.11}
\]

\textit{Source: LOO-v2, p.\ 11, equations (19)--(23).  Equations
(T3.6)--(T3.10) write every substitution explicitly.}

The centering identities follow directly.  For a fixed worker \(i\),
\[
\sum_jq_jh_{ij}
=
\sum_jq_jm_{ij}
-\sum_jq_jm_{ij}
-\sum_jq_j\sum_\ell p_\ell m_{\ell j}
+\sum_\ell\sum_kp_\ell q_km_{\ell k}
=0.
\tag{T3.12}
\]
The last two terms cancel because
\[
\sum_jq_j\sum_\ell p_\ell m_{\ell j}
=
\sum_\ell\sum_jp_\ell q_jm_{\ell j}.
\]
The analogous column calculation gives
\[
\sum_ip_ih_{ij}=0.
\tag{T3.13}
\]
Also,
\[
\sum_ip_ia_i=0,
\qquad
\sum_jq_jb_j=0.
\tag{T3.14}
\]
These identities make \((\mu,a,b,h)\) unique for the declared product weights.

\textit{Source: LOO-v2, p.\ 11.  The displayed cancellations are derived
from the definitions.}

In the no-pair-covariate formulation, or after applying the same covariate
residualization used to define the target schedule, the maintained low-rank
model restricts the last term:
\[
h_{ij}
=
\mathbf u_i'\Lambda\mathbf v_j,
\qquad
\operatorname{rank}(\Lambda)=r_0.
\tag{T3.15}
\]
Here \(r_0\) is the maintained interaction rank; it is not the BS discount
rate \(r\).  The completed \(m\), rather than the primitive factors by
themselves, is then used to form \(Q_F,H_F,A_h\), and
\(C^w_{\mathrm{assign}}\).
If pair-varying covariates remain in \(m_{ij}\), their double-centered
contribution must also be included when computing the canonical \(h_{ij}\);
the primitive factor term and canonical interaction are not identical by
notation alone.
\textit{Source: LOO-v2, pp.\ 21--23, especially the qualification following
equation (65), and pp.\ 12--17.}

\subsection{3.5 Step 2: BS20 is a functional of that interface}

Assume \(p_i>0\) and \(q_j>0\).  Conditional probabilities in the observed
matched economy are
\[
\Pr(J=j\mid I=i)
=
\frac{w_{ij}}{p_i},
\qquad
\Pr(I=i\mid J=j)
=
\frac{w_{ij}}{q_j}.
\tag{T3.16}
\]
Therefore the BS20 individual wage types implied by the interface are
\[
\lambda_i^{\mathrm{BS}}
=
\sum_j\frac{w_{ij}}{p_i}m_{ij},
\tag{T3.17}
\]
\[
\mu_j^{\mathrm{BS}}
=
\sum_i\frac{w_{ij}}{q_j}m_{ij}.
\tag{T3.18}
\]
These are discrete versions of BS20's integrals defining
\(\lambda(x)\) and \(\mu(y)\).
\textit{Source: BS20, p.\ 5, equations (1)--(2); discrete formulas are
derived by conditional expectation.}

Let the common observed mean be
\[
\bar w
=
\sum_i\sum_jw_{ij}m_{ij}.
\tag{T3.19}
\]
By iterated expectations,
\[
\sum_ip_i\lambda_i^{\mathrm{BS}}
=
\sum_i\sum_jw_{ij}m_{ij}
=
\bar w
\tag{T3.20}
\]
and, symmetrically,
\[
\sum_jq_j\mu_j^{\mathrm{BS}}=\bar w.
\tag{T3.21}
\]
The BS covariance and variances are consequently
\[
c_{\mathrm{BS}}
=
\sum_i\sum_jw_{ij}
\bigl(\lambda_i^{\mathrm{BS}}-\bar w\bigr)
\bigl(\mu_j^{\mathrm{BS}}-\bar w\bigr),
\tag{T3.22}
\]
\[
\sigma_\lambda^2
=
\sum_ip_i\bigl(\lambda_i^{\mathrm{BS}}-\bar w\bigr)^2,
\qquad
\sigma_\mu^2
=
\sum_jq_j\bigl(\mu_j^{\mathrm{BS}}-\bar w\bigr)^2,
\tag{T3.23}
\]
and
\[
T_{\mathrm{BS}}(m,P^{\mathrm{obs}}_{IJ})
=
\rho_{\mathrm{BS}}
=
\frac{c_{\mathrm{BS}}}{\sigma_\lambda\sigma_\mu},
\tag{T3.24}
\]
when both variances are positive.

\textit{Source: BS20, pp.\ 5--7, equations (3)--(6).  Equations
(T3.19)--(T3.24) are the finite-interface form.}

This proves the functional-inclusion statement: once
\((m,P^{\mathrm{obs}}_{IJ})\) is specified, every population quantity on the
right side of (T3.24) is specified.  No additive or low-rank restriction is
needed to \emph{define} the BS functional.
\textit{Derived from (T3.16)--(T3.24).}

\subsubsection{Why BS types are not the project main effects}

Substitute (T3.11) into (T3.17):
\begin{align}
\lambda_i^{\mathrm{BS}}
&=
\sum_j\frac{w_{ij}}{p_i}
\bigl(\mu+a_i+b_j+h_{ij}\bigr) \notag\\
&=
\mu+a_i
+\sum_j\frac{w_{ij}}{p_i}b_j
+\sum_j\frac{w_{ij}}{p_i}h_{ij} \notag\\
&=
\mu+a_i
+\mathbb E_{\mathrm{obs}}[b_J+h_{iJ}\mid I=i].
\tag{T3.25}
\end{align}
The second line uses
\(\sum_jw_{ij}/p_i=1\).  Similarly,
\[
\mu_j^{\mathrm{BS}}
=
\mu+b_j
+\mathbb E_{\mathrm{obs}}[a_I+h_{Ij}\mid J=j].
\tag{T3.26}
\]

\textit{Derived by substituting the LOO-v2 canonical decomposition into the
BS20 definitions.}

Equations (T3.25)--(T3.26) show exactly what the project main effects omit from
the BS types.  A BS worker type includes the composition of firms employing
that worker and the worker's selected interaction cells.  A BS firm type
includes the composition of its workers and its selected interaction cells.
The project components \(a_i,b_j\), by contrast, average every row and column
using the common product weights \(q_j,p_i\).
\textit{Derived from (T3.7)--(T3.8) and (T3.25)--(T3.26).}

Under independent assignment, \(w_{ij}=p_iq_j\).  Then
\[
\sum_j\frac{w_{ij}}{p_i}b_j=\sum_jq_jb_j=0,
\qquad
\sum_j\frac{w_{ij}}{p_i}h_{ij}=\sum_jq_jh_{ij}=0.
\]
Hence
\[
\lambda_i^{\mathrm{BS}}=\mu+a_i.
\tag{T3.27}
\]
The symmetric calculation gives
\[
\mu_j^{\mathrm{BS}}=\mu+b_j.
\tag{T3.28}
\]
Thus independence is a transparent special case in which BS wage types and
the project main effects coincide up to their common mean.  Wage additivity
alone does not give (T3.27)--(T3.28) under dependent assignment, because the
conditional counterpart-composition terms can still vary.
\textit{Derived from the centering identities and BS20's observed-assignment
definition.}

\subsubsection{Why functional inclusion is not estimator inclusion}

BS20's estimator additionally needs repeated wages from distinct matches and
the auxiliary independence conditions used to make distinct-observation
cross-products unbiased.  The static interface (T3.4) contains neither a
panel observation law nor those independence conditions.  Therefore the
project contains the BS20 \emph{population estimand} as a functional, but it
does not automatically contain the BS20 estimator's sampling model or
consistency theorem.
\textit{Source: BS20, pp.\ 18--24, Auxiliary Assumptions 1--4 and
Propositions 2--5; derived scope distinction.}

\subsection{3.6 Step 3: the later BS model generates the interface}

The later BS model starts before a match exists.  For a meeting of worker type
\(x\) and firm type \(y\), flow output would be
\[
zf_{x,y}.
\tag{T3.29}
\]
The worker's and firm's outside flow values are \(\bar w_x\) and
\(\bar\pi_y\).  The net flow surplus is therefore
\[
zf_{x,y}-\bar w_x-\bar\pi_y.
\tag{T3.30}
\]
The match forms if and only if this expression is nonnegative.  Because
\(f_{x,y}>0\), dividing by \(f_{x,y}\) preserves the inequality:
\[
zf_{x,y}\geq\bar w_x+\bar\pi_y
\quad\Longleftrightarrow\quad
z\geq
\frac{\bar w_x+\bar\pi_y}{f_{x,y}}
\equiv\bar z_{x,y}.
\tag{T3.31}
\]
Thus
\[
p_{x,y}
=
\Pr(z\geq\bar z_{x,y}\mid x,y)
=
1-G(\bar z_{x,y}).
\tag{T3.32}
\]

\textit{Source: BS24/25 revision, pp.\ 6--7, equations (5)--(6).  The
inequality steps are written explicitly.}

Conditional on acceptance, Nash bargaining yields
\[
W_{x,y}(z)
=
\bar w_x
+\gamma\bigl(zf_{x,y}-\bar w_x-\bar\pi_y\bigr),
\qquad z\geq\bar z_{x,y}.
\tag{T3.33}
\]
This is the realized wage in an accepted match.  It is not the conditional
mean wage because \(z\) remains random above the threshold.
\textit{Source: BS24/25 revision, p.\ 7, equations (8)--(9).}

In steady state, inflows into the \((x,y)\) match stock equal outflows:
\[
\rho_{\mathrm{meet}}u_xv_yp_{x,y}
=
\delta\bar\phi_{x,y}.
\tag{T3.34}
\]
Solving for the stock gives
\[
\bar\phi_{x,y}
=
\frac{\rho_{\mathrm{meet}}}{\delta}
u_xv_yp_{x,y}.
\tag{T3.35}
\]
After normalization, the observed joint type distribution is
\[
\Pi^{\mathrm{obs}}_{x,y}
=
\frac{\bar\phi_{x,y}}
{\sum_{x'\in\mathcal X}\sum_{y'\in\mathcal Y}
\bar\phi_{x',y'}}.
\tag{T3.36}
\]
Its type marginals are
\[
\pi_x=\sum_y\Pi^{\mathrm{obs}}_{x,y},
\qquad
\kappa_y=\sum_x\Pi^{\mathrm{obs}}_{x,y}.
\tag{T3.36a}
\]

\textit{Source: BS24/25 revision, p.\ 8, equation (13).  Equations
(T3.35)--(T3.36a) are algebraic normalization and marginalization.}

The distribution of \(z\) among accepted \((x,y)\) matches is the lower-truncated
distribution
\[
\Pr(z\leq t\mid z\geq\bar z_{x,y},x,y)
=
\frac{G(t)-G(\bar z_{x,y})}
{1-G(\bar z_{x,y})},
\qquad t\geq\bar z_{x,y}.
\tag{T3.37}
\]
Combining this distribution with (T3.33) generates the accepted-wage
distribution in the cell.  Its log-wage conditional mean is
\[
\ell_{x,y}
=
\frac{
\displaystyle\int_{\bar z_{x,y}}^\infty
\log W_{x,y}(z)g(z)\,dz}
{1-G(\bar z_{x,y})}.
\tag{T3.38}
\]

\textit{Source: BS24/25 revision, pp.\ 8, 17--18, equations (14) and (31);
(T3.38) restates the paper's average log-wage definition.}

If the structural model treats individuals with the same \(x\) or \(y\) as
exchangeable, lift the type-cell mean to an individual pair:
\[
m^{\log}_{ij}=\ell_{x_i,y_j}.
\tag{T3.39}
\]
Equations (T3.36) and (T3.39), together with a within-type allocation rule,
produce the interface \((m,P^{\mathrm{obs}}_{IJ})\).  The within-type rule is
needed because BS determines type-cell masses directly, whereas the project
indexes individual workers and firms.
\textit{Derived map from the BS type model to the LOO individual interface.}

The complete direction of construction is now
\[
\begin{array}{c}
(f,G,r,\delta,\gamma,\rho_{\mathrm{meet}},\text{type populations})\\
\Downarrow\ \text{equilibrium}\\
(\bar w_x,\bar\pi_y,u_x,v_y)\\
\Downarrow\ \text{(T3.31)--(T3.32)}\\
(\bar z_{x,y},p_{x,y})\\
\Downarrow\ \text{(T3.34)--(T3.36)}\\
\Pi^{\mathrm{obs}}_{x,y}\\
\Downarrow\ \text{(T3.33), (T3.37)--(T3.39)}\\
(m^{\log},P^{\mathrm{obs}}_{IJ})\\
\Downarrow\\
\text{project and BS20 functionals}.
\end{array}
\tag{T3.40}
\]
Every downward arrow after the first is an explicit equation above.  The first
arrow requires solving the BS equilibrium reservation-value, steady-state,
and population-accounting equations.
\textit{Source for the equilibrium system: BS24/25 revision, pp.\ 5--9,
equations (1)--(16); the complete chain is derived.}

The chain also makes the relation between the two BS papers exact.  The later
model's implied BS20 worker type is
\[
\lambda^{\mathrm{BS}}(x)
=
\sum_y
\frac{\Pi^{\mathrm{obs}}_{x,y}}{\pi_x}
\ell_{x,y},
\qquad \pi_x>0,
\tag{T3.40a}
\]
and its implied BS20 firm type is
\[
\mu^{\mathrm{BS}}(y)
=
\sum_x
\frac{\Pi^{\mathrm{obs}}_{x,y}}{\kappa_y}
\ell_{x,y},
\qquad \kappa_y>0.
\tag{T3.40b}
\]
These are simply the conditional mean log wages among accepted matches.
Inserting (T3.40a)--(T3.40b) and the joint weights
\(\Pi^{\mathrm{obs}}_{x,y}\) into the covariance, variance, and correlation
formulas (T3.22)--(T3.24) produces the
\(\rho_{\mathrm{BS}}\) implied by the later structural equilibrium.
\textit{Source for the type definitions: BS20, p.\ 5, equations (1)--(2);
the substitution of later-model objects is derived.}

\subsection{3.7 Step 4: the Pareto level-wage result maps exactly into an additive project schedule}

Assume
\[
G(z)=1-\left(\frac{z}{z_0}\right)^{-\theta},
\qquad
\theta>1,
\tag{T3.41}
\]
and that every type-pair threshold is interior as in BS Proposition 2.  BS
show that the accepted wage in levels satisfies
\[
\omega_{i,j}
=
\alpha^*_{x_i}+\psi^*_{y_j}+\varepsilon_{i,j},
\tag{T3.42}
\]
with
\[
\mathbb E[
\varepsilon_{i,j}\mid x_i,y_j,\text{accepted match}
]=0.
\tag{T3.43}
\]
The components are
\[
\alpha_x^*
=
\left(1+\frac{\gamma}{\theta-1}\right)\bar w_x,
\qquad
\psi_y^*
=
\frac{\gamma}{\theta-1}\bar\pi_y.
\tag{T3.44}
\]

\textit{Source: BS24/25 revision, pp.\ 11--13, Proposition 2,
equations (21)--(23).}

Taking the conditional expectation of (T3.42) and applying (T3.43) yields
\[
m^L_{ij}
\equiv
\mathbb E[
\omega_{i,j}\mid i,j,\text{accepted match}
]
=
\alpha^*_{x_i}+\psi^*_{y_j}.
\tag{T3.45}
\]
This expectation step is the exact bridge from the realized BS wage to a
project-style conditional mean schedule in levels.
\textit{Derived directly from BS Proposition 2.}

Define
\[
A_i\equiv\alpha^*_{x_i},
\qquad
B_j\equiv\psi^*_{y_j},
\tag{T3.46}
\]
\[
\bar A\equiv\sum_ip_iA_i,
\qquad
\bar B\equiv\sum_jq_jB_j.
\tag{T3.47}
\]
The product grand mean of (T3.45) is
\begin{align}
\mu^L
&=
\sum_i\sum_jp_iq_j(A_i+B_j) \notag\\
&=
\left(\sum_ip_iA_i\right)\left(\sum_jq_j\right)
+\left(\sum_jq_jB_j\right)\left(\sum_ip_i\right) \notag\\
&=
\bar A+\bar B.
\tag{T3.48}
\end{align}
The product row component is
\begin{align}
a_i^L
&=
\sum_jq_j(A_i+B_j)-\mu^L \notag\\
&=
A_i+\bar B-(\bar A+\bar B) \notag\\
&=
A_i-\bar A.
\tag{T3.49}
\end{align}
The product column component is
\[
b_j^L=B_j-\bar B.
\tag{T3.50}
\]
Finally,
\begin{align}
h_{ij}^L
&=
m_{ij}^L-\mu^L-a_i^L-b_j^L \notag\\
&=
A_i+B_j-(\bar A+\bar B)
-(A_i-\bar A)-(B_j-\bar B) \notag\\
&=0.
\tag{T3.51}
\end{align}
This proves exact rank-zero interaction for the accepted conditional mean in
\emph{wage levels}.
\textit{Derived by applying the LOO canonical definitions to BS Proposition
2.}

The corresponding level-wage project functionals simplify to
\[
Q_F^L=\operatorname{Var}_{P_J}(B_J),
\qquad
H_F^L=0,
\qquad
A_h^L=0,
\tag{T3.52}
\]
and
\[
C_{\mathrm{assign}}^{w,L}
=
\operatorname{Cov}_{P^{\mathrm{obs}}_{IJ}}
(A_I-\bar A,B_J-\bar B).
\tag{T3.53}
\]
Equation (T3.53) follows because the project assignment decomposition has
only its common-component covariance when \(h^L=0\).
\textit{Source for the additive special case: LOO-v2, p.\ 17,
Section 4.5; substitution from (T3.49)--(T3.51) is derived.}

This algebra does \emph{not} say production is additive.  The production
surface \(f_{x,y}\) remains inside the acceptance threshold in (T3.31), hence
inside \(p_{x,y}\) and the observed match mass in (T3.35).  Pareto
scale-invariance removes \(f_{x,y}\) from the conditional accepted level-wage
distribution after selection; it does not remove \(f_{x,y}\) from assignment.
Therefore complementary production, strongly assortative assignment, and an
additive accepted level-wage mean can coexist.
\textit{Source: BS24/25 revision, pp.\ 11--14 and 18--22, Propositions 2 and
4; LOO-v2, pp.\ 37--38 and 42--44.}

\subsection{3.8 Step 5: the original log-wage project requires a different bridge}

LOO-v2 defines \(Y\) as log wage.  Applying a logarithm to (T3.42) does not
preserve additivity:
\[
\log\omega_{i,j}
=
\log\bigl(
\alpha^*_{x_i}+\psi^*_{y_j}+\varepsilon_{i,j}
\bigr),
\tag{T3.54}
\]
and in general
\[
\mathbb E[\log\omega_{i,j}\mid x_i,y_j,\text{accepted}]
\neq
\log\bigl(\alpha^*_{x_i}+\psi^*_{y_j}\bigr).
\tag{T3.55}
\]
The inequality reflects both the nonlinear logarithm and the
heteroskedastic residual.  Consequently, (T3.51) cannot be transferred from
levels to logs.
\textit{Source: BS24/25 revision, pp.\ 17--18, Section 3.4; the expectation
inequality is the standard nonlinear-transformation implication.}

BS instead define the type-cell mean log wage \(\ell_{x,y}\) in (T3.38).
Under their Pareto and monotonicity assumptions, it is increasing in both
types and strictly submodular:
\[
\ell_{x_1,y_2}+\ell_{x_2,y_1}
>
\ell_{x_1,y_1}+\ell_{x_2,y_2}
\quad
\text{when }x_1<x_2,\ y_1<y_2.
\tag{T3.56}
\]
Equivalently, the high-firm-type log-wage gain is smaller for the high worker
type:
\[
\ell_{x_2,y_2}-\ell_{x_2,y_1}
<
\ell_{x_1,y_2}-\ell_{x_1,y_1}.
\tag{T3.57}
\]
\textit{Source: BS24/25 revision, pp.\ 17--18, Proposition 3.}

To map this surface into the project without a gap, align the type marginals
in (T3.36a) with the individual marginals:
\[
\pi_x=\sum_{i:x_i=x}p_i,
\qquad
\kappa_y=\sum_{j:y_j=y}q_j.
\tag{T3.58}
\]
Each equality is the consistency condition between
the type-level assignment law and its individual-level lift.
Then define
\[
\ell_{x\cdot}=\sum_y\kappa_y\ell_{x,y},
\qquad
\ell_{\cdot y}=\sum_x\pi_x\ell_{x,y},
\tag{T3.59}
\]
\[
\ell_{\cdot\cdot}
=
\sum_x\sum_y\pi_x\kappa_y\ell_{x,y}.
\tag{T3.60}
\]
The double-centered interaction is
\[
M^\ell_{x,y}
=
\ell_{x,y}-\ell_{x\cdot}-\ell_{\cdot y}
+\ell_{\cdot\cdot}.
\tag{T3.61}
\]

\textit{Derived by applying LOO-v2 equations (19)--(23) at the BS type-cell
level.}

Now lift the type decomposition to individuals:
\[
\mu=\ell_{\cdot\cdot},
\tag{T3.62}
\]
\[
a_i=\ell_{x_i\cdot}-\ell_{\cdot\cdot},
\qquad
b_j=\ell_{\cdot y_j}-\ell_{\cdot\cdot},
\tag{T3.63}
\]
\[
h_{ij}=M^\ell_{x_i,y_j}.
\tag{T3.64}
\]
Adding (T3.62)--(T3.64) gives
\begin{align}
\mu+a_i+b_j+h_{ij}
&=
\ell_{\cdot\cdot}
+\ell_{x_i\cdot}-\ell_{\cdot\cdot}
+\ell_{\cdot y_j}-\ell_{\cdot\cdot} \notag\\
&\quad
+\ell_{x_i,y_j}-\ell_{x_i\cdot}
-\ell_{\cdot y_j}+\ell_{\cdot\cdot} \notag\\
&=
\ell_{x_i,y_j}
=m^{\log}_{ij}.
\tag{T3.65}
\end{align}
Every row-mean, column-mean, and grand-mean term cancels explicitly.  Thus the
project decomposition exactly reconstructs the BS accepted mean log-wage
surface.
\textit{Derived from (T3.59)--(T3.64).}

\subsubsection{The exact rank statement}

By construction,
\[
\sum_y\kappa_yM^\ell_{x,y}=0
\quad\text{for every }x,
\qquad
\sum_x\pi_xM^\ell_{x,y}=0
\quad\text{for every }y.
\tag{T3.66}
\]
If all type weights are positive, the nonzero vectors \(\kappa\) and \(\pi\)
are right and left null vectors.  Therefore
\[
K_\ell
\equiv
\operatorname{rank}(M^\ell)
\leq
\min\{X-1,Y-1\}.
\tag{T3.67}
\]
\textit{Derived from weighted double centering and the rank-nullity theorem.}

Take a thin singular-value decomposition
\[
M^\ell=L_\ell\Sigma_\ell R_\ell',
\tag{T3.68}
\]
where \(L_\ell\) has \(K_\ell\) columns, \(\Sigma_\ell\) is a
\(K_\ell\times K_\ell\) positive diagonal matrix, and \(R_\ell\) has
\(K_\ell\) columns.  Define individual factors
\[
\mathbf u_i^\ell
=
\bigl(L_\ell\Sigma_\ell^{1/2}\bigr)_{x_i,\cdot},
\qquad
\mathbf v_j^\ell
=
\bigl(R_\ell\Sigma_\ell^{1/2}\bigr)_{y_j,\cdot}.
\tag{T3.69}
\]
Then
\begin{align}
(\mathbf u_i^\ell)'\mathbf v_j^\ell
&=
\bigl(L_\ell\Sigma_\ell^{1/2}\bigr)_{x_i,\cdot}
\bigl(R_\ell\Sigma_\ell^{1/2}\bigr)'_{\cdot,y_j} \notag\\
&=
\bigl(L_\ell\Sigma_\ell R_\ell'\bigr)_{x_i,y_j} \notag\\
&=
M^\ell_{x_i,y_j}
=h_{ij}.
\tag{T3.70}
\end{align}
Thus the project model represents the finite BS mean log-wage surface exactly
with interaction rank \(K_\ell\), using \(\Lambda=I_{K_\ell}\).
\textit{Derived from the SVD.}

If the project fixes \(r_0<K_\ell\), it does not exactly contain that BS
surface.  It imposes an approximation or restriction by retaining only
\(r_0\) singular directions.  Conversely, submodularity in (T3.56) restricts
cross-difference signs but does not force \(K_\ell\) to be one or any other
small number.  Shape and rank are distinct restrictions.
\textit{Derived from (T3.56)--(T3.70).}

\subsection{3.9 Step 6: what each project's functional means inside the BS model}

Once BS has generated \((m^{\log},P^{\mathrm{obs}}_{IJ})\), the project can
calculate
\[
Q_F
=
\mathbb E_{I\sim P_I}
\left[
\operatorname{Var}_{J\sim P_J}(m^{\log}_{IJ})
\right],
\tag{T3.71}
\]
\[
H_F
=
\mathbb E_{J,K\stackrel{\mathrm{iid}}{\sim}P_J}
\left[
\operatorname{Var}_{I\sim P_I}
\bigl(m^{\log}_{IK}-m^{\log}_{IJ}\bigr)
\right],
\tag{T3.72}
\]
\[
A_h
=
\mathbb E_{P^{\mathrm{obs}}_{IJ}}[h_{IJ}],
\tag{T3.73}
\]
and
\[
C^w_{\mathrm{assign}}
=
\frac12\left\{
\operatorname{Var}_{P^{\mathrm{obs}}_{IJ}}(m^{\log}_{IJ})
-
\operatorname{Var}_{P_I\otimes P_J}(m^{\log}_{IJ})
\right\}.
\tag{T3.74}
\]
\textit{Source: LOO-v2, pp.\ 12--17, equations (25), (30), (42), and
(46).}

Their meanings remain at the accepted-wage interface:

\begin{enumerate}
\item \(Q_F\) asks how much the accepted mean log wage for the same worker
  varies across firms in the declared product reference population.
\item \(H_F\) asks how much workers differ in their accepted mean log-wage
  contrasts between the same two firms.
\item \(A_h\) asks whether observed matches concentrate in positive
  accepted-wage interaction cells.
\item \(C^w_{\mathrm{assign}}\) asks how observed dependence between worker
  and firm identities changes accepted systematic log-wage variance while
  holding the completed wage schedule fixed.
\end{enumerate}

\textit{Source: LOO-v2, pp.\ 12--17.  Inserting ``accepted'' states the BS
conditioning explicitly and is derived.}

None of these quantities is automatically BS production, total surplus, ATT,
or ITT.  In particular, (T3.74) holds the wage schedule fixed and replaces the
observed assignment with an independent coupling.  A BS policy counterfactual
changes meeting values, acceptance probabilities, unemployment, vacancies,
match masses, and possibly reservation values.  It therefore need not hold
either \(m\) or the marginals fixed.
\textit{Source: LOO-v2, pp.\ 15, 38; BS24/25 revision, pp.\ 8--14.
The contrast is derived from the two counterfactual definitions.}

The structural model nevertheless explains how a project statistic may arise.
For example, nonadditive \(f_{x,y}\) changes \(\bar z_{x,y}\), which changes
\(p_{x,y}\), which changes \(\bar\phi_{x,y}\), which changes
\(P^{\mathrm{obs}}_{IJ}\), which changes both
\(C^w_{\mathrm{assign}}\) and \(\rho_{\mathrm{BS}}\).  This is a mechanism for
the reduced-form statistics, not an identity equating them to productivity or
welfare.
\textit{Derived from equations (T3.31)--(T3.40), (T3.24), and (T3.74).}

\subsection{3.10 Step 7: exact model-level inclusion statements}

The word ``nest'' must specify both the object space and the direction of
inclusion.  The following statements are exact.

\begin{longtable}{p{0.28\textwidth}p{0.24\textwidth}p{0.38\textwidth}}
\textbf{Claim} & \textbf{Verdict} & \textbf{Reason} \\
\hline
The project contains the BS20 population sorting estimand. &
Yes, as a functional. &
Equations (T3.16)--(T3.24) construct
\(\rho_{\mathrm{BS}}\) from
\((m,P^{\mathrm{obs}}_{IJ})\). \\
The project contains the BS20 estimator and consistency theorem. &
No, not by the wage layer alone. &
The estimator also needs repeated distinct matches, duration and observation
rules, and BS20's auxiliary independence conditions. \\
The project contains the later BS accepted mean level-wage surface under
Pareto selection. &
Yes, after changing the declared outcome from log wages to wage levels. &
Equations (T3.45)--(T3.51) show exact additive, rank-zero representation. \\
The original log-wage project contains the later BS accepted mean log-wage
surface. &
Yes if \(r_0\geq K_\ell\); otherwise only under an additional low-rank
restriction or approximation. &
Equations (T3.61)--(T3.70) give the exact finite-rank representation. \\
The project contains the full later BS structural model. &
No. &
The static wage layer omits meetings, rejected draws, value functions,
bargaining, unemployment, vacancies, steady state, and policy equilibrium. \\
The later BS Pareto model contains every project low-rank wage surface. &
No. &
Its level conditional mean is additive and its log surface obeys additional
search, bargaining, Pareto, and shape restrictions. \\
The later BS model supplies a structural interpretation for every project
interaction. &
No. &
Project interactions can also arise from comparative advantage, outside
options, bargaining, amenities, or departures from Pareto selection. \\
\end{longtable}

\textit{Sources: BS20, Sections 2--4; BS24/25 revision, Sections 2--4;
LOO-v2, Sections 2--6 and Appendix F.  The verdicts are derived from the
explicit maps above.}

The safest single sentence is:
\[
\boxed{
\begin{array}{c}
\text{The project contains BS20's sorting estimand as an interface
functional and contains}\\
\text{the later BS finite accepted-mean wage layer at sufficient rank, but it
does not contain}\\
\text{either paper's complete dynamic, sampling, or structural model.}
\end{array}}
\tag{T3.75}
\]

\subsection{3.11 Step 8: BS's closest semistructural section versus ours}

Neither BS paper labels a formal section ``the semistructural section.''  Two
parts are nevertheless close enough to require comparison.

\subsubsection{BS20: models as laboratories}

BS20 first defines \(\rho_{\mathrm{BS}}\) from expected observed-match wages.
Section 3 then evaluates that same measure in several model laboratories:
correctly specified AKM, two-sided search with match-specific shocks, and a
discrete-choice model.  The models discipline the interpretation of an
already defined statistic; the statistic is not redefined for each model.
Section 4 then supplies a repeated-measurement estimator under auxiliary
statistical assumptions.
\textit{Source: BS20, pp.\ 5--18, Sections 2--3, and pp.\ 18--24,
Section 4.}

This is semistructural in the relevant sense:
\[
\text{portable economic functional}
\quad+\quad
\text{multiple structural laboratories}
\quad+\quad
\text{separate statistical measurement model}.
\tag{T3.76}
\]
The project follows the same broad discipline but retains a richer interface:
instead of reducing the matched economy immediately to one scalar correlation,
it keeps the pairwise wage schedule and assignment law long enough to measure
\(Q_F,H_F,A_h,C^w_{\mathrm{assign}}\), and
\(\rho_{\mathrm{BS}}\).
\textit{Derived comparison of BS20 Sections 2--4 with LOO-v2 Sections
3--4 and 12.}

\subsubsection{BS24/25: structural model projected onto empirical wage decompositions}

The later paper derives accepted-wage restrictions from a single search,
selection, and bargaining model and then studies what conventional two-way
fixed-effects, BLM, KSS, and mover designs recover inside that model.  This is
an empirical projection of a structural model, not an assumption-light
definition of a portable wage functional.  Its strength is that it can
distinguish accepted-match effects from meeting-level policy effects and can
explain selection bias.
\textit{Source: BS24/25 revision, Sections 2--6, especially pp.\ 11--18 and
26--31.}

The project's direction is the reverse:
\[
\text{observed accepted wages}
\ \longrightarrow\
\text{maintained mean schedule}
\ \longrightarrow\
\text{support-aware functionals}
\ \longrightarrow\
\text{optional structural interpretations}.
\tag{T3.77}
\]
It does not select one production, search, bargaining, or assignment mechanism
as the baseline truth.  Its strength is portability across structural
mechanisms; its cost is that wage data alone cannot identify production,
surplus, ATT, ITT, or equilibrium welfare.
\textit{Source: LOO-v2, pp.\ 37--38 and 42--44; derived comparison.}

\subsubsection{Side-by-side comparison}

\begin{longtable}{p{0.18\textwidth}p{0.25\textwidth}p{0.25\textwidth}p{0.23\textwidth}}
\textbf{Dimension} & \textbf{BS20} & \textbf{BS24/25} &
\textbf{Project} \\
\hline
Starting object &
Observed matched wages. &
Production, search, outside values, bargaining, and shocks. &
Accepted conditional mean wage schedule, observed graph, and assignment. \\
Primary target &
\(\rho_{\mathrm{BS}}\). &
Selection mechanism; ATT/ITT and structural interpretation. &
\(Q_F,H_F,A_h,C^w_{\mathrm{assign}}\), with
\(\rho_{\mathrm{BS}}\) available as an added functional. \\
Role of a wage model &
Not required to define the target. &
Derived from the structural equilibrium. &
Maintained low-rank conditional-mean completion. \\
Role of assignment &
Observed joint law used inside the types and their correlation. &
Endogenous equilibrium outcome of meetings and acceptance. &
Observed law compared with declared fixed-schedule benchmarks. \\
Missing cells &
Avoided for the scalar type functional by observed-match averaging. &
Filled at the type level by structural exchangeability and equilibrium
restrictions when every type pair has positive acceptance probability. &
Handled explicitly through target support, graph conditions, and low-rank
completion. \\
Causal content &
Limited without a structural laboratory. &
Can discuss production, meeting surplus, ATT, ITT, and equilibrium. &
Measures wage-schedule objects; structural labels require a separate bridge. \\
Estimator logic &
Distinct-match cross-products estimate aggregate latent moments. &
Studies conventional estimators inside a selected structural DGP. &
Leave-out estimation for target functionals after estimating a low-rank
schedule. \\
\end{longtable}

\textit{Sources: BS20, Sections 2--4; BS24/25 revision, Sections 2--6;
LOO-v2, Sections 3--12.  Table entries synthesize the cited sections.}

\subsection{3.12 Step 9: missing cells look different at the three layers}

The individual worker--firm graph can be sparse even when every BS type pair
has positive match probability.  The distinction is:
\[
\text{type-cell support}
\neq
\text{individual-pair support}.
\tag{T3.78}
\]
If \(p_{x,y}>0\) for all types, the structural model predicts a positive mass
for every type cell through (T3.35), provided \(u_xv_y>0\).  But a finite
administrative panel can still leave most individual pairs \((i,j)\)
unobserved.
\textit{Derived from BS equation (13) and the difference between type and
individual indices.}

BS fills the gap by imposing that individuals of the same structural type
share \(f_{x,y}\), thresholds, accepted-wage distributions, and match rates.
The project instead preserves individual heterogeneity in
\((a_i,b_j,\mathbf u_i,\mathbf v_j)\) and uses low rank plus graph overlap to
propagate information.  These are different completion assumptions.
\textit{Source: BS24/25 revision, Sections 2--4; LOO-v2, Sections 5--7;
derived comparison.}

Therefore BS does not make the Task 1 product/missing-cell issue disappear
without cost.  It replaces individual missing cells with type-level
exchangeability and a full structural equilibrium.  That can be useful as a
sensitivity model, but it changes the maintained model class and can change
the estimand.
\textit{Derived from Task 1 and the type-versus-individual map above.}

\subsection{3.13 Recommended division of labor in the project}

The three components should be used in the following order.

\begin{enumerate}
\item \textbf{Headline semistructural layer.}  Estimate the project's declared
  accepted-log-wage functionals with explicit support and rank diagnostics.
\item \textbf{Portable sorting benchmark.}  Compute the BS20
  \(T_{\mathrm{BS}}(m,P^{\mathrm{obs}}_{IJ})\) from the fitted interface and
  compare it with the direct BS20 estimator when the repeated-match
  assumptions are credible.
\item \textbf{Structural sensitivity layer.}  Use BS24/25 to ask whether the
  observed combination of assignment, accepted level-wage additivity, and
  accepted log-wage submodularity can be rationalized by Pareto selection.
\item \textbf{Interpretation boundary.}  Report production, ATT, ITT, or
  equilibrium welfare only after estimating or calibrating the additional BS
  primitives and defending random meetings, the shock distribution,
  bargaining, and equilibrium restrictions.
\end{enumerate}

\textit{Derived recommendation from the established inclusion and
interpretation results.}

The resulting conceptual ordering is
\[
\boxed{
\begin{array}{c}
\text{BS20 scalar observed-match measurement}\\
\subset\\
\text{project accepted-wage schedule and functional interface}\\
\subset\\
\text{BS24/25-style structural interpretation}
\end{array}}
\tag{T3.79}
\]
only in the sense of \emph{amount of maintained economic structure}.  It is
not literal set inclusion of full statistical models: BS20 has a repeated
sampling model absent from the static project, and the project admits wage
surfaces excluded by the Pareto BS model.
\textit{Derived scope statement.}

\subsection{3.14 Final answer}

The exact connection is now linear:

\[
\boxed{
\begin{array}{c}
\text{BS24/25 explains how meetings become accepted matches and wages;}\\
\text{its equilibrium therefore generates }(m,P^{\mathrm{obs}}_{IJ}).\\[0.4em]
\text{The project decomposes and, under maintained low rank, completes }m,\\
\text{then evaluates wage-dispersion and assignment functionals.}\\[0.4em]
\text{BS20 is one further projection of the same interface:}\\
\rho_{\mathrm{BS}}
=
T_{\mathrm{BS}}(m,P^{\mathrm{obs}}_{IJ}).
\end{array}}
\tag{T3.80}
\]

At the wage-layer level, the project exactly represents the later BS Pareto
mean wage in levels with \(h^L=0\), and exactly represents its finite mean
log-wage type table at rank
\(K_\ell\leq\min\{X-1,Y-1\}\).  At the full-model level, there is no nesting in
either direction.  BS20 is the closest semistructural companion; BS24/25 is
the structural mechanism and interpretation boundary.
\textit{Derived from Sections 3.3--3.13.}

\clearpage
\section{Task 4: Does BS extend to Kline and the other nested models?}

\subsection{4.1 The four questions and the answer}

For every model, this section answers the same four questions in the same
order:

\begin{enumerate}
\item \textbf{Population measure.}  Does the model generate the accepted or
  observed wage and assignment objects needed to define the BS20 wage-type
  correlation?
\item \textbf{Observed-data estimator.}  Do the model's assumptions imply the
  repeated-match conditions needed for the BS20 estimator, or are additional
  statistical assumptions required?
\item \textbf{Selection logic.}  Does the model contain a potential
  meeting/access event and an acceptance rule, so that the later BS
  meeting-versus-accepted-match logic and its ITT/ATT distinction are
  meaningful?
\item \textbf{Project nesting.}  Which layer of the model is contained in the
  project's low-rank wage interface, which layer is not, and what assumptions
  or model changes would close the gap?
\end{enumerate}

These are the four bullets used throughout Task 4.  They must remain separate.
A model can supply the BS population correlation without satisfying the BS
estimator assumptions; it can satisfy a low-rank wage equation without
containing a meeting/acceptance process; and it can generate observed wages
without being fully nested by the project's static interface.
\textit{Derived organizational rule from Tasks 1--3.}

The overall answer is:

\begin{enumerate}
\item The \textbf{BS20 population measure extends very broadly}.  Any model
  that generates a joint distribution of observed matches and square-integrable
  observed wages generates expected-wage worker and firm types and hence a
  model-implied \(\rho_{\mathrm{BS}}\).
\item The \textbf{BS20 estimator does not extend automatically}.  AKM, Kline,
  Tukey/Crippa, BLM, GKLP, Shimer--Smith, and Eeckhout--Kircher do not, merely
  by specifying their economic wage or mobility equations, imply all of
  BS20's repeated-measurement and cross-match independence assumptions.
\item The \textbf{later BS selection logic extends literally only to models
  with a meeting/access and acceptance margin}.  It applies directly to
  Shimer--Smith and Eeckhout--Kircher, although their acceptance is
  deterministic conditional on types.  It applies only by analogy to Kline,
  BLM, and GKLP unless those models are augmented with an offer/meeting and
  acceptance law.
\item The \textbf{project nests wage-equation layers, not every full model}.
  It exactly contains AKM, Tukey, BLM's static finite type-cell mean layer,
  GKLP perfect-information comparative advantage, and the leading
  Eeckhout--Kircher example.  It contains Shimer--Smith only when the relevant
  production interaction is low rank.  It does not contain Kline's complete
  unrestricted edge model without a separate transition layer.
\end{enumerate}

\textit{Sources: BS20, Sections 2--4; Kline, pp.\ 15--35; BLM,
pp.\ 702--706; GKLP, pp.\ 6--16; Shimer--Smith, pp.\ 346--349;
Eeckhout--Kircher, pp.\ 882--889; Crippa, Sections 2--3; LOO-v2,
pp.\ 21, 42--44.  Each conclusion is proved or qualified below.}

\subsection{4.2 Notation chart for Task 4}

Tasks 1--3 notation remains unchanged.  The chart defines every new symbol
used in Task 4.  Superscripts identify model-specific objects and prevent
collisions with the project's canonical \((\mu,a_i,b_j,h_{ij})\).

\begin{longtable}{p{0.20\textwidth}p{0.52\textwidth}p{0.21\textwidth}}
\textbf{Symbol} & \textbf{Meaning} & \textbf{Source or status} \\
\hline
\(\mathfrak M\) & A generic candidate economic or statistical model. &
Derived notation. \\
\(m^{\mathfrak M}_{ij},P^{\mathfrak M}_{IJ}\) & Conditional mean observed
wage schedule and joint observed-match assignment law generated by
\(\mathfrak M\). & Derived generic interface. \\
\(p_i^{\mathfrak M},q_j^{\mathfrak M}\) & Worker and firm marginals of
\(P^{\mathfrak M}_{IJ}\). & Derived. \\
\(\lambda_i^{\mathfrak M},\nu_j^{\mathfrak M}\) & BS expected-wage worker and
firm types implied by \(\mathfrak M\).  The symbol \(\nu\) replaces the BS
symbol \(\mu\) locally to avoid collision with the project grand mean. &
Derived translation of BS20, p.\ 5. \\
\(\rho_{\mathrm{BS}}^{\mathfrak M}\) & BS20 population correlation implied
by model \(\mathfrak M\). & Derived model-specific notation. \\
\(\rho_{\mathrm{BS}}^{\mathrm{data}}\) & BS20 population correlation in the
empirical matched-wage distribution, used as a model-validation target. &
Derived label. \\
\(\mathsf M\) & Indicator that a potential meeting/access event becomes an
observed match.  It is unrelated to the interaction matrices called \(M\) in
earlier sections. & Derived selection notation. \\
\(P^{\mathrm{meet}}_{IJ},s_{ij}\) & Distribution of potential meetings and
conditional acceptance probability
\(s_{ij}=\Pr(\mathsf M=1\mid I=i,J=j,\text{meeting})\). & Derived generic
selection notation. \\
\(d_{ij}^{\mathrm{obs}}\) & Expected duration, survival, or observation
weight of an accepted \((i,j)\) match.  It is common in a one-shot sample or
when all match-separation rates are equal. & Derived stock-sampling
correction. \\
\(G_{ij},\operatorname{ITT}_{ij},\operatorname{ATT}_{ij}\) & Worker gain from
a potential meeting, its mean before acceptance is known, and its conditional
mean given acceptance. & Derived generalization of BS24/25,
pp.\ 9--14. \\
\(D_{it},Y_{it}(j)\) & Firm or sector occupied by worker \(i\) at time \(t\),
and the potential log wage at firm/sector \(j\). & Kline, pp.\ 29--34;
GKLP, pp.\ 12--14. \\
\(X_t,\beta\) & Kline time-varying covariates and their coefficient in the
two-period adjusted wage change. & Kline, pp.\ 15, 18--19. \\
\(R_i,R\) & Covariate-adjusted wage change for mover \(i\), and its vector. &
Kline, pp.\ 18--19. \\
\(\mathcal E_F,|\mathcal E_F|\) & Directed firm-mobility edge set and its
number of edges.  This is distinct from the worker--firm observation graph
\(\mathcal E\) in Task 1. & Kline, pp.\ 16--19; subscript is derived. \\
\(E_{\mathrm{edge}}\) & Mover-by-edge indicator matrix, denoted \(E\) by
Kline. & Kline, p.\ 19; subscript is a clarification. \\
\(B_{\mathrm{inc}}\) & Firm-by-edge incidence matrix, denoted \(B\) by
Kline. & Kline, pp.\ 17--19; subscript is a clarification. \\
\(\Delta,\Delta_{jk}\) & Vector of unrestricted Kline edge effects and the
effect on the directed transition \(j\to k\). & Kline, p.\ 19. \\
\(u^{K}_i,u^K\) & Kline edge-regression residual and its vector.  The
superscript distinguishes it from project interaction factors. & Kline,
p.\ 19; superscript is derived. \\
\(c_{\mathrm{cyc}},\mathcal C,C_{\mathrm{cyc}},\dot\psi,\dot\eta\) & A cycle
vector, a directed cycle, matrix of fundamental cycle directions, projected
firm coefficients, and cycle coefficients in Kline's edge decomposition. &
Kline, pp.\ 18--20; subscripts clarify notation. \\
\(\Delta^{LR}_{jk},\Delta^K_{jk}\) & Edge effect implied by the project
low-rank level surface and an arbitrary target Kline edge effect. & Derived. \\
\(\bar{\mathbf u}_{jk}\) & Mean project worker interaction factor among
workers observed moving from \(j\) to \(k\). & Derived. \\
\(d_u,d_v\) & Dimensions of project worker and firm interaction factors. &
Derived project notation. \\
\(\tau_{jk},\xi_i\) & Unrestricted transition wedge and mover-level residual
in the proposed level-plus-edge supermodel. & Derived extension. \\
\(\alpha_i,\psi_j\) & Primitive additive project worker and firm effects in
the low-rank wage equation used in Task 4.  The canonical product-weighted
components remain \(a_i,b_j\); they need not equal these primitives when
pair-varying covariates remain. & LOO-v2, pp.\ 21--23. \\
\(\psi_j^{\mathrm{ATE}}\) & Population mean potential-wage level of firm
\(j\), relative to a reference firm, under Kline's no-selection-on-treatment-
effects restriction. & Derived notation from Kline, pp.\ 32--34. \\
\(\alpha_i^A,\psi_j^A,\varepsilon_{it}^A\) & AKM worker effect, firm effect,
and residual. & Kline, p.\ 15; superscripts are clarifications. \\
\(m^A_{ij},\lambda_i^A,\nu_j^A\) & AKM conditional mean wage and the BS
worker and firm wage types it implies. & Derived. \\
\(\alpha_i^T,\psi_j^T,\beta_0\) & Tukey worker coordinate, firm coordinate,
and scalar interaction parameter. & Crippa, Section 2.2. \\
\(m^T_{ij},\lambda_i^T,\nu_j^T\) & Tukey conditional mean wage and the BS
worker and firm wage types it implies. & Derived. \\
\(\bar\psi_i^{\,T},\bar\alpha_j^{\,T}\) & Observed-assignment conditional
means of the Tukey firm and worker coordinates for worker \(i\) and firm
\(j\). & Derived. \\
\(f_m,\alpha_i^S,\psi_j^S\) & Crippa seriation interaction function and its
scalar worker and firm coordinates. & Crippa, Section 2.3.2. \\
\(m^S_{ij}\) & Conditional mean outcome in the Crippa seriation model. &
Crippa, Section 2.3.2. \\
\(\zeta_{k\ell}\) & BLM static conditional mean wage for worker type \(k\)
and firm class \(\ell\). & Derived relabeling of BLM, pp.\ 702--704. \\
\(\Pi^B_{k\ell},\pi_k^B,\kappa_\ell^B\) & BLM observed type-class match
distribution and its worker-type and firm-class marginals. & Derived. \\
\(\lambda_k^B,\nu_\ell^B,\rho_{\mathrm{BS}}^B\) & BLM type-level BS worker
type, firm type, and sorting correlation. & Derived from BS20 and the BLM
interface. \\
\(\bar\zeta^B,\sigma_{\lambda,B},\sigma_{\nu,B}\) & BLM matched mean wage and
standard deviations of the two BLM-implied BS wage types. & Derived. \\
\(M^B,K_B,L_B\) & Double-centered BLM type-class interaction matrix and
numbers of worker types and firm classes. & Derived from BLM,
pp.\ 702--704. \\
\(\zeta_{k\cdot},\zeta_{\cdot\ell},\zeta_{\cdot\cdot}\) & Product-weighted
BLM type-row, class-column, and grand means. & Derived. \\
\(\alpha_k^B,a_\ell,b_\ell,\bar\alpha^B,\bar b\) & Worker-type value,
firm-class intercept and slope in BLM's simple regression, and the
corresponding weighted worker-type and slope means. & BLM, p.\ 704; weighted
means are derived. \\
\(X_{it},w_{ijt}\) & GKLP observed worker covariates and wage offered to
worker \(i\) by sector \(j\) at time \(t\). & GKLP, pp.\ 6--14. \\
\(y_{ijt},\psi_{ijt}\) & GKLP output and its log unobserved-productivity
component. & GKLP, pp.\ 6--7. \\
\(Z_i,\eta_i,\varepsilon_{ijt}^G\) & GKLP general ability, sector-specific
ability, and productivity shock. & GKLP, pp.\ 6--7. \\
\(\beta_j,b_j^G,c_j^G,\sigma_\varepsilon^2\) & GKLP sector-specific observed-
skill coefficient, return to unobserved skill, intercept, and productivity-
shock variance. & GKLP, pp.\ 6--13; superscripts clarify collisions. \\
\(m_{i,t-1}^G,\sigma_t^2,\sigma_{it}^2\) & GKLP posterior mean of
sector-specific skill before period \(t\), its common time-slice variance
term, and the state-augmented worker-time version used in the derived
embedding. & GKLP, pp.\ 8--14; superscripts clarify the project \(m_{ij}\). \\
\(\mathbf u_{it},I_2\) & Time-varying project worker factor used for the
GKLP learning slice and the two-dimensional identity matrix. & Derived. \\
\(\mathcal O_{it}\) & A proposed random sector/job offer or access set used
to augment GKLP with a literal meeting/access margin. & Derived extension. \\
\(m_{ijt}^{G,PI},m_{ijt}^{G,L}\) & GKLP perfect-information and learning
conditional mean log-wage surfaces. & Derived labels for GKLP equations
(9)--(10), pp.\ 12--13. \\
\(x,y,f^{SS}(x,y)\) & Shimer--Smith agent types and deterministic type-pair
flow output. & Shimer--Smith, pp.\ 346--349. \\
\(W^{SS}(x),w^{SS}(x)\) & Shimer--Smith unmatched asset value and its flow
value \(w^{SS}(x)=rW^{SS}(x)\). & Shimer--Smith, pp.\ 348--349. \\
\(S^{SS}(x\mid y),\pi^{SS}(x\mid y)\) & Shimer--Smith personal match surplus
and matched flow payoff to type \(x\). & Shimer--Smith, pp.\ 348--349. \\
\(\mathcal M^{SS}(x),a^{SS}_{xy}\) & Shimer--Smith mutually acceptable type
set and deterministic acceptance indicator. & Shimer--Smith,
pp.\ 347--349. \\
\(\Phi^{SS}_{xy}\) & Stationary observed joint distribution of
Shimer--Smith matched types, after normalization. & Derived from their
steady-state match stocks. \\
\(f^{EK}(x,y),w^{EK}(x,y),c_{\mathrm{search}}\) & Eeckhout--Kircher
production, accepted wage, and search cost. & Eeckhout--Kircher,
pp.\ 882--884; search subscript is a clarification. \\
\(\mathcal A_y^{EK},\mathcal B_x^{EK}\) & Eeckhout--Kircher worker and firm
acceptance sets, relabeled to avoid collision with Kline's incidence matrix. &
Eeckhout--Kircher, p.\ 882. \\
\(a_{\mathrm{EK}},h_{\mathrm{EK}},g_{\mathrm{EK}}\) & Coefficient and
worker- and firm-additive functions in the leading production example
\(f^{EK}(x,y)=a_{\mathrm{EK}}xy+h_{\mathrm{EK}}(x)+g_{\mathrm{EK}}(y)\). &
Eeckhout--Kircher, pp.\ 883--885; subscripts clarify notation. \\
\(\bar x,\bar y,h^{EK}(x,y)\) & Product-reference type means and the
double-centered Eeckhout--Kircher wage interaction. & Derived. \\
\end{longtable}

\subsection{4.3 A general extension theorem}

Suppose a model \(\mathfrak M\) generates an observed match
\((I,J)\sim P^{\mathfrak M}_{IJ}\) and an observed wage \(Y\) with a finite
second moment.  Define
\[
m^{\mathfrak M}_{ij}
=
\mathbb E_{\mathfrak M}[Y\mid I=i,J=j,\text{observed match}],
\tag{T4.1}
\]
\[
p_i^{\mathfrak M}=\sum_jP^{\mathfrak M}_{IJ}(i,j),
\qquad
q_j^{\mathfrak M}=\sum_iP^{\mathfrak M}_{IJ}(i,j).
\tag{T4.2}
\]
For positive marginals, define
\[
\lambda_i^{\mathfrak M}
=
\sum_j
\frac{P^{\mathfrak M}_{IJ}(i,j)}
{p_i^{\mathfrak M}}
m^{\mathfrak M}_{ij},
\tag{T4.3}
\]
\[
\nu_j^{\mathfrak M}
=
\sum_i
\frac{P^{\mathfrak M}_{IJ}(i,j)}
{q_j^{\mathfrak M}}
m^{\mathfrak M}_{ij}.
\tag{T4.4}
\]
The model-implied BS correlation is
\[
\rho_{\mathrm{BS}}^{\mathfrak M}
=
\operatorname{Corr}_{P^{\mathfrak M}_{IJ}}
\left(
\lambda_I^{\mathfrak M},
\nu_J^{\mathfrak M}
\right),
\tag{T4.5}
\]
provided both type variances are positive.
\textit{Source: BS20, pp.\ 5--7, equations (1)--(6); equations
(T4.1)--(T4.5) are the finite-model restatement.}

\paragraph{Population-extension proposition.}
Equations (T4.1)--(T4.5) prove that the BS20 population measure extends to
every model satisfying the stated moment and nondegeneracy conditions.  The
proof needs neither wage additivity, low rank, random matching, nor a
meeting/acceptance model.  It only needs the model's observed matched-wage
distribution.
\textit{Derived directly from the definitions.}

This proposition is deliberately about the population functional.  The BS20
observed-data estimator additionally requires repeated wages from distinct
matches, conditional stability of worker and firm wage types, suitable
independence of distinct-match errors and durations, cross-side leave-out
independence, and weak enough cross-unit dependence for a law of large
numbers.  These conditions are not consequences of (T4.1)--(T4.5).
\textit{Source: BS20, pp.\ 18--24, Auxiliary Assumptions 1--4 and
Propositions 2--5.}

\paragraph{General selection identity.}
Now suppose the model also generates a potential meeting distribution
\(P^{\mathrm{meet}}_{IJ}\), acceptance indicator \(\mathsf M\), and worker
gain \(G_{ij}\), with \(G_{ij}=0\) when the meeting is rejected.  Let
\[
s_{ij}
=
\Pr(\mathsf M=1\mid i,j,\text{meeting}),
\tag{T4.6}
\]
\[
\operatorname{ITT}_{ij}
=
\mathbb E[G_{ij}\mid i,j,\text{meeting}],
\tag{T4.7}
\]
\[
\operatorname{ATT}_{ij}
=
\mathbb E[G_{ij}\mid i,j,\text{meeting},\mathsf M=1].
\tag{T4.8}
\]
The law of total expectation gives
\begin{align}
\operatorname{ITT}_{ij}
&=
s_{ij}
\mathbb E[G_{ij}\mid i,j,\mathsf M=1]
+(1-s_{ij})
\mathbb E[G_{ij}\mid i,j,\mathsf M=0] \notag\\
&=
s_{ij}\operatorname{ATT}_{ij},
\tag{T4.9}
\end{align}
because the second conditional mean is zero.  In a one-shot match sample, or
when accepted matches have a common duration/separation rate, the observed
assignment law is proportional to
\[
P^{\mathfrak M}_{IJ}(i,j)
\propto
P^{\mathrm{meet}}_{IJ}(i,j)s_{ij}.
\tag{T4.10}
\]
More generally, a cross-sectional stock sample also weights matches by their
expected duration, survival, or observation frequency:
\[
P^{\mathfrak M}_{IJ}(i,j)
\propto
P^{\mathrm{meet}}_{IJ}(i,j)s_{ij}d_{ij}^{\mathrm{obs}}.
\tag{T4.10a}
\]
\textit{Source: BS24/25 revision, pp.\ 9--14; equations (T4.6)--(T4.10)
are the generic law-of-total-expectation version of its ITT/ATT and selection
identities.  Equation (T4.10a) is the derived stock-sampling correction; in
BS's baseline common-separation model its common duration factor cancels.}

Models without \(P^{\mathrm{meet}}_{IJ}\) and \(\mathsf M\) cannot use
(T4.9) literally.  Calling observed workers or movers ``treated'' does not
create a meeting-level ITT; the intervention and the rejection outcome must
first be defined.
\textit{Derived scope condition.}

\subsection{4.4 Kline first: the model we are especially interested in}

\subsubsection{4.4.1 Refresher: Kline's edge model}

For two-period movers, Kline defines the covariate-adjusted wage change vector
\[
R
=
Y_2-Y_1-(X_2-X_1)'\beta.
\tag{T4.11}
\]
Let \(E_{\mathrm{edge}}\) indicate the directed origin-destination edge
traversed by each mover.  Kline's unrestricted edge model is
\[
R=E_{\mathrm{edge}}\Delta+u^K,
\qquad
\mathbb E[u^K\mid E_{\mathrm{edge}}]=0.
\tag{T4.12}
\]
It gives every observed directed transition \(j\to k\) its own mean
\(\Delta_{jk}\).
\textit{Source: Kline, pp.\ 18--19.}

Let \(B_{\mathrm{inc}}\) be the firm-by-edge incidence matrix: an edge column
has \(-1\) at its origin and \(+1\) at its destination.  AKM restricts
\[
\Delta=B_{\mathrm{inc}}'\psi^A.
\tag{T4.13}
\]
For every cycle vector \(c_{\mathrm{cyc}}\) satisfying
\(B_{\mathrm{inc}}c_{\mathrm{cyc}}=0\),
\[
c_{\mathrm{cyc}}'\Delta
=
c_{\mathrm{cyc}}'B_{\mathrm{inc}}'\psi^A
=
(B_{\mathrm{inc}}c_{\mathrm{cyc}})'\psi^A
=0.
\tag{T4.14}
\]
Kline shows that, in a connected graph, these zero-cycle restrictions exhaust
the restrictions that AKM places on edge effects.  He writes
\[
\Delta
=
B_{\mathrm{inc}}'\dot\psi
+C_{\mathrm{cyc}}\dot\eta,
\tag{T4.15}
\]
where AKM sets \(\dot\eta=0\).
\textit{Source: Kline, pp.\ 18--20.}

\subsubsection{4.4.2 How the current project maps into Kline}

Suppress covariates and write the project's primitive low-rank wage surface as
\[
m_{ij}
=
\alpha_i+\psi_j+\mathbf u_i'\Lambda\mathbf v_j.
\tag{T4.16}
\]
For worker \(i\), the wage contrast between destination \(k\) and origin \(j\)
is
\[
m_{ik}-m_{ij}
=
\psi_k-\psi_j
+\mathbf u_i'\Lambda(\mathbf v_k-\mathbf v_j).
\tag{T4.17}
\]
Average (T4.17) over the workers observed moving on edge \(j\to k\).  With
\[
\bar{\mathbf u}_{jk}
\equiv
\mathbb E[
\mathbf u_i\mid D_{i1}=j,D_{i2}=k
],
\tag{T4.18}
\]
the project implies
\[
\Delta_{jk}^{LR}
=
\psi_k-\psi_j
+\bar{\mathbf u}_{jk}'\Lambda
(\mathbf v_k-\mathbf v_j).
\tag{T4.19}
\]
\textit{Derived from the project wage surface; source for the edge
conditional-mean object: Kline, p.\ 19.}

Around a directed cycle \(\mathcal C\),
\begin{align}
\sum_{(j,k)\in\mathcal C}\Delta_{jk}^{LR}
&=
\underbrace{
\sum_{(j,k)\in\mathcal C}(\psi_k-\psi_j)
}_{=0}
\notag\\
&\quad+
\sum_{(j,k)\in\mathcal C}
\bar{\mathbf u}_{jk}'\Lambda
(\mathbf v_k-\mathbf v_j) \notag\\
&=
\sum_{(j,k)\in\mathcal C}
\bar{\mathbf u}_{jk}'\Lambda
(\mathbf v_k-\mathbf v_j).
\tag{T4.20}
\end{align}
Thus the project can generate Kline cycle effects, but only through shared
worker factors, firm factors, and edge-specific mover compositions.
\textit{Derived from (T4.19).}

\subsubsection{4.4.3 Why unrestricted Kline is not currently nested}

There are two independent failures.

\paragraph{Object-space failure.}
The project maps an individual-firm pair to a wage level,
\[
(i,j)\longmapsto m_{ij},
\]
whereas Kline's unrestricted model maps a directed firm transition to the
selected movers' mean change,
\[
(j,k)\longmapsto\Delta_{jk}.
\]
An edge table alone does not specify stayer wages, cross-sectional wage levels,
or individual counterfactual wage cells.  It therefore does not determine the
project wage surface or the BS wage types.
\textit{Source for the two objects: LOO-v2, pp.\ 9--10; Kline, p.\ 19;
non-equivalence is derived.}

\paragraph{Shared-factor failure.}
Conditional on the mover-composition moments, (T4.19) has at most
\[
J+Jd_v+d_ud_v
\tag{T4.21}
\]
raw shared parameters before normalizations.  Kline's unrestricted edge table
has \(|\mathcal E_F|\) free entries.  When the edge count is large relative to
the fixed factor dimensions, a generic \(\Delta\) cannot satisfy (T4.19).
Increasing the factors until every edge can be fit removes the intended
low-rank restriction.
\textit{Derived dimension comparison; source for one parameter per edge:
Kline, p.\ 19.}

This parameter count is conditional on the mover-composition vectors
\(\{\bar u_{jk}\}_{(j,k)\in\mathcal E_F}\).  If every \(\bar u_{jk}\) is instead
treated as a free edge-specific parameter, those vectors contribute
\(|\mathcal E_F|d_u\) further degrees of freedom and can absorb much of the
unrestricted edge table.  That maneuver does not establish that Kline's edge
model is nested in the project's wage layer: it relocates the unrestricted
variation to an otherwise unmodeled edge-selection or mover-composition layer.
A disciplined factor model must therefore derive \(\bar u_{jk}\) from the
assignment and selection mechanism, or impose cross-edge restrictions on it.
\textit{Derived qualification to the dimension count.}

\subsubsection{4.4.4 Two rigorous routes to close the Kline gap}

\paragraph{Route A: add assumptions and restrict Kline.}
Kline's Assumption 4 imposes no selection on treatment effects:
\[
Y_{i2}(k)-Y_{i2}(j)
\ \perp\
(D_{i1},D_{i2})
\qquad
\text{for every }j\neq k.
\tag{T4.22}
\]
Together with Kline's exclusion, parallel-trends, stationarity, and
connectedness conditions, this gives
\[
\Delta_{jk}
=
\mathbb E[Y_{i2}(k)-Y_{i2}(j)]
=
\psi_k^{\mathrm{ATE}}-\psi_j^{\mathrm{ATE}}.
\tag{T4.23}
\]
The edge table is then a gradient of population wage levels, all cycle sums
are zero, and the rank-zero project/AKM layer nests it.
\textit{Source: Kline, pp.\ 29--34, Assumptions 1--4 and Proposition 2.}

A weaker restriction is to posit the factor-rationalizable edge equation
(T4.19) directly.  This permits cycle effects but asserts that they arise from
worker-factor composition and firm-factor differences.  It nests a structured
subclass of Kline, not Kline's unrestricted edge model.
\textit{Derived restriction; Kline, pp.\ 34--35, explicitly suggests
interactive factor models as a way to rationalize uncoarsened edge effects.}

\paragraph{Route B: change the project and keep Kline unrestricted.}
Add a transition layer:
\[
\Delta_{jk}
=
\psi_k-\psi_j
+\bar{\mathbf u}_{jk}'\Lambda(\mathbf v_k-\mathbf v_j)
+\tau_{jk},
\tag{T4.24}
\]
\[
R_i
=
\Delta_{D_{i1},D_{i2}}+\xi_i.
\tag{T4.25}
\]
For any desired Kline edge table \(\Delta^K\), choose
\[
\tau_{jk}
=
\Delta^K_{jk}
-\psi_k+\psi_j
-\bar{\mathbf u}_{jk}'\Lambda(\mathbf v_k-\mathbf v_j).
\tag{T4.26}
\]
Then (T4.24) reproduces \(\Delta^K\) exactly.  In particular, setting
\(\psi=0\) and \(\Lambda=0\) gives \(\tau_{jk}=\Delta^K_{jk}\).
\textit{Derived constructive nesting proof.}

The unrestricted \(\tau_{jk}\) must be called a \emph{transition wedge}, not
a time-invariant firm wage level.  If its sum around a cycle is nonzero, no
collection of firm levels can generate it.  Economically it may represent
origin dependence, switching costs, contracts, tenure/history, or selected
mover gains.  Giving it one of those structural interpretations requires an
appropriate dynamic potential-outcome or mobility model.
\textit{Derived from Kline's cycle characterization, pp.\ 18--20.}

The recommended supermodel is therefore
\[
\boxed{
\begin{array}{ll}
\text{Level layer:}
&
m_{ij}=\alpha_i+\psi_j+\mathbf u_i'\Lambda\mathbf v_j,\\[0.3em]
\text{Transition layer:}
&
\Delta_{jk}
=
\psi_k-\psi_j
+\bar{\mathbf u}_{jk}'\Lambda(\mathbf v_k-\mathbf v_j)
+\tau_{jk}.
\end{array}}
\tag{T4.27}
\]
Low rank remains a restriction on the wage-level interaction; complete Kline
nesting is obtained by leaving the transition residual unrestricted.
\textit{Derived recommendation.}

\subsubsection{4.4.5 The four BS questions for Kline}

\begin{enumerate}
\item \textbf{Population measure: not from the edge model alone.}
  The unrestricted \(\Delta\) records selected mover changes, not the
  cross-sectional wage levels and assignment frequencies needed in
  (T4.1)--(T4.5).  Two economies can have the same mover edge changes while
  having different stayer wages or different level distributions, and hence
  different \(\rho_{\mathrm{BS}}\).  Adding the level layer
  \((m,P^{\mathrm{obs}}_{IJ})\) makes the BS population measure well defined.
  \textit{Derived object-space argument.}
\item \textbf{Estimator: additional assumptions are required.}
  Kline's parallel-trends and edge-error conditions do not imply BS20's
  distinct-match repeated-measurement conditions.  Applying
  \(\widehat\rho_{\mathrm{BS}}\) requires the BS panel construction,
  at least two suitable distinct matches per worker and firm, conditional
  type stability, cross-match leave-out independence, and the correct
  cross-sectional match weighting.
  \textit{Source: Kline, pp.\ 29--35; BS20, pp.\ 18--24.}
\item \textbf{Selection logic: a close analogue, not a literal meeting ITT.}
  Under Kline's causal assumptions,
  \(\Delta_{jk}\) is a treatment effect for workers selected onto edge
  \(j\to k\).  Assumption 4 removes selection on treatment effects and converts
  it to the population effect in (T4.23).  This is analogous to an
  ATT-versus-ATE selection gap.  It becomes a literal BS ITT/ATT distinction
  only after defining a job-access or meeting intervention, rejection, and
  acceptance.
  \textit{Source: Kline, pp.\ 29--35; comparison to BS is derived.}
\item \textbf{Project nesting: use (T4.27).}
  AKM and the factor-rationalizable edge subclass are nested without
  \(\tau\).  The complete unrestricted Kline edge model is nested only after
  adding \(\tau_{jk}\), or after imposing assumptions such as (T4.22) that
  restrict Kline back to a level-rationalizable subclass.
  \textit{Derived from Sections 4.4.2--4.4.4.}
\end{enumerate}

\subsection{4.5 AKM and KSS}

The AKM conditional mean is
\[
m^A_{ij}=\alpha_i^A+\psi_j^A.
\tag{T4.28}
\]
It is the rank-zero case of the project: set \(\Lambda=0\).
\textit{Source: Kline, p.\ 15; LOO-v2, p.\ 42.}

The implied BS worker type is
\begin{align}
\lambda_i^A
&=
\mathbb E[m^A_{iJ}\mid I=i] \notag\\
&=
\alpha_i^A+\mathbb E[\psi_J^A\mid I=i],
\tag{T4.29}
\end{align}
and the implied firm type is
\[
\nu_j^A
=
\psi_j^A+\mathbb E[\alpha_I^A\mid J=j].
\tag{T4.30}
\]
Thus BS wage types are not generally the primitive AKM fixed effects.  They
also contain the observed composition of counterpart effects.  BS20 shows
that its correlation and the AKM correlation coincide in magnitude under the
additional linear-conditional-expectation conditions in its Proposition 1;
additivity alone is insufficient.
\textit{Source: BS20, pp.\ 8--10, Proposition 1; equations
(T4.29)--(T4.30) are derived.}

The four conclusions are:

\begin{enumerate}
\item \textbf{Population measure: yes.}  Combine (T4.28) with the observed
  assignment law and use (T4.29)--(T4.30).
\item \textbf{Estimator: not supplied by AKM/KSS.}  KSS corrects limited-
  mobility bias for AKM quadratic functionals.  It does not by itself impose
  the BS repeated-match independence conditions or estimate
  \(\rho_{\mathrm{BS}}\).
\item \textbf{Selection logic: absent from the wage equation.}  AKM strict
  exogeneity restricts wage errors and mobility but does not specify random
  meetings, rejected offers, or an acceptance probability.  A literal BS
  ITT requires adding those objects.
\item \textbf{Project nesting: exact rank zero.}  The AKM mean layer is
  exactly nested.  KSS is an inference method for that layer, not a different
  wage surface, so ``nesting KSS'' should mean preserving its leave-out
  inferential target rather than embedding a new economic model.
\end{enumerate}

\textit{Sources: Kline, Sections 2 and 4.2; Kline--Saggio--S{\o}lvsten as
summarized in LOO-v2, p.\ 42; BS20, Sections 2 and 4.  The distinctions are
derived.}

\subsection{4.6 Tukey/Crippa and the seriation boundary}

Crippa's Tukey mean surface is
\[
m^T_{ij}
=
\alpha_i^T+\psi_j^T+\beta_0\alpha_i^T\psi_j^T.
\tag{T4.31}
\]
After product centering, its interaction is an outer product and therefore
has rank at most one.  The project nests it exactly with one interaction
factor.
\textit{Source: Crippa, Section 2.2; LOO-v2, p.\ 42.}

Define the observed counterpart means
\[
\bar\psi_i^{\,T}
\equiv
\mathbb E[\psi_J^T\mid I=i],
\qquad
\bar\alpha_j^{\,T}
\equiv
\mathbb E[\alpha_I^T\mid J=j].
\tag{T4.32}
\]
Taking conditional expectations of (T4.31) gives the exact BS types
\[
\lambda_i^T
=
\alpha_i^T
+\left(1+\beta_0\alpha_i^T\right)
\bar\psi_i^{\,T},
\tag{T4.33}
\]
\[
\nu_j^T
=
\psi_j^T
+\left(1+\beta_0\psi_j^T\right)
\bar\alpha_j^{\,T}.
\tag{T4.34}
\]
The BS sorting statistic is the observed-match correlation of (T4.33) and
(T4.34).  It is not generally
\(\operatorname{Corr}(\alpha_I^T,\psi_J^T)\), because counterpart composition
and the interaction both enter the wage types.
\textit{Derived from the Tukey equation and BS20's definitions.}

The four conclusions are:

\begin{enumerate}
\item \textbf{Population measure: yes.}  Equations (T4.33)--(T4.34) give it
  directly for any observed assignment law.
\item \textbf{Estimator: additional BS assumptions are required.}  Crippa's
  cycle and graph conditions identify Tukey parameters; they do not imply
  that distinct wage observations are conditionally independent repeated
  measures of fixed BS wage types.
\item \textbf{Selection logic: absent.}  The bipartite interaction graph says
  which outcomes are observed but does not specify potential meetings,
  rejection, or acceptance.  Add an assignment/offer block to obtain a
  meeting-level ITT.
\item \textbf{Project nesting: exact rank one for Tukey.}  This is wage-layer
  inclusion only; Crippa's graph-identification and estimator theory are
  separate.
\end{enumerate}

\textit{Source: Crippa, Sections 2--4; BS20, Section 4; nesting and estimator
comparison are derived.}

Crippa's seriation model instead writes
\[
m^S_{ij}=f_m(\alpha_i^S,\psi_j^S),
\tag{T4.35}
\]
where \(f_m\) is increasing in both arguments.  It identifies latent
productivity rankings, under its strong graph condition, only up to strictly
increasing reparameterizations; cardinal latent differences are not
identified.
\textit{Source: Crippa, Sections 2.3.2 and 3.4, Proposition 1 and
Theorem 4.}

This creates an important split.  The BS observed-wage correlation can still
be estimated directly from repeated observed wages without assigning a
cardinal scale to \(\alpha_i^S\) or \(\psi_j^S\).  But seriation alone does not
identify a completed cardinal wage surface or the project's missing-cell
functionals.  A finite seriation surface is representable by the project at
its full double-centered rank; a fixed small rank is an additional
restriction, not a consequence of monotonicity.
\textit{Derived from BS20's observed-wage definition, Crippa's ordinal
identification result, and the SVD argument in Task 3.}

\subsection{4.7 BLM}

\subsubsection{4.7.1 Static type-class wage layer}

Let \(k\in\{1,\ldots,K_B\}\) index BLM worker types and
\(\ell\in\{1,\ldots,L_B\}\) index firm classes.  Let
\[
\zeta_{k\ell}
=
\mathbb E[
Y\mid\text{worker type }k,\text{ firm class }\ell
]
\tag{T4.36}
\]
be the static conditional mean, and let
\(\Pi^B_{k\ell}\) be the observed match probability of that cell.  Its
marginals are
\[
\pi_k^B=\sum_\ell\Pi^B_{k\ell},
\qquad
\kappa_\ell^B=\sum_k\Pi^B_{k\ell}.
\tag{T4.37}
\]
\textit{Source for BLM worker types, firm classes, and conditional earnings
distributions: BLM, pp.\ 702--704; the finite probability notation is
derived.}

The BLM-implied BS wage types are
\[
\lambda_k^B
=
\sum_\ell
\frac{\Pi^B_{k\ell}}{\pi_k^B}
\zeta_{k\ell},
\qquad
\pi_k^B>0,
\tag{T4.38}
\]
\[
\nu_\ell^B
=
\sum_k
\frac{\Pi^B_{k\ell}}{\kappa_\ell^B}
\zeta_{k\ell},
\qquad
\kappa_\ell^B>0.
\tag{T4.39}
\]
Define
\[
\bar\zeta^B
=
\sum_k\sum_\ell\Pi^B_{k\ell}\zeta_{k\ell}.
\tag{T4.40}
\]
Then
\[
\sigma_{\lambda,B}^2
=
\sum_k\pi_k^B(\lambda_k^B-\bar\zeta^B)^2,
\qquad
\sigma_{\nu,B}^2
=
\sum_\ell\kappa_\ell^B(\nu_\ell^B-\bar\zeta^B)^2,
\tag{T4.41}
\]
and
\[
\rho_{\mathrm{BS}}^B
=
\frac{
\displaystyle
\sum_k\sum_\ell
\Pi^B_{k\ell}
(\lambda_k^B-\bar\zeta^B)
(\nu_\ell^B-\bar\zeta^B)
}{
\sigma_{\lambda,B}\sigma_{\nu,B}
}.
\tag{T4.42}
\]
Equations (T4.38)--(T4.42) are an exact plug-in map from a fitted BLM
type-class matched distribution to the BS20 population correlation.
\textit{Derived from BS20, pp.\ 5--7, and BLM's static type-class objects.}

If the desired BS object is defined at individual firm identity rather than
firm class, (T4.39) is sufficient only when firms in a class share the same
relevant conditional wage distribution, or when the target is explicitly
class-level sorting.  Otherwise, within-class firm heterogeneity must be
restored before calculating the firm-identity BS statistic.
\textit{Derived scope condition from BLM's class construction,
pp.\ 702--704.}

\subsubsection{4.7.2 Exact wage-layer nesting}

Using the product weights \(\pi_k^B\kappa_\ell^B\), define the row, column, and
grand means
\[
\zeta_{k\cdot}
=
\sum_\ell\kappa_\ell^B\zeta_{k\ell},
\qquad
\zeta_{\cdot\ell}
=
\sum_k\pi_k^B\zeta_{k\ell},
\tag{T4.43}
\]
\[
\zeta_{\cdot\cdot}
=
\sum_k\sum_\ell
\pi_k^B\kappa_\ell^B\zeta_{k\ell}.
\tag{T4.44}
\]
The double-centered interaction is
\[
M^B_{k\ell}
=
\zeta_{k\ell}
-\zeta_{k\cdot}
-\zeta_{\cdot\ell}
+\zeta_{\cdot\cdot}.
\tag{T4.45}
\]
It satisfies
\[
\sum_\ell\kappa_\ell^BM^B_{k\ell}=0,
\qquad
\sum_k\pi_k^BM^B_{k\ell}=0.
\tag{T4.46}
\]
With positive weights, these give nonzero left and right null vectors, hence
\[
\operatorname{rank}(M^B)
\leq
\min\{K_B-1,L_B-1\}.
\tag{T4.47}
\]
Lifting type and class rows to individual workers and firms reproduces the
static BLM mean surface in the project at that rank.
\textit{Derived weighted double-centering proof; source for the BLM static
surface: BLM, pp.\ 702--704.}

For BLM's displayed simple regression
\[
\zeta_{k\ell}=a_\ell+b_\ell\alpha_k^B,
\tag{T4.48}
\]
double centering yields
\[
M^B_{k\ell}
=
(\alpha_k^B-\bar\alpha^B)
(b_\ell-\bar b),
\tag{T4.49}
\]
which is an outer product and has rank at most one.
\textit{Source for the regression: BLM, p.\ 704; equation (T4.49) is the
derived weighted-centering result proved in the BLM handoff.}

\subsubsection{4.7.3 The four BS questions for BLM}

\begin{enumerate}
\item \textbf{Population measure: yes.}  The exact class-level formula is
  (T4.42).  BLM can therefore be evaluated against the BS sorting moment even
  though BLM estimates a much richer wage distribution.
\item \textbf{Estimator: not automatically.}  One route is to calculate
  (T4.42) from the fitted BLM model.  A separate route is to use the BS20
  distinct-match estimator.  BLM's dynamic framework allows mobility to
  depend on current earnings and future earnings to depend on previous
  earnings and classes; this flexibility does not imply BS20's stationary,
  independent repeated-measurement conditions.
\item \textbf{Selection logic: assignment yes, meeting/acceptance no.}
  BLM models class assignment and mobility, including endogenous mobility in
  its dynamic version, but it does not decompose infrequent matches into
  infrequent meetings versus low acceptance probabilities.  A literal BS
  ITT requires an offer/meeting distribution and acceptance block.
\item \textbf{Project nesting: static mean layer only.}  Equations
  (T4.43)--(T4.47) give exact finite-rank inclusion.  The full BLM earnings
  distributions, latent-class estimation, type proportions, and dynamic
  mobility laws are not contained in the static project interface.
\end{enumerate}

\textit{Sources: BLM, pp.\ 699--706; BS20, Section 4; the four classifications
are derived.}

\subsection{4.8 Gibbons--Katz--Lemieux--Parent (GKLP)}

\subsubsection{4.8.1 Perfect-information comparative advantage}

GKLP write output as
\[
y_{ijt}
=
\exp(X_{it}'\beta_j+\psi_{ijt}),
\qquad
\psi_{ijt}
=
Z_i+b_j^G(\eta_i+\varepsilon_{ijt}^G)+c_j^G.
\tag{T4.50}
\]
\textit{Source: GKLP, pp.\ 6--7, equations (1)--(2).}

Under perfect information and competitive wage setting, their log wage is
\[
m_{ijt}^{G,PI}
=
X_{it}'\beta_j
+Z_i
+b_j^G\eta_i
+c_j^G
+\frac12(b_j^G)^2\sigma_\varepsilon^2.
\tag{T4.51}
\]
\textit{Source: GKLP, p.\ 12, equation (9).}

If the sector-specific observed-skill term \(X_{it}'\beta_j\) is included in
the project's covariate component, the remaining wage surface has the exact
rank-one representation
\[
\alpha_i=Z_i,
\qquad
\psi_j=c_j^G+\frac12(b_j^G)^2\sigma_\varepsilon^2,
\tag{T4.52}
\]
\[
\mathbf u_i=(\eta_i),
\qquad
\mathbf v_j=(b_j^G),
\qquad
\Lambda=(1).
\tag{T4.53}
\]
Substitution of (T4.52)--(T4.53) into the project equation reproduces every
term in (T4.51) after the covariate component.
\textit{Derived exact embedding of GKLP equation (9); LOO-v2, p.\ 43 states
the same rank-one classification.}

\subsubsection{4.8.2 Learning}

With learning, the log wage becomes
\[
m_{ijt}^{G,L}
=
X_{it}'\beta_j
+Z_i
+b_j^Gm_{i,t-1}^G
+c_j^G
+\frac12(b_j^G)^2\sigma_t^2.
\tag{T4.54}
\]
\textit{Source: GKLP, p.\ 13, equation (10).}

For a fixed time or experience slice in which \(\sigma_t^2\) is common across
workers, the last term is firm-specific and
\(b_j^Gm_{i,t-1}^G\) is rank one.  If the relevant posterior-variance state is
worker-time-specific, the slice is represented with two factors:
\[
\mathbf u_{it}
=
\begin{pmatrix}
m_{i,t-1}^G\\
\sigma_{it}^2
\end{pmatrix},
\qquad
\mathbf v_j
=
\begin{pmatrix}
b_j^G\\
\frac12(b_j^G)^2
\end{pmatrix},
\qquad
\Lambda=I_2.
\tag{T4.55}
\]
\textit{Derived slice embedding of GKLP equation (10).}

The pooled learning panel is not a static worker-firm surface with a
time-invariant \(\mathbf u_i\).  The posterior mean changes as output histories
arrive, and sector choice depends on the current belief.  Full nesting
therefore requires time-indexed factors and the Bayesian state-transition and
sector-choice laws, not merely a larger static rank.
\textit{Source: GKLP, pp.\ 8, 13--16; derived boundary statement.}

\subsubsection{4.8.3 The four BS questions for GKLP}

\begin{enumerate}
\item \textbf{Population measure: yes, after fixing the unit and time
  target.}  For a fixed \(t\), apply (T4.3)--(T4.5) to the observed joint
  distribution of workers and sectors and to (T4.51) or (T4.54).  If \(j\)
  denotes a sector rather than a firm, the result is worker--sector wage-type
  sorting, not automatically worker--firm sorting.
\item \textbf{Estimator: not across learning states without modification.}
  Under learning, the conditional expected wage of the same worker changes
  with \(m_{i,t-1}^G\).  Repeated wages are therefore not repeated noisy
  measurements of one fixed worker wage type.  One may apply BS within a
  common time/experience slice, residualize the modeled learning state, or
  estimate the model and calculate its implied correlation.  With a small
  fixed number of sectors, firm-side asymptotics also differ from BS20's
  many-firm design.
\item \textbf{Selection logic: Roy choice, not meeting acceptance.}  Under
  perfect information the worker selects the sector offering the highest
  wage; under learning the choice depends on the current posterior mean.
  This selects observed sector wages on comparative advantage, but there is
  no random meeting or rejected offer in the baseline model.  A literal BS
  extension can add a random accessible offer set \(\mathcal O_{it}\) and an
  acceptance/choice rule over that set.
\item \textbf{Project nesting: exact static slices, dynamic extension for the
  full model.}  Perfect information is rank one after handling
  \(X_{it}'\beta_j\); learning is rank one or two per declared slice.  The
  full pooled panel requires state-augmented, time-varying factors and a
  mobility/learning block.
\end{enumerate}

\textit{Sources: GKLP, pp.\ 6--16; Kline, pp.\ 32--35, notes that
comparative-advantage models generally violate no selection on treatment
effects; BS20, Section 4.  The BS extension classifications are derived.}

\subsection{4.9 Shimer--Smith}

\subsubsection{4.9.1 Wage/payoff map}

Shimer--Smith is a symmetric-agent search model.  When types \(x\) and \(y\)
match, they produce deterministic flow output \(f^{SS}(x,y)\).  Let
\(W^{SS}(x)\) be the unmatched asset value and
\(w^{SS}(x)=rW^{SS}(x)\) its flow value.  Their matched Bellman and Nash
surplus-splitting equations imply the wage-like flow payoff
\[
\pi^{SS}(x\mid y)
=
\frac12f^{SS}(x,y)
+\frac12w^{SS}(x)
-\frac12w^{SS}(y).
\tag{T4.56}
\]
\textit{Source: Shimer--Smith, pp.\ 348--349, equations (3)--(5);
(T4.56) is the derived rearrangement documented in the search handoff.}

After additive terms are removed, the interaction in
\(\pi^{SS}(x\mid y)\) is one half of the double-centered production surface.
Consequently,
\[
\operatorname{rank}
\bigl(
\text{payoff interaction}
\bigr)
=
\operatorname{rank}
\bigl(
\text{double-centered }f^{SS}
\bigr).
\tag{T4.57}
\]
The project nests the payoff layer exactly if that rank does not exceed the
maintained project rank.  Shimer--Smith supermodularity conditions discipline
matching patterns but do not imply low rank.
\textit{Source: Shimer--Smith, pp.\ 344--363; equation (T4.57) is derived
from (T4.56).}

\subsubsection{4.9.2 The four BS questions for Shimer--Smith}

\begin{enumerate}
\item \textbf{Population measure: yes after orienting the two sides.}  Given
  the stationary matched-type law \(\Phi^{SS}_{xy}\), treat one side as the
  worker side and the other as the firm side, use
  \(\pi^{SS}(x\mid y)\) as the wage-like outcome, and apply
  (T4.3)--(T4.5).  Because the original model is symmetric, this is a
  correlation of expected side-specific flow payoffs, not literally a
  worker wage and a firm's wage bill unless an employment interpretation is
  imposed.
\item \textbf{Estimator: additional sampling assumptions are required.}
  The equilibrium model supplies stationary types and matches but not the
  exact finite-panel observation process used in BS20.  The BS estimator
  additionally needs persistent identities, repeated distinct matches,
  suitable no-recall/independence conditions, duration treatment, and weak
  cross-agent dependence.
\item \textbf{Selection logic: yes, but deterministic by type pair.}
  Shimer--Smith acceptance is
  \[
  a^{SS}_{xy}
  =
  \mathbf 1\{
  f^{SS}(x,y)\geq w^{SS}(x)+w^{SS}(y)
  \}.
  \tag{T4.58}
  \]
  For a type pair with \(a^{SS}_{xy}=1\), every such meeting is accepted and
  its meeting-level ITT equals its ATT.  For
  \(a^{SS}_{xy}=0\), no accepted wage is observed.  Adding a match-specific
  draw \(z\) makes acceptance probabilistic within a type pair and creates
  the genuine BS ITT/ATT wedge.
\item \textbf{Project nesting: conditional wage-layer inclusion.}
  Equation (T4.57) is the exact condition.  Generic Shimer--Smith is not
  automatically low rank, and the project's static wage layer does not
  contain the equilibrium matching sets, unmatched density, value functions,
  or steady-state laws.
\end{enumerate}

\textit{Sources: Shimer--Smith, pp.\ 346--363; BS20, Section 4; BS24/25,
Sections 2--3; derived comparison.}

\subsection{4.10 Eeckhout--Kircher}

\subsubsection{4.10.1 Selection and the leading rank-one wage example}

Eeckhout--Kircher study a two-stage environment.  Workers and firms are
initially paired; they accept the first-stage match when its surplus exceeds
the value of waiting and paying search cost \(c_{\mathrm{search}}\).  Rejected
agents proceed to the second-stage frictionless allocation.  Their first-stage
acceptance sets are \(\mathcal A_y^{EK}\) and
\(\mathcal B_x^{EK}\).
\textit{Source: Eeckhout--Kircher, pp.\ 882--883, equations (11)--(13).}

For the leading production example
\[
f^{EK}(x,y)
=
a_{\mathrm{EK}}xy
+h_{\mathrm{EK}}(x)
+g_{\mathrm{EK}}(y),
\tag{T4.59}
\]
equal bargaining gives the accepted first-stage wage
\[
w^{EK}(x,y)
=
\frac{a_{\mathrm{EK}}}{4}x^2
-\frac{a_{\mathrm{EK}}}{4}y^2
+\frac{a_{\mathrm{EK}}}{2}xy
+h_{\mathrm{EK}}(x).
\tag{T4.60}
\]
\textit{Source: Eeckhout--Kircher, pp.\ 882--884; LOO-v2, p.\ 43,
Appendix F.5.}

Let \(\bar x\) and \(\bar y\) be the product-reference means of the two types.
Double centering removes the \(x\)-only and \(y\)-only terms and changes the
bilinear term into
\[
h^{EK}(x,y)
=
\frac{a_{\mathrm{EK}}}{2}
(x-\bar x)(y-\bar y).
\tag{T4.61}
\]
This is an outer product and therefore rank one.  The project exactly nests
this wage layer.
\textit{Derived from (T4.60).}

\subsubsection{4.10.2 The four BS questions for Eeckhout--Kircher}

\begin{enumerate}
\item \textbf{Population measure: yes.}  The two-stage process generates
  observed matches and wages, so (T4.3)--(T4.5) applies.  This measure should
  be interpreted carefully: Eeckhout--Kircher prove that wage data alone may
  fail to distinguish positive from negative productive sorting under a
  reversal of the unobserved firm-type order.  A BS wage-type correlation is
  therefore not a sign test for production sorting.
\item \textbf{Estimator: promising data structure, but not automatic.}
  Eeckhout--Kircher explicitly consider repeated observations and a long
  panel for workers and jobs.  That helps estimate expected wages, but their
  long-panel argument is not the same as BS20's few-observation
  cross-product consistency theorem.  Applying the BS estimator still
  requires distinct-match error and cross-side independence conditions.
\item \textbf{Selection logic: yes.}  The first-stage random pairing is a
  meeting and the acceptance sets define \(\mathsf M\).  Conditional on
  \((x,y)\), acceptance is deterministic in the baseline, so accepted type
  pairs have \(s_{xy}=1\) and first-stage meeting ITT equals ATT; rejected
  pairs have \(s_{xy}=0\) and proceed to the continuation allocation.  A
  match-specific shock would create \(0<s_{xy}<1\).
\item \textbf{Project nesting: exact for the leading example, not for the
  full search model.}  Equation (T4.61) is rank one.  More general
  Eeckhout--Kircher wage surfaces require their own rank check or
  approximation.  The two-stage search, continuation allocation, and
  identification of search costs and production are outside the static
  project wage layer.
\end{enumerate}

\textit{Sources: Eeckhout--Kircher, pp.\ 872--889, especially Propositions
3--4 and the repeated-observation discussion on pp.\ 888--889; BS20,
Section 4; derived nesting comparison.}

\subsection{4.11 Model-by-model decision table}

\begin{longtable}{p{0.16\textwidth}p{0.17\textwidth}p{0.19\textwidth}p{0.20\textwidth}p{0.19\textwidth}}
\textbf{Model} & \textbf{BS population measure} &
\textbf{BS estimator} & \textbf{Literal BS selection logic} &
\textbf{Project inclusion} \\
\hline
Kline unrestricted edges &
Not from \(\Delta\) alone; yes after adding wage levels and assignment. &
Needs BS repeated-match assumptions. &
Mover ATT/ATE analogue; literal ITT needs job access/meeting. &
Add \(\tau_{jk}\) for complete edge nesting, or restrict selection/edge
structure. \\
AKM/KSS &
Yes; types include counterpart composition. &
KSS inference does not replace BS assumptions. &
No meeting or acceptance block. &
Exact rank zero; KSS is inference, not a new surface. \\
Tukey/Crippa &
Yes, by (T4.33)--(T4.34). &
Graph identification does not imply BS independence. &
No meeting or acceptance block. &
Exact rank one for Tukey; seriation needs rank check. \\
BLM &
Yes, by the plug-in formula (T4.42). &
Use BLM plug-in or separately impose BS assumptions. &
Mobility selection, but no meeting-versus-acceptance decomposition. &
Exact static type-cell mean layer; not full distributions/dynamics. \\
GKLP &
Yes for a declared time/sector target. &
Learning breaks fixed-type repeated measurement unless sliced or modeled. &
Roy sector choice, not random meeting acceptance. &
Perfect information rank one; learning rank one/two per slice; full panel
dynamic. \\
Shimer--Smith &
Yes for oriented side-specific payoffs. &
Needs a panel sampling and independence layer. &
Yes; deterministic acceptance by type pair. &
Only when double-centered production is low rank; full search is outside. \\
Eeckhout--Kircher &
Yes, but not a sign test for productive sorting. &
Long panels help; BS independence still must be imposed. &
Yes; deterministic first-stage acceptance bands. &
Leading wage example exact rank one; full two-stage search is outside. \\
\end{longtable}

\textit{Sources: model sections above.  Every table entry is a compressed
restatement of its cited four-part analysis.}

\subsection{4.12 The clean project extension}

The comparison suggests a four-block architecture:
\[
\boxed{
\begin{array}{ll}
\text{Wage block \(W\):}
&
m_{ij}=\mu+a_i+b_j+\mathbf u_i'\Lambda\mathbf v_j,\\[0.3em]
\text{Assignment block \(G\):}
&
P^{\mathrm{obs}}_{IJ}(i,j)
\propto
P^{\mathrm{meet}}_{IJ}(i,j)s_{ij}d_{ij}^{\mathrm{obs}},\\[0.3em]
\text{Edge block \(E\):}
&
\Delta_{jk}
=
\psi_k-\psi_j
+\bar{\mathbf u}_{jk}'\Lambda(\mathbf v_k-\mathbf v_j)
+\tau_{jk},\\[0.3em]
\text{Sampling block \(S\):}
&
\text{repeated-match observation and independence conditions.}
\end{array}}
\tag{T4.62}
\]

Each block has one job:

\begin{enumerate}
\item \(W\) contains the static wage-equation cores of AKM, Tukey, BLM, GKLP,
  and restricted search-model wage surfaces.
\item \(G\) adds the meeting/acceptance distinction needed for BS-style ITT,
  assignment counterfactuals, and literal selection mechanisms.
\item \(E\) retains Kline's directed mover interface and allows unrestricted
  transition residuals without pretending they are wage levels.
\item \(S\) states the statistical assumptions under which the BS20
  observed-data estimator, project leave-out estimators, or alternative
  procedures are consistent.
\end{enumerate}

\textit{Derived synthesis from the model audit.}

The blocks should not be activated indiscriminately.  The core
semistructural paper can estimate \(W\), calculate the BS20 population
functional, and use \(E\) as a diagnostic.  The assignment block \(G\) is
needed only for claims about meetings, acceptance, ITT, production, or policy.
An unrestricted \(\tau\) is needed only if complete Kline edge inclusion is a
stated goal rather than a factor-rationalizable approximation.
\textit{Derived recommendation.}

\subsection{4.13 What BS does and does not buy across models}

The BS measure is valuable across all of these models for one common reason:
it supplies the same observed-match sorting functional after each model has
generated wages and assignment.  It can therefore be used as a
model-comparison moment:
\[
\rho_{\mathrm{BS}}^{\mathrm{data}}
\quad\text{versus}\quad
\rho_{\mathrm{BS}}^{\mathfrak M}
\tag{T4.63}
\]
for AKM, Tukey, BLM, GKLP, or a search model.
\textit{Derived application of the general extension theorem.}

It does not perform three stronger tasks:

\begin{enumerate}
\item It does not identify the project's missing wage cells or
  \(C^w_{\mathrm{assign}}(m)\); it changes the target to an observed-support
  wage-type correlation.
\item It does not identify production or productive sorting.  The same wage-
  type statistic can arise from comparative advantage, outside options,
  bargaining, selection, or different production orders.
\item It does not make its estimator valid under every nested model.  The
  repeated-match observation and independence conditions must be checked
  against each model's dynamics, selection, common shocks, and unit of
  observation.
\end{enumerate}

\textit{Source for item 1: Task 1.  Sources for item 2:
Eeckhout--Kircher, pp.\ 872--889; BS24/25, Sections 2--4.  Source for item 3:
BS20, Section 4.}

\subsection{4.14 Final classification}

The rigorous final answer is:

\[
\boxed{
\begin{array}{c}
\text{BS20's population correlation extends to every nested model once that
model}\\
\text{generates an observed matched-wage distribution.}\\[0.3em]
\text{Its estimator extends only after the BS repeated-match assumptions are
verified.}\\[0.3em]
\text{The later BS meeting/acceptance logic extends literally only when the
model}\\
\text{contains, or is augmented with, a potential-access and acceptance
block.}\\[0.3em]
\text{Kline additionally requires either restrictions making its edges
level-rationalizable}\\
\text{or a separate unrestricted transition wedge in the project
supermodel.}
\end{array}}
\tag{T4.64}
\]

Kline should remain a central model for the project, not a footnote.  Its edge
effects are the natural mover-level statistical interface, its cycle
restrictions show exactly where AKM compresses the data, and its selection
discussion explains why causal edge effects need not form a global firm
ranking.  The project's low-rank wage layer supplies one disciplined
rationalization of Kline cycle effects; the proposed \(\tau_{jk}\) layer
preserves the option to nest the unrestricted edge model when that broader
coverage is required.
\textit{Source: Kline, Sections 3.1--3.3; project implication is derived.}

\clearpage
\section{Task 5: Exogenous mobility versus additive separability}

\subsection{5.1 Resolution in one paragraph}

There is no substantive contradiction once the assumptions are separated into
atomic pieces.  ``Additive separability'' can mean either:

\begin{enumerate}
\item the \emph{pure wage-function restriction} that permanent worker--firm
  interactions are zero; or
\item that restriction \emph{plus} an exogeneity condition on the remaining
  transitory error.
\end{enumerate}

Volpe uses the first meaning and therefore presents additive separability and
exogenous mobility as two separate pillars.  Sorkin--Warwar use the second,
composite meaning.  Their ``additive separability'' assumption contains the
same transitory-error exogeneity condition as their ``exogenous mobility''
assumption and replaces mean-zero match effects by zero match effects.
Consequently, Sorkin--Warwar additive separability is strictly stronger than
their exogenous mobility.  Kline's review makes both viewpoints visible:
his statistical AKM model combines an additive expected-wage map with strict
exogeneity, while his potential-outcome analysis separates causal
identification of mover effects, heterogeneity, and selection on gains.

\textit{Sources: Sorkin--Warwar, ``Paired movers,'' March 2026 draft,
pp.\ 6--7; Sorkin Lecture 10, PDF pp.\ 97--105, especially slide 49/79;
Volpe, ``Firms and Earnings Inequality,'' PDF p.\ 21, slide 4/35; Kline,
pp.\ 14--16 and 29--34.  The reconciliation is derived below.}

\subsection{5.2 Notation chart for Task 5}

Tasks 1--4 notation remains in force.  In particular, \(D_{it}\) is worker
\(i\)'s firm at time \(t\), \(Y_{it}(j)\) is the potential log wage at firm
\(j\), and \((\alpha_i,\psi_j,X_{it},\beta)\) are the worker effect, firm
effect, observed covariates, and covariate coefficients.  The following chart
defines every new symbol.

\begin{longtable}{p{0.20\textwidth}p{0.52\textwidth}p{0.21\textwidth}}
\textbf{Symbol} & \textbf{Meaning} & \textbf{Source or status} \\
\hline
\(\omega_{ij}\) & Permanent statistical worker--firm match effect.  This
relabels Sorkin--Warwar's \(\mu_{ij}\) to avoid collision with the project
grand mean \(\mu\). & Sorkin--Warwar, p.\ 6; symbol changed. \\
\(\widetilde\varepsilon_{it}(j)\) & Transitory potential-wage remainder at
firm \(j\), after removing the permanent match component. & Sorkin--Warwar,
pp.\ 6--7; potential-outcome index is added. \\
\(\varepsilon^A_{it}\) & Observed residual from an additive AKM equation:
\(\varepsilon^A_{it}=\omega_{iD_{it}}+
\widetilde\varepsilon_{it}(D_{it})\). & Derived decomposition. \\
\(\mathcal H_{is},\mathcal H_i\) & The period-\(s\) assignment/covariate cell
\((D_{is},X_{is})\), and the entire assignment/covariate history
\(\{(D_{is},X_{is})\}_{s=1}^T\). & Derived shorthand. \\
\(\mathsf A_0\) & Pure additive separability:
\(\omega_{ij}=0\) for every potential worker--firm pair. & Derived label for
Sorkin--Warwar equation (6) and Volpe's second pillar. \\
\(\mathsf E_\varepsilon\) & Mean exogeneity of the transitory potential-wage
remainder with respect to mobility/covariates. & Derived atomic assumption. \\
\(\mathsf E_\omega\) & Mean exogeneity of permanent match effects with
respect to mobility/covariates. & Derived atomic assumption. \\
\(\mathsf{EM}_{SW},\mathsf{AS}_{SW}\) & Sorkin--Warwar exogenous mobility and
their composite additive-separability assumption. & Derived labels. \\
\(\mathsf{EM}_{V},\mathsf{AS}_{V}\) & Volpe's mobility pillar and pure
additive-separability pillar. & Derived labels. \\
\(\mathsf{AKM}_{V}\) & Intersection of Volpe's two pillars: mobility
exogeneity and pure additive separability. & Derived label. \\
\(\mathsf{SE}_{K}\) & Kline's strict-exogeneity condition for the statistical
AKM equation. & Kline, pp.\ 14--15. \\
\(g_i(j,k)\) & Worker \(i\)'s permanent conditional-mean wage gain from firm
\(j\) to firm \(k\). & Derived. \\
\(G_{i2}(j,k)\) & Kline's period-2 realized potential-wage gain
\(Y_{i2}(k)-Y_{i2}(j)\). & Derived shorthand for Kline, p.\ 32. \\
\(\mathsf{NSG}_{K}\) & Kline's no-selection-on-treatment-effects assumption:
\(G_{i2}(j,k)\perp(D_{i1},D_{i2})\). & Kline, p.\ 32; label derived. \\
\(Z,\theta,u\) & Generic OLS regressor vector, coefficient vector, and
regression residual in \(Y=Z'\theta+u\). & Standard notation. \\
\(D,Y(0),Y(1),\tau_i\) & Binary treatment, two potential outcomes, and
individual treatment effect \(\tau_i=Y_i(1)-Y_i(0)\). & Standard causal
notation. \\
\(d,d_1,d_2\) & Generic firm choices; \(d_t\) is the firm relevant in period
\(t\) in Kline's exclusion restriction. & Kline, pp.\ 29--30. \\
\(\operatorname{ATE}\) & Average treatment effect
\(\mathbb E[\tau_i]\). & Standard notation. \\
\(\theta_0,\theta_1\) & Intercept and binary-treatment slope in the OLS
representation of the average potential-outcome contrast. & Derived. \\
\(\Delta A_i\) & First difference \(A_{i2}-A_{i1}\) for any panel variable
\(A\); \(i'\) denotes a second worker. & Standard shorthand. \\
\(\xi_i,e_{it}\) & Mean-zero match heterogeneity and a transitory shock used
in the counterexamples. & Derived. \\
\end{longtable}

\subsection{5.3 The common wage decomposition}

Write the potential adjusted log wage of worker \(i\) at firm \(j\) as
\[
Y_{it}(j)
=
\alpha_i+\psi_j+X_{it}'\beta
+\omega_{ij}
+\widetilde\varepsilon_{it}(j).
\tag{T5.1}
\]
Consistency connects potential and observed wages:
\[
Y_{it}=Y_{it}(D_{it}).
\tag{T5.2}
\]
If the researcher estimates an additive AKM equation that omits
\(\omega_{ij}\), its observed residual is
\[
\varepsilon^A_{it}
=
\omega_{iD_{it}}
+\widetilde\varepsilon_{it}(D_{it}).
\tag{T5.3}
\]

Sorkin--Warwar define the permanent match component as the mean of the raw
worker--firm residual over all potential draws, not merely the periods in
which the pair is observed.  Equation (T5.1) is their equation (2), with
\(\omega_{ij}\) replacing their \(\mu_{ij}\) and with the potential firm
argument made explicit.
\textit{Source: Sorkin--Warwar, p.\ 6; notation translation is derived.}

The permanent conditional-mean gain from moving \(j\) to \(k\) is
\[
\begin{aligned}
g_i(j,k)
&\equiv
\left(\alpha_i+\psi_k+\omega_{ik}\right)
-\left(\alpha_i+\psi_j+\omega_{ij}\right)\\
&=
\psi_k-\psi_j+\omega_{ik}-\omega_{ij}.
\end{aligned}
\tag{T5.4}
\]
Thus:

\begin{enumerate}
\item If \(\omega_{ij}=0\) for all \(i,j\), then
  \(g_i(j,k)=\psi_k-\psi_j\) for every worker.  The permanent
  conditional-mean firm effect is constant across workers.
\item If \(\omega_{ij}\) can vary, firm effects can be heterogeneous even
  when match effects average to zero in every mobility cell.
\end{enumerate}

\textit{Derived directly from (T5.1).  ``Constant'' here concerns the
permanent conditional-mean firm contrast; it does not assert that every
realized transitory potential-wage draw is identical.}

\subsection{5.4 The two atomic dimensions}

\subsubsection{5.4.1 Pure additive separability is an outcome restriction}

Define
\[
\mathsf A_0:
\qquad
\omega_{ij}=0
\quad\text{for every }(i,j).
\tag{T5.5}
\]
This is a restriction on the wage surface.  It rules out a permanent
worker--firm interaction.  By itself it says nothing about whether workers
move in response to \(\widetilde\varepsilon_{it}(j)\).

\textit{Source: Sorkin--Warwar, equation (6), p.\ 7; Volpe, PDF p.\ 21.
The separation from mobility is derived.}

\subsubsection{5.4.2 Exogenous mobility is a selection restriction}

The pairwise conditioning convention used by Sorkin--Warwar and Kline can be
written
\[
\mathsf E_\varepsilon:
\qquad
\mathbb E[
\widetilde\varepsilon_{it}(j)
\mid
D_{is}=d,\ X_{is}=x
]
=0
\tag{T5.6}
\]
for all admissible \(j,d,s,t,x\), and
\[
\mathsf E_\omega:
\qquad
\mathbb E[
\omega_{ij}
\mid
D_{is}=d,\ X_{is}=x
]
=0.
\tag{T5.7}
\]
These are selection restrictions: within every conditioning cell, mobility
does not select the mean transitory shock or the mean permanent match
component.  They do not require either component to be zero for every
worker--firm pair.

Sorkin--Warwar state (T5.6)--(T5.7) with their realized-error notation and
for every \((s,t)\).  A conventional panel strict-exogeneity condition may
instead condition jointly on the full history \(\mathcal H_i\):
\[
\mathbb E[
\widetilde\varepsilon_{it}(D_{it})
\mid
\mathcal H_i
]
=0.
\tag{T5.8}
\]
Equation (T5.8) is generally stronger than a list of pairwise conditional
mean restrictions unless further conditions make the two conditioning
schemes equivalent.

\textit{Source for pairwise conditions: Sorkin--Warwar equations (3)--(4),
p.\ 6; Kline, pp.\ 14--15.  The distinction between pairwise and full-history
conditioning is derived.}

\subsubsection{5.4.3 Why the two atomic restrictions are not nested}

Neither \(\mathsf A_0\) nor \(\mathsf E_\varepsilon\) implies the other.
Two counterexamples prove both directions.

\paragraph{Exogenous mobility without additive separability.}
Let there be two firms, let
\[
\omega_{i0}=\xi_i,
\qquad
\omega_{i1}=-\xi_i,
\qquad
\mathbb E[\xi_i]=0,
\tag{T5.9}
\]
set \(\widetilde\varepsilon_{it}(j)=0\), and assign workers randomly,
independently of \(\xi_i\).  Then
\[
\mathbb E[\omega_{ij}\mid D_{is}=d]=0
\tag{T5.10}
\]
and mobility is exogenous to match quality, but
\(\omega_{ij}\neq0\) with positive probability.  Moreover,
\[
g_i(0,1)=\psi_1-\psi_0-2\xi_i,
\tag{T5.11}
\]
so the firm effect is heterogeneous.  Exogenous mobility holds and pure
additive separability fails.

\paragraph{Additive separability without exogenous mobility.}
Set \(\omega_{ij}=0\) for every pair, let \(e_{it}\sim N(0,1)\), and suppose
the current firm indicator satisfies
\[
D_{it}=\mathbf 1\{e_{it}>0\},
\qquad
\widetilde\varepsilon_{it}(j)=e_{it}.
\tag{T5.12}
\]
The wage surface is additively separable, but
\[
\mathbb E[e_{it}\mid D_{it}=1]
=
\mathbb E[e_{it}\mid e_{it}>0]
=
\sqrt{\frac{2}{\pi}}
\neq0.
\tag{T5.13}
\]
Mobility is endogenous to the wage fluctuation.  Pure additive separability
holds and exogenous mobility fails.

\textit{Both counterexamples are derived.  They formalize the two violations
listed on Volpe's slide: match interactions versus moves driven by wage
shocks.}

\subsection{5.5 Sorkin--Warwar: why their assumptions are nested}

Using the atomic assumptions above, Sorkin--Warwar's definitions are
\[
\mathsf{EM}_{SW}
\equiv
\mathsf E_\varepsilon\cap\mathsf E_\omega,
\tag{T5.14}
\]
and
\[
\mathsf{AS}_{SW}
\equiv
\mathsf E_\varepsilon\cap\mathsf A_0.
\tag{T5.15}
\]
But \(\mathsf A_0\) mechanically implies \(\mathsf E_\omega\), because the
conditional expectation of an identically zero variable is zero.  Therefore
\[
\boxed{
\mathsf{AS}_{SW}
\subset
\mathsf{EM}_{SW}.
}
\tag{T5.16}
\]
The inclusion is strict: (T5.9)--(T5.11) satisfy
\(\mathsf{EM}_{SW}\) but violate \(\mathsf{AS}_{SW}\).

This proves the Sorkin statement:

\begin{enumerate}
\item Under \(\mathsf{EM}_{SW}\), match effects may be nonzero and individual
  firm gains may be heterogeneous, but those effects average to zero within
  the mobility conditioning cells.
\item Under \(\mathsf{AS}_{SW}\), permanent match effects are zero
  pointwise, so permanent firm contrasts are constant across workers.
\end{enumerate}

Neither assumption requires random matching of worker and firm fixed
effects.  Workers may sort on \(\alpha_i\) and on the vector of average firm
effects \(\{\psi_j\}_{j=1}^J\); the restriction concerns selection on the
remaining match and transitory components.

\textit{Source: Sorkin--Warwar, pp.\ 6--7; Sorkin Lecture 10, PDF
pp.\ 97--105.  Set representation and strictness proof are derived.}

\subsection{5.6 Volpe: why the two restrictions are separate pillars}

Translating Volpe's \(\log W_{it}\) into the common outcome notation, his AKM
equation is
\[
Y_{it}
=
\alpha_i+\psi_{D_{it}}+X_{it}'\beta+\varepsilon^A_{it}
\tag{T5.17}
\]
and places two restrictions beneath it:

\begin{enumerate}
\item exogenous mobility: moves are independent of unobserved wage
  fluctuations; and
\item additive separability: worker--firm interactions are absent.
\end{enumerate}

In the atomic notation, the intended distinction is
\[
\mathsf{EM}_{V}\approx\mathsf E_\varepsilon,
\qquad
\mathsf{AS}_{V}\equiv\mathsf A_0.
\tag{T5.18}
\]
These are separate and non-nested by (T5.9)--(T5.13).  The AKM causal
interpretation requires both:
\[
\mathsf{AKM}_{V}
\equiv
\mathsf{EM}_{V}\cap\mathsf{AS}_{V}.
\tag{T5.19}
\]

There is consequently no disagreement with Sorkin--Warwar:
\[
\underbrace{\mathsf{AS}_{V}}_{\text{pure functional form}}
\neq
\underbrace{\mathsf{AS}_{SW}}_{\text{functional form plus error exogeneity}}.
\tag{T5.20}
\]
Indeed, after aligning conditioning sets,
\(\mathsf{AKM}_{V}\) and \(\mathsf{AS}_{SW}\) contain the same two atomic
ingredients.  The different diagram follows from attaching the name
``additive separability'' to different bundles.

Volpe's slide states a full-history form of strict exogeneity,
\[
\mathbb E[
\varepsilon^A_{it}
\mid
\alpha_i,\{\psi_{D_{is}},X_{is}\}_{s=1}^T
]
=0.
\tag{T5.21}
\]
Once \(\mathsf A_0\) is imposed, (T5.3) reduces to
\(\varepsilon^A_{it}=\widetilde\varepsilon_{it}(D_{it})\), so (T5.21) is the
mobility restriction on the remaining error.  Without \(\mathsf A_0\),
(T5.21) is merely a zero-conditional-mean restriction on the composite
\(\omega_{iD_{it}}+\widetilde\varepsilon_{it}(D_{it})\); it does not
separately prove that either component is zero.

\paragraph{Volpe's mover-change implication.}
For worker \(i\) moving from \(j\) to \(k\), (T5.1) gives
\[
\Delta Y_i
=
\psi_k-\psi_j+\Delta X_i'\beta
+\omega_{ik}-\omega_{ij}
+\Delta\widetilde\varepsilon_i.
\tag{T5.21a}
\]
Take two workers \(i\) and \(i'\) making the same transition and having the
same \(\Delta X\).  Under transitory-error exogeneity,
\[
\begin{aligned}
&\mathbb E[
\Delta Y_i-\Delta Y_{i'}
\mid
i,i'\text{ make }j\to k,\ \Delta X_i=\Delta X_{i'}]\\
&\qquad =
\mathbb E[
(\omega_{ik}-\omega_{ij})
-(\omega_{i'k}-\omega_{i'j})
\mid
i,i'\text{ make }j\to k,\ \Delta X_i=\Delta X_{i'}].
\end{aligned}
\tag{T5.21b}
\]
Pure additive separability makes the right-hand side zero.  Mobility
exogeneity is separately needed to remove the conditional mean of the
transitory changes.  This is the algebra behind the conditional
mover-change equality displayed on Volpe's slide.

\textit{Source: Volpe, PDF p.\ 21, slide 4/35.  The decomposition and
comparison to Sorkin--Warwar are derived.}

\subsection{5.7 Kline: the multiple layers}

\subsubsection{5.7.1 Layer 1: the statistical AKM model}

Kline begins with
\[
Y_{it}
=
\alpha_i+\psi_{D_{it}}+X_{it}'\beta+\varepsilon^A_{it}
\tag{T5.22}
\]
and the pairwise strict-exogeneity condition
\[
\mathsf{SE}_{K}:
\qquad
\mathbb E[
\varepsilon^A_{it}
\mid
D_{is}=j,\ X_{is}=x
]
=0
\quad
\text{for all }(s,t,j,x).
\tag{T5.23}
\]
Kline explicitly says that this condition embeds both:

\begin{enumerate}
\item mobility not driven by time-varying wage fluctuations; and
\item an additively separable mapping from worker and firm heterogeneity to
  expected log wages.
\end{enumerate}

This is the same bundling issue.  Equation (T5.22) supplies the additive
conditional-mean map, while (T5.23) prevents the remaining disturbance from
being selected by mobility.  If a permanent interaction is silently folded
into \(\varepsilon^A_{it}\), zero conditional mean of the \emph{sum} in
(T5.3) does not algebraically imply pointwise
\(\omega_{ij}=0\); the additive interpretation is a maintained restriction
on the wage map.

\textit{Source: Kline, pp.\ 14--16.  The last decomposition is derived.}

\subsubsection{5.7.2 Layer 2: causal edge effects without additive
separability}

Kline then drops the global AKM restriction and studies two-period potential
wages.  His first three causal assumptions are:

\begin{enumerate}
\item \textbf{Exclusion:}
  \(Y_{it}(d_1,d_2)=Y_{it}(d_t)\).
\item \textbf{Parallel trends:}
  \[
  \mathbb E[
  Y_{i2}(j)-Y_{i1}(j)
  \mid
  D_{i1}=j,D_{i2}=k
  ]
  =0.
  \tag{T5.24}
  \]
\item \textbf{Stationarity of the selected movers' mean gain:}
  \[
  \begin{aligned}
  &\mathbb E[
  Y_{i1}(k)-Y_{i1}(j)
  \mid D_{i1}=j,D_{i2}=k]\\
  &\qquad =
  \mathbb E[
  Y_{i2}(k)-Y_{i2}(j)
  \mid D_{i1}=j,D_{i2}=k]
  \equiv\Delta_{jk}.
  \end{aligned}
  \tag{T5.25}
  \]
\end{enumerate}

Under these assumptions,
\[
\mathbb E[
Y_{i2}-Y_{i1}
\mid
D_{i1}=j,D_{i2}=k
]
=
\Delta_{jk}.
\tag{T5.26}
\]
The edge effect is therefore an average treatment effect for workers who
actually move on edge \(j\to k\).  Treatment effects may be heterogeneous,
the wage surface may be non-additive, and different edges may concern
different selected worker populations.  Consequently, the edge effects need
not obey AKM cycle restrictions or define a global firm ranking.

\textit{Source: Kline, pp.\ 29--32, Assumptions 1--3 and Proposition 1.}

\subsubsection{5.7.3 Layer 3: no selection on treatment gains}

Kline next assumes
\[
\mathsf{NSG}_{K}:
\qquad
G_{i2}(j,k)
\equiv
Y_{i2}(k)-Y_{i2}(j)
\ \perp\
(D_{i1},D_{i2})
\quad
\text{for every }j\neq k.
\tag{T5.27}
\]
Then
\[
\Delta_{jk}
=
\mathbb E[G_{i2}(j,k)]
=
\psi_k^{\mathrm{ATE}}-\psi_j^{\mathrm{ATE}},
\tag{T5.28}
\]
so all edges refer to the same population and inherit transitivity.

Crucially, (T5.27) is not a constant-effects assumption.  For example,
\[
G_{i2}(j,k)
=
\psi_k^{\mathrm{ATE}}-\psi_j^{\mathrm{ATE}}+\xi_i(j,k),
\qquad
\xi_i(j,k)\perp(D_{i1},D_{i2}),
\tag{T5.29}
\]
allows nondegenerate treatment-effect heterogeneity while satisfying
\(\mathsf{NSG}_{K}\).

\textit{Source: Kline, pp.\ 32--34, Assumption 4 and Proposition 2;
equation (T5.29) is a derived example.}

\subsubsection{5.7.4 Kline's no-selection assumption is not identical to
Sorkin--Warwar exogenous mobility}

The restrictions act on different objects:

\begin{enumerate}
\item Sorkin--Warwar impose conditional \emph{mean} restrictions separately
  on permanent match levels and transitory errors.
\item Kline's Assumption 4 imposes full statistical independence between a
  pairwise treatment \emph{gain} and the two-period assignment.
\end{enumerate}

Therefore there is no unconditional nesting:

\begin{enumerate}
\item Mean independence does not imply full independence.  Mobility cells
  can have the same mean gain but different gain variances, satisfying the
  relevant mean restriction while violating (T5.27).
\item Independence of gains does not imply exogeneity of levels or
  transitory trends.  Mobility can select on a component common to
  \(Y_{i2}(j)\) and \(Y_{i2}(k)\) that cancels from their difference.
\end{enumerate}

With additional assumptions---joint distributional independence of all
potential match components, exclusion, and the relevant time restrictions---
a strengthened Sorkin-style mobility condition can imply Kline's
no-selection-on-gains condition.  That implication does not follow from the
published mean restrictions alone.

\textit{Derived comparison of Sorkin--Warwar equations (3)--(4) with Kline
Assumption 4.}

\subsection{5.8 Relation to ordinary OLS exogeneity}

\subsubsection{5.8.1 What \(\mathbb E[u\mid Z]=0\) says statistically}

Consider
\[
Y=Z'\theta+u.
\tag{T5.30}
\]
The conditional-mean restriction
\[
\mathbb E[u\mid Z]=0
\tag{T5.31}
\]
is equivalent to
\[
\mathbb E[Y\mid Z]=Z'\theta.
\tag{T5.32}
\]
Thus it simultaneously says that the conditional mean is correctly
specified by the chosen regressors and that no remaining conditional-mean
component is correlated with them.

The weaker OLS normal equation
\[
\mathbb E[Zu]=0
\tag{T5.33}
\]
is sufficient to characterize a population linear projection under the
usual rank condition, but it does not imply (T5.31).  Moreover, neither
(T5.31) nor (T5.33), by itself, gives \(\theta\) a causal interpretation.
A causal interpretation also requires a treatment definition, consistency,
and a design or selection assumption connecting observed and potential
outcomes.

\textit{Standard regression algebra; statements are derived.}

\subsubsection{5.8.2 Heterogeneous effects can satisfy OLS exogeneity}

For a binary treatment,
\[
Y_i
=
Y_i(0)+D_i\tau_i,
\qquad
\tau_i=Y_i(1)-Y_i(0).
\tag{T5.34}
\]
Suppose the potential outcomes are mean independent of assignment:
\[
\mathbb E[Y_i(d)\mid D_i]
=
\mathbb E[Y_i(d)]
\quad
\text{for }d\in\{0,1\}.
\tag{T5.35}
\]
Then
\[
\begin{aligned}
\mathbb E[Y_i\mid D_i=1]
-\mathbb E[Y_i\mid D_i=0]
&=
\mathbb E[Y_i(1)]-\mathbb E[Y_i(0)]\\
&=
\operatorname{ATE},
\end{aligned}
\tag{T5.36}
\]
even when \(\tau_i\) varies across workers.

Set
\[
\theta_0=\mathbb E[Y_i(0)],
\qquad
\theta_1=\operatorname{ATE},
\tag{T5.37}
\]
and define
\[
u_i
=
Y_i(0)-\theta_0
+D_i(\tau_i-\theta_1).
\tag{T5.38}
\]
Under (T5.35),
\[
\mathbb E[u_i\mid D_i]=0.
\tag{T5.39}
\]
Hence a correctly specified conditional-mean regression can identify an
average causal effect without constant individual treatment effects.  This
is the ordinary-OLS analogue of Sorkin--Warwar exogenous mobility.

This binary-treatment regression is saturated: its conditional mean has
only two treatment cells.  With continuous treatments, multiple treatments,
or additional covariates, a linear OLS coefficient may instead be a
particular weighted average of heterogeneous effects, and
\(\mathbb E[u\mid Z]=0\) additionally imposes the stated linear conditional-
mean form.

\textit{Derived potential-outcome calculation.}

\subsubsection{5.8.3 Constant effects do not eliminate selection bias}

Conversely, suppose
\[
\tau_i=\tau
\quad\text{for every }i,
\tag{T5.40}
\]
but assignment selects on the untreated potential outcome:
\[
\mathbb E[Y_i(0)\mid D_i=1]
\neq
\mathbb E[Y_i(0)\mid D_i=0].
\tag{T5.41}
\]
Then
\[
\begin{aligned}
&\mathbb E[Y_i\mid D_i=1]
-\mathbb E[Y_i\mid D_i=0]\\
&\qquad =
\tau
+\mathbb E[Y_i(0)\mid D_i=1]
-\mathbb E[Y_i(0)\mid D_i=0],
\end{aligned}
\tag{T5.42}
\]
which is not generally \(\tau\).  Constant effects are an outcome
restriction; exogeneity is a selection restriction.  One cannot substitute
for the other.

\textit{Derived potential-outcome calculation.}

\subsubsection{5.8.4 The exact AKM--OLS dictionary}

In an AKM regression, \(Z\) contains worker indicators, firm indicators, and
observed covariates.  The omitted-interaction decomposition is
\[
u_{it}
=
\omega_{iD_{it}}
+\widetilde\varepsilon_{it}(D_{it}).
\tag{T5.43}
\]
The correspondence is:

\begin{longtable}{p{0.29\textwidth}p{0.31\textwidth}p{0.31\textwidth}}
\textbf{Ordinary causal regression} &
\textbf{AKM analogue} &
\textbf{What it does not imply} \\
\hline
Correct conditional-mean functional form &
Additive worker plus firm expected-wage map &
Exogenous mobility \\
Mean independence of potential outcomes or gains &
Mobility does not select unobserved wage components or firm gains &
Constant worker-specific firm effects \\
Constant individual treatment effect &
No permanent worker--firm interaction &
Mobility exogeneity \\
\(\mathbb E[u\mid Z]=0\) for the maintained equation &
Strict exogeneity of the correctly defined AKM residual &
A decomposition showing separately why the residual has mean zero \\
\end{longtable}

The last row is important.  If \(u\) is defined as a composite residual,
\(\mathbb E[u\mid Z]=0\) can result from averaging or cancellation.  It does
not logically establish that the structural interaction is zero and the
transitory shock is separately exogenous.  Those are additional
interpretive restrictions.

\textit{Derived dictionary from (T5.1)--(T5.43).}

\subsection{5.9 Implications for the semistructural project}

The project's low-rank wage equation deliberately permits
worker--firm interactions:
\[
Y_{it}(j)
=
\alpha_i+\psi_j
+\mathbf u_i'\Lambda\mathbf v_j
+X_{it}'\beta
+\widetilde\varepsilon_{it}(j).
\tag{T5.44}
\]
Therefore the project generally rejects pure AKM additive separability
\(\mathsf A_0\).  This does \emph{not} force mobility to be endogenous.
The assignment process can still satisfy a mean-independence or
no-selection-on-gains condition.

The assumptions should be recorded in separate blocks:

\begin{enumerate}
\item \textbf{Wage-surface block:} state whether
  \(\mathbf u_i'\Lambda\mathbf v_j\) is zero, low rank, or unrestricted.
\item \textbf{Mobility block:} state whether assignment is mean independent
  of transitory shocks, permanent interaction terms, potential wage levels,
  or pairwise gains.  These are different restrictions.
\item \textbf{Time block:} state exclusion, parallel trends, and stationarity
  when mover changes are given a causal interpretation.
\item \textbf{Target block:} distinguish a selected-mover edge effect,
  an unconditional average firm effect, a constant individual firm effect,
  and a descriptive observed-match wage statistic.
\end{enumerate}

This separation yields four consequences:

\begin{enumerate}
\item The BS20 observed-wage sorting functional does not require additive
  separability or a causal mobility assumption merely to exist as a
  descriptive observed-match object.  Its estimator still requires the
  BS20 sampling conditions.
\item A Kline edge effect can have a causal selected-mover interpretation
  under exclusion, parallel trends, and stationarity even when the project
  interaction is nonzero.
\item Turning edge effects into a common cardinal firm ranking requires an
  additional restriction such as no selection on treatment gains; low rank
  alone does not provide that restriction.
\item Imposing a zero interaction term does not make the firm effect causal
  unless the mobility/error condition is also imposed.
\end{enumerate}

\textit{Items 1--4 are derived from Tasks 1--4, Sorkin--Warwar pp.\ 6--7,
Volpe PDF p.\ 21, and Kline pp.\ 29--34.}

\subsection{5.10 Final hierarchy}

The complete reconciliation is:
\[
\boxed{
\begin{array}{c}
\text{Volpe: }\quad
\mathsf{AS}_{V}=\mathsf A_0,\qquad
\mathsf{EM}_{V}\approx\mathsf E_\varepsilon.\\[0.25em]
\text{Sorkin--Warwar: }\quad
\mathsf{EM}_{SW}=\mathsf E_\varepsilon\cap\mathsf E_\omega.\\[0.25em]
\text{Sorkin--Warwar: }\quad
\mathsf{AS}_{SW}=\mathsf E_\varepsilon\cap\mathsf A_0
\subset\mathsf{EM}_{SW}.\\[0.25em]
\text{Kline statistical AKM:}\\
\text{additive expected-wage map plus strict exogeneity.}\\[0.25em]
\text{Kline causal edges: exclusion plus parallel trends plus
stationarity.}\\[0.25em]
\text{Kline common-population ranking:}\\
\text{causal edge conditions plus no selection on gains.}
\end{array}}
\tag{T5.45}
\]

Thus the slogan ``additive separability is stronger than exogenous
mobility'' is correct under the Sorkin--Warwar bundled definitions.  The
slogan ``additive separability and exogenous mobility are separate pillars''
is correct when additive separability denotes only the wage-function
restriction, as in Volpe's slide.  Ordinary OLS makes the same distinction:
a homogeneous or additive outcome equation restricts treatment effects,
whereas \(\mathbb E[u\mid Z]=0\) restricts the conditional relation between
assignment and the residual.  A causal regression commonly needs both a
correct outcome specification and a credible assignment restriction, but
average effects can be identified under exogenous assignment without
constant effects.

\textit{Source for the named definitions: Sorkin--Warwar, pp.\ 6--7;
Sorkin Lecture 10, slide 49/79; Volpe, slide 4/35; Kline, pp.\ 14--16 and
29--34.  The hierarchy and final OLS comparison are derived.}

\end{document}
